

\documentclass[sigconf]{acmart}
% \setlength{\parindent}{1em}
% \citestyle{acmauthoryear}
\usepackage{graphicx}
\usepackage{subcaption}
\usepackage{enumitem}
\usepackage{xcolor}
%\usepackage{minted}
% \usepackage[frozencache,cachedir=.]{minted}
% \usepackage{tcolorbox}
\usepackage{xspace}
\usepackage{multirow}
\usepackage{algorithm}
% \usepackage{algpseudocode}
\usepackage{amsmath}
\usepackage{stmaryrd} % semantic brackets
\usepackage{dsfont} % indicator function

\usepackage{listings}
\usepackage{soul} % For text highlighting
\usepackage{cleveref}

% \usepackage{tabularx}
\usepackage[utf8]{inputenc}
\usepackage{xcolor}
\definecolor{softblue}{RGB}{65, 105, 225} % SteelBlue-like, softer than default
\DeclareMathSymbol{\mlq}{\mathord}{operators}{``}
\DeclareMathSymbol{\mrq}{\mathord}{operators}{`'}


\definecolor{lightblue}{RGB}{245, 245, 220}
\definecolor{grey}{RGB}{169,169,169}
\setlength{\parskip}{0pt}
\setlength{\parindent}{0pt}

\definecolor{darkgreen}{RGB}{54, 156, 90}
\definecolor{lightblue}{RGB}{235, 247, 255}
\definecolor{niceblue}{RGB}{52, 164, 235}
\definecolor{nicered}{RGB}{219, 42, 48}
\definecolor{niceorange}{RGB}{224, 114, 54}
\definecolor{lightred}{RGB}{255, 227, 228}
\definecolor{lightorange}{RGB}{252, 235, 227}
\definecolor{orange1}{RGB}{214, 178, 131}
\definecolor{brown}{RGB}{163, 117, 57}
\definecolor{lightbrown}{RGB}{224, 163, 83}
\definecolor{commentcolor}{RGB}{101, 163, 178}
\definecolor{lightpurple}{RGB}{203, 195, 227}
%224, 163, 83
%252, 213, 164
\xdefinecolor{highlight_gold}{RGB}{204, 153, 0}

\newcommand{\sergey}[1]{{\color{red}\textbf{(Sergey:} #1\textbf{)}}}
\newcommand{\dimitris}[1]{{\color{red}\textbf{(Dimitris:} #1\textbf{)}}}
\newcommand{\yihan}[1]{{\color{red}\textbf{(Yihan:} #1\textbf{)}}}

\definecolor{nicegreen}{RGB}{52, 207, 98}
\definecolor{lightgreen}{RGB}{223, 247, 230}
% \usepackage[framemethod=tikz]{mdframed}
% \mdfsetup{
%     roundcorner=10pt,
%     linewidth=1pt,
%     backgroundcolor=gray!20,
%     innertopmargin=12pt,
%     innerbottommargin=12pt,
%     innerrightmargin=12pt,
%     innerleftmargin=12pt,
%     skipabove=12pt,
%     skipbelow=12pt,
%     nobreak=false,
%     align=justify  % Ensures text justification and line wrapping
% }
\usepackage[most]{tcolorbox}



% Define a new tighter highlight box
\newtcbox{\highlight}{on line,
    arc=1pt, % Rounded corners (adjust as needed)
    colframe=red!20, % Border color
    colback=red!20, % Background color
    boxrule=0.0pt, % Border thickness
    boxsep=0pt, % Remove padding inside the box
    left=1pt, right=1pt, top=1pt, bottom=1pt % Minimal margins
}
\newtcbox{\highlightblue}{on line,
    arc=1pt, % Rounded corners (adjust as needed)
    colframe=blue!12, % Border color
    colback=blue!12, % Background color
    boxrule=0.0pt, % Border thickness
    boxsep=0pt, % Remove padding inside the box
    left=1pt, right=1pt, top=1pt, bottom=1pt % Minimal margins
}
\definecolor{codehighlight}{RGB}{237, 173, 12} % Light orange

\definecolor{highlightbg}{RGB}{237, 230, 12} % Light orange

% Define highlight colors
\definecolor{marker-yellow}{RGB}{255, 255, 150} % Soft yellow like a highlighter
\definecolor{marker-green}{RGB}{200, 255, 200} % Light green for variation

% Define highlight colors
\definecolor{highlight-yellow}{RGB}{255, 255, 150} % Yellow marker
\definecolor{highlight-blue}{RGB}{180, 210, 255} % Blue marker


\newcommand{\projName}{HoarePrompt\xspace}
\newcommand{\benchName}{CoCoClaNeL\xspace}
% \setminted{
%     fontsize=\footnotesize,
%     frame=none,
%     breaklines=true,
%     breakautoindent=true,
%     showspaces=false,
%     linenos=false
% }
\lstset{
    language=Python,
    basicstyle=\ttfamily\footnotesize,
    keywordstyle=\bfseries\color{violet},
    stringstyle=\color{magenta},
    commentstyle=\itshape\color{gray},
    showstringspaces=false,
    breaklines=true,          
    flexiblecolumns=true,     
    columns=fullflexible,     
    frame=none,               
    xleftmargin=0pt,          
    xrightmargin=0pt,         
    numbers=none,             
     moredelim=[is][\color{red}]{@@}{@@} , % Highlight with yellow,  % Allows LaTeX commands in lstlisting
      moredelim=[is][\color{blue}]{!}{!}  % Highlight with yellow,  % Allows LaTeX commands in lstlisting
}

% from there: https://tex.stackexchange.com/questions/21598/how-to-color-math-symbols
\makeatletter
\def\mathcolor#1#{\@mathcolor{#1}}
\def\@mathcolor#1#2#3{%
  \protect\leavevmode
  \begingroup
    \color#1{#2}#3%
  \endgroup
}
\makeatother


% Larger font listing style
\lstdefinestyle{largePython}{
    language=Python,
    basicstyle=\ttfamily\normalsize,
    keywordstyle=\bfseries\color{darkgrey},
    stringstyle=\color{darkgrey},
    commentstyle=\itshape\color{gray},
    showstringspaces=false,
    breaklines=true,
    frame=none,
    numbers=none,
     moredelim=[is][\color{red}]{@}{@},   % Highlight specific words in red
    moredelim=[is][\color{blue}]{!}{!}
}


% === ACM Metadata for ICSE 2026 ===
\copyrightyear{2026}
\acmYear{2026}
\setcopyright{cc}
\setcctype{by}
\acmConference[ICSE '26]{2026 IEEE/ACM 48th International Conference on Software Engineering}{April 12--18, 2026}{Rio de Janeiro, Brazil}
\acmBooktitle{2026 IEEE/ACM 48th International Conference on Software Engineering (ICSE '26), April 12--18, 2026, Rio de Janeiro, Brazil}
\acmPrice{}
\acmDOI{10.1145/3744916.3773206}
\acmISBN{979-8-4007-2025-3/2026/04}


\begin{document}
\setlength{\parindent}{1em} 
\setlength{\parskip}{0pt} 
\title{\projName: Structural Reasoning About Program Correctness in Natural Language}

\author{Dimitrios Stamatios Bouras}
\email{2501112125@stu.pku.edu.cn}
\affiliation{%
  \institution{Key Laboratory of HCST (PKU), MoE, SCS, Peking University}
  \city{Beijing}
  \country{China}
}

\author{Yihan Dai}
\email{2501112020@stu.pku.edu.cn}
\affiliation{%
  \institution{Key Laboratory of HCST (PKU), MoE, SCS, Peking University}
  \city{Beijing}
  \country{China}
}

\author{Tairan Wang}
\email{tairan.wang.22@ucl.ac.uk}
\affiliation{%
  \institution{University College London}
  \city{London}
  \country{United Kingdom}
}

\author{Yingfei Xiong}
\email{xiongyf@pku.edu.cn}
\affiliation{%
  \institution{Key Laboratory of HCST (PKU), MoE, SCS, Peking University}
  \city{Beijing}
  \country{China}
}

\author{Sergey Mechtaev}
\email{mechtaev@pku.edu.cn}
\authornote{Sergey Mechtaev is the corresponding author.}
\affiliation{%
  \institution{Key Laboratory of HCST (PKU), MoE, SCS, Peking University}
  \city{Beijing}
  \country{China}
}

\authorsaddresses{} % keep empty for camera-ready if your venue requests it

\newcommand{\jp}[1]{\textcolor{blue}{#1}}

\newcommand{\changes}[1]{%
% \textcolor{red}{#1}\xspace
#1\xspace
}


%\author{Anonymous}
\authorsaddresses{} % No addresses for? now

% Abstract
\begin{abstract}
While software requirements are often expressed in natural language, verifying the correctness of a program against such requirements is a hard and underexplored problem. Large language models (LLMs) are promising candidates for addressing this challenge, however, our experience shows that they are ineffective in this task, often failing to detect even straightforward bugs. To address this gap, we introduce \projName, a novel approach that adapts fundamental ideas from program verification to natural language artifacts. Inspired by the strongest postcondition calculus, \projName employs a systematic, step-by-step process in which an LLM generates natural language descriptions of reachable program states at various code points. To manage loops, we propose few-shot-driven k-induction, an adaptation of the k-induction method widely used in model checking. Once program states are described, \projName leverages the LLM to assess whether the program, annotated with these state descriptions, conforms to the natural language requirements. For evaluating the quality of classifiers of program correctness with respect to natural language requirements, we constructed \benchName, a challenging dataset of solutions to programming competition problems. Our experiments show that \projName improves the MCC by 61\% compared to directly using Zero-shot-CoT prompts for correctness classification. Furthermore, \projName outperforms a classifier that assesses correctness via LLM-based test generation by an MCC increase of 106\%. The inductive reasoning mechanism contributes a 26\% boost to MCC, underscoring its effectiveness in managing loops.
\end{abstract}

\maketitle

% Introduction Section
\section{Introduction}
\label{sec:introduction}

Assessing whether a program meets its requirements is a key part of software quality assurance. In practice, these requirements are often written in natural language, geared towards human-to-human communication or designed as prompts for LLMs. Verifying program correctness against such requirements is an inherently difficult and unexplored challenge. One approach is address it is to generate tests based on requirements using an LLM, and use these tests to check program correctness~\cite{chen2022codet}. More generally, natural language requirements can be translated in formal ones, and then used in conjunction with verification tools~\cite{endres2024can}. Despite the promise of these approaches, our experiments show that LLMs tend to generate incorrect or incomplete assertions when handling complex requirements, which reduces their accuracy. Another option is to directly employ LLMs to assess a program's conformance to requirements, relying on LLMs' ability to reason about both natural language and code~\cite{liu2024codemind}. Our experiments show that the state-of-the-art LLMs struggle with this task, often failing to identify even simple bugs. We attribute this limitation to the models' inability to reason about the intricacies of program semantics.

To address this gap, we introduce \projName, a program analysis approach for reasoning about program correctness w.r.t. natural language requirements that synergizes LLM reasoning techniques with ideas from program verification. In its core is the \emph{natural Hoare triple}, a variant of the traditional Hoare triple~\cite{hoare1969axiomatic} with pre- and postconditions expressed in natural language; for example, $$\mathcolor{softblue}{\{ \text{x is a full name} \}} \,\mathtt{x = x.split()[0]}\, \mathcolor{softblue}{\{ \text{x is a first name} \}}.$$ Inspired by the strongest postcondition calculus~\cite{dijkstra1976discipline}, \projName employs a step-by-step process to generate natural language descriptions of reachable program states at various points in the code. This is achieved by sampling \emph{natural strongest postconditions (NSP)} from $\mathcal{SP}$, our LLM-based natural Hoare triple completion model defining the conditional probability distribution $\mathcal{SP}(\mathcolor{softblue}{\mathrm{Post}}\,\mid\,\mathcolor{softblue}{\mathrm{Pre}}, C),$
where $\mathcolor{softblue}{\mathrm{Pre}}, \mathcolor{softblue}{\mathrm{Post}}$ are natural pre- and postconditions, $C$ is a program fragment. Inspired by methods to enhance LLM reasoning about program execution~\cite{ni2024next}, \projName uses these state descriptions to annotate a program, and then prompts an LLM to assess if the annotated program aligns with the given requirements.

\begin{figure*}[t]
  \begin{center}
    \includegraphics[width=\textwidth]{motivating_example.png}
  \end{center} 
  \caption{Annotating code with natural language descriptions of reachable program states, which are automatically inferred by \projName, enables an LLM to find a bug that is missed by the straightforward Vanilla prompt.\label{fig:motivating_example}}
  \vspace{-2mm}
\end{figure*}

\projName's key novelty in comparison with traditional program analysis and verification is in using natural language as the medium of reasoning about program semantics, as opposed to, e.g., first-order logic~\cite{filliatre2013why3} or monotonic functions over lattices~\cite{cousot1977abstract}. This enables a direct application of \projName to assessing program correctness w.r.t. natural language requirements without hard and time-consuming formal modeling. In comparison to previous techniques for reasoning about code with LLMs~\cite{ni2024next,liu2024codemind,jiang2024ledex}, \projName reasons about all reachable states simultaneously instead of individual concrete states. This enables \projName to operate statically, without executing the program and without reliance on provided failing tests. \changes{In comparison with directly prompting LLMs to analyze programs, \projName simplifies this task by decomposing it into small, independent subtasks, which are solved step-by-step in the spirit of chain-of-thought reasoning~\cite{wei2022chain}.}

Loops pose a fundamental challenge in program analysis and verification because they give rise to infinite execution paths~\cite{si2018learning}. Our experiments demonstrate that LLMs also struggle to summarize the semantics of complex loops, producing incorrect natural postconditions. To address this problem, we introduce a novel method, \emph{few-shot-driven k-induction}, synergizing few-shot learning~\cite{brown2020language} with the idea of k-induction~\cite{donaldson2011software} widely used in model checking. This method starts with iteratively unwinding a loop $k$ times to compute natural postconditions after each iteration. After that, it performs an induction inference step by prompting an LLM to compute the postcondition for the entire loop, using the computed natural Hoare triples for individual iterations as few-shot examples.

\begin{table}[t]
\centering
\caption{Success rates for detecting the bug in \Cref{fig:motivating_example} via prompts with and without state annotations, using four LLMs, across 60 responses for each, with the temperature 0.5.}
\vspace{-2mm}
\label{tab:motivating_example_results}
\begin{tabular}{l c c}
\toprule
\textbf{Model} & \textbf{Not Annotated} & \textbf{Annotated} \\
\midrule
Qwen2.5-7B          & 0\%   & 11.7\% \\
Qwen2.5-Coder-32B   & 15\%  & 85\%  \\
Qwen2.5-72B    & 35\% & 85\% \\
\changes{Llama3.1-70B}    & \changes{50\%} & \changes{100\%} \\
\bottomrule
\end{tabular}
\vspace{-3mm}
\end{table}


\begin{figure*}[t]
  \begin{center}
    \includegraphics[width=\textwidth]{loop_motivation.png}
  \end{center}
  \vspace{-2mm}  
  \caption{Few-shot-driven $k$-induction unrolls the loop three times, computes the natural strongest postconditions (NSP) for each iteration sequentially, and uses the resulting natural Hoare triples as few-shot examples for inductively inferring the postcondition for the entire loop, which enhances precision compared to non-inductive inference, or only a single unrolling.\label{fig:kinduction_example}}
\end{figure*} 
To evaluate \projName, we created \benchName, a dataset of program and natural language requirement pairs, where some programs meet the requirements and others do not due to human-introduced bugs. All programs and requirements were published since January 2024, after the training cut-off date of many SOTA LLMs to prevent data leakage. We formulated program correctness assessment as a binary classification problem. Our experiments show that \projName, at the expense of a higher token use, improves the Matthews Correlation Coefficient (MCC) by 61\% compared to directly using Zero-shot-CoT prompts. Furthermore, \projName outperforms a classifier that assesses correctness via LLM-generated tests by increasing the MCC by 106\%. The few-shot-driven $k$-induction mechanism contributes a 26\% boost to MCC. \changes{We also show that \projName improves code generation when used as a feedback mechanism, outperforming both direct generation (by 19.4\%) and test-based feedback (by 12.7\%).}

In summary, the paper makes the following contributions:
\begin{itemize}[nosep, topsep=0pt, partopsep=0pt, parsep=0pt, itemsep=0pt]
\item \projName, the first program analysis approach using natural language as the medium for reasoning about program semantics.
\item Few-shot-driven $k$-induction, the first approach that enhances LLM's ability to summarize loops.
\item \benchName, a challenging benchmark for program correctness classification.
\item An empirical study of classifying program correctness w.r.t. natural language requirements \changes{and guiding code generation}, demonstrating the effectiveness of \projName.
\end{itemize}
All code, scripts, and data, to reproduce this work are available at 
\changes{ Figshare: \href{https://figshare.com/s/5d39b388bd9ff2c7ffab}{https://figshare.com/s/5d39b388bd9ff2c7ffab}}. 
\vspace{-2mm}
\section{Motivating Examples}
\label{sec:motivating}

To motivate our approach, we first show that annotating a program with natural language descriptions of reachable program states enhances LLM's reasoning about program correctness, and then show how our approach precisely summarizes loops.
\vspace{-2mm}
\subsection{State Descriptions Enhance LLM Reasoning}
\label{sec:motivating_states}

Consider the following problem: ``given an integer list, the task is to find the maximum product of any sublist.'' \Cref{fig:motivating_example} (left) presents a flawed implementation of Kadane's algorithm to solve this problem. The issue lies in incorrect indentation, which results in \texttt{best\_so\_far} being updated only once, after the loop ends, instead of being updated during each iteration.

Despite the bug's simplicity, SOTA LLMs struggle to identify it consistently, and thus often wrongly classify the program as correct. Providing evidence that an LLM fails to solve a task is inherently challenging due to the model's non-deterministic behavior~\cite{ouyang2025empirical} and its sensitivity to prompt design~\cite{sclar2023quantifying}. To address these challenges, we made every effort to present objective evidence by crafting six distinct prompts: a straightforward (Vanilla) prompt, a Zero-shot-CoT prompt~\cite{kojima2022large}, an open-ended prompt~\cite{arora2022ask}, and combinations of these strategies, each containing the problem description and the code snippet. For each prompt, we generated 10 responses from four different LLMs, employing various temperature settings. More details about the experimental setup and results for this example can be found in the supplementary materials (\Cref{appendix:motivating_example}).

The results in \Cref{tab:motivating_example_results} demonstrate that LLMs, using the above six prompts, labelled collectively as ``Not Annotated'', detect the bug with only a small probability. An example of an incorrect response for the Vanilla prompt is presented in \Cref{fig:motivating_example} on the left.

\projName addresses this problem as follows. It first extracts the program's natural precondition by prompting an LLM to summarize the set of valid input values based on the problem description. In this case, the inferred precondition is $\mathcolor{softblue}{\{ \text{xs is a list of integers} \}}.$ Then, it iteratively queries its natural strongest postcondition (NSP) model to propagate reachable state descriptions along the program's control flow graph (CFG). For example, the Hoare Triplet for the first if statement:
% computing the NSP of the first if statement based on the program's precondition will result in the following natural Hoare triple:
\setlength{\jot}{0pt}
\begin{align*}
  &\mathcolor{softblue}{\{ \text{xs is a list of integers} \}}\\
  &\texttt{if (not xs):}\\
  &\quad\texttt{return float('-inf')}\\
  &\mathcolor{softblue}{\{ \text{xs is a list of integers, and xs is not empty} \}}
\end{align*}
\noindent capturing the fact that if an execution reaches the program point after the if statement, the list is guaranteed to be non-empty.

\projName continues to propagate the states until it reaches the end of the program. Loops, which present a particular challenge, are handled via inductive reasoning as illustrated in \Cref{sec:motivating_example_induction}.

Finally, \projName annotates the program with the reachable states at key program points (return statements, after top-level if statements and loops) as presented in \Cref{fig:motivating_example} on the right. Adding such annotations to the six prompts significantly increases LLM's success rate of correctness classification as shown in \Cref{tab:motivating_example_results}. \changes{This improvement stems from \projName’s compositional reasoning strategy. Since it reasons about individual blocks in isolation, the loop is interpreted based only on its local semantics. In this case, \projName accurately captures that the loop updates only consider sublists ending at the final element, with the resulting postcondition annotation enabling better error detection.}
% highlights this narrow behavior, acting as an "alarm bell" for the LLM. Despite the overall code appearing structurally similar to a correct implementation, this annotation helps the LLM realize that the semantics diverge significantly, enabling better error detection.}

\vspace{-3mm}
\subsection{Inductively Reasoning About Loops}
\label{sec:motivating_example_induction}

Reasoning about loops is challenging, because they induce infinite behaviours~\cite{si2018learning}, and our experience shows that LLMs also struggle to summarize them. To alleviate this problem, we take inspiration for the $k$-induction technique~\cite{donaldson2011software} widely used in model checking, which first checks that the property holds for all states in the first $k$ steps of the program's execution, and then proves that if a property holds for any k consecutive steps, then it holds for the next step.

To illustrate our adaptation of $k$-induction, which we refer to as \emph{few-shot-driven $k$-induction}, consider the program with a loop in \Cref{fig:kinduction_example}, labelled as ``Program''. Assume that its precondition is given in ``Precondition''. Inferring its natural strongest postcondition (NSP) using the straightforward ``Non-Inductive NSP Prompt'' with Qwen2.5-72B results in an incorrect natural postcondition mistakenly stating the \texttt{min\_diff} is the smallest absolute difference between any two consecutive elements in the list.

\projName unrolls the loop $k$ times, with the semantics of each iteration represented by code snippets in ``Iteration 1'', ..., ``Iteration k''. Then, it sequentially computes natural strongest postconditions for each iteration, using the postcondition of the previous iteration as the precondition of the next one (the first iteration uses the precondition of the entire loop). The resulting triples
\begin{align*}
  \mathcolor{softblue}{\{ \text{Precondition} \}} \,&\mathtt{Iteration\ 1}\, \mathcolor{softblue}{\{ \text{Postcondition 1} \}},\\
  &...\\
  \mathcolor{softblue}{\{ \text{Postcondition k-1} \}} \,&\mathtt{Iteration\ k}\, \mathcolor{softblue}{\{ \text{Postcondition k} \}}
\end{align*}
are used as few-shot examples to infer the postcondition of the entire loop in ``Inductive NSP Prompt (k=3)''. The intuition of this step is that LLMs generalize from few relevant examples~\cite{brown2020language}, and the examples above are relevant by construction. As shown in the corresponding LLM response, the enhanced prompt enabled the LLM to correctly infer that \texttt{min\_diff} is the smallest absolute difference between any  element in the list and the first element.

Interestingly, with k=1, the mistake persisted, as the model failed to generalize from a single example. Based on our pilot study (\Cref{sec:hyperparameters}), we chose $k=3$ as the best performing configuration.

\begin{figure}[t]
    \centering
    \includegraphics[width=\linewidth]{example1-sec5_blue.png}
    \vspace{-3mm}
    \caption{\projName succinctly describes program behavior in human-centric terms, which helps align formal semantics with informal requirements.
    \label{fig:motivating_expressive}}
    \vspace{-5mm}
\end{figure}
\vspace{-3mm}


\subsection{Expressive Power of Natural Language}
\label{sec:motivating_NL}
LLMs exhibit limitations in symbolic reasoning due to their reliance on natural language~\cite{xu2024faithful,can_llm_reason}. Consequently, \projName avoids formal logical inference, opting instead for natural language as the reasoning medium. While this approach does not provide formal correctness guarantees, it generates structured informal proofs, which, as demonstrated in \Cref{sec:evaluation}, improve the  LLMs' ability to effectively align program semantics with informal requirements.

% We hypothesize that the efficacy of this method stems from natural language's expressive power. An example in \Cref{fig:motivating_expressive} shows that Qwen2.5-72B computes a natural strongest postcondition for a code fragment, succinctly summarizing the variable \texttt{sat} using human-centric terms rather than formal execution semantics, leveraging its extensive training data. This expressivity is hard to achieve in symbolic logic due to its inherently limited vocabulary. More examples are given in the supplementary materials (\Cref{appendix:NL}).

 The method’s effectiveness stems from natural language’s expressivity. In \Cref{fig:motivating_expressive}, Qwen2.5-72B derives an NSP that succinctly describes \texttt{sat} in human terms, beyond formal semantics—an expressivity symbolic logic lacks. Further examples are in \Cref{appendix:NL}.


% \begin{figure}[t]
%     \centering
%     \includegraphics[width=0.9\linewidth]{example2-sec5.png}
%     \caption{Example 2 - Unique Email Domains \label{fig:unique_email}}
%     \vspace{-2mm}
% \end{figure}

% \begin{figure}[t]
%     \centering
%     \includegraphics[width=0.9\linewidth]{example3-sec5.png}
%     \caption{Example 3 - First Name Extraction \label{fig:first_name}}
%     \vspace{-2mm}
% \end{figure}

% Background Section
\vspace{-2mm}
\section{Background}
\label{sec:background}

Hoare logic~\cite{hoare1969axiomatic} is a formalism for reasoning about program correctness, at the core of which are specifications called \emph{Hoare triples} $\{\mathrm{Pre}\}\;C\;\{\mathrm{Post}\},$ meaning that the postcondition $\mathrm{Post}$ holds in any state reachable by executing the code $C$ from an initial state in which the precondition $\mathrm{Pre}$ holds, provided that the code terminates. Pre- and postconditions are predicates over program states, that are typically expressed in first-order logic~\cite{filliatre2013why3}. Strongest postcondition calculus~\cite{dijkstra1976discipline} automates verification using Hoare logic. The strongest postcondition (SP) of a program \( C \) w.r.t. a precondition \( \text{Pre} \), denoted as \( \text{sp}(\text{Pre}, C) \), is the most precise logical formula describing all possible states reachable after executing \( C \) from a state satisfying \( \text{Pre} \). Formally, $$\models \{\text{Pre}\}\;C\;\{\text{Post}\} \quad \text{iff} \quad \text{sp}(\text{Pre}, C) \implies \text{Post}.$$ The following algorithm for computing SP is defined compositionally for \textsc{While}~\cite{nielson2015principles}, a simple programming language that has only assignments ($x := e$), sequences ($S_1; S_2$), if statements, one loop instruction (while) and a single type (integer):
% \begin{align*}
%   &\mathrm{sp}(\mathrm{Pre}, x := e) = \mathrm{Pre}[x \mapsto e] \\[1mm]
%   &\mathrm{sp}(\mathrm{Pre}, S_1; S_2) = \mathrm{sp}(\mathrm{sp}(\mathrm{Pre}, S_1), S_2) \\[1mm]
%   &\mathrm{sp}(\mathrm{Pre}, \texttt{if } b \texttt{ then } S_1 \texttt{ else } S_2) =\\
%   &\quad \mathrm{sp}(\mathrm{Pre \land b}, S_1)  \vee \mathrm{sp} (\mathrm{Pre  \land \neg b}, S_2) \\[1mm]
%   &\mathrm{sp}(\mathrm{Pre}, \texttt{while } b \texttt{ do } S) =\\
%   &\quad (\mathrm{Pre} \land \neg b) \lor \mathrm{sp}(\mathrm{sp}(\mathrm{Pre} \land b, S), \texttt{while } b \texttt{ do } S)
% \end{align*}
{\setlength{\abovedisplayskip}{2pt}
 \setlength{\belowdisplayskip}{2pt}
 \setlength{\jot}{1pt}
\begin{align*}
  \mathrm{sp}(\mathrm{Pre}, x := e) &= \mathrm{Pre}[x \mapsto e] \\
  \mathrm{sp}(\mathrm{Pre}, S_1; S_2) &= \mathrm{sp}(\mathrm{sp}(\mathrm{Pre}, S_1), S_2) \\
  \mathrm{sp}(\mathrm{Pre}, \texttt{if } b \texttt{ then } S_1 \texttt{ else } S_2) &=
      \mathrm{sp}(\mathrm{Pre \land b}, S_1) \vee
      \mathrm{sp}(\mathrm{Pre \land \neg b}, S_2) \\
  \mathrm{sp}(\mathrm{Pre}, \texttt{while } b \texttt{ do } S) &=
      (\mathrm{Pre} \land \neg b) \,\vee {}\\[-1mm]
  &\quad \mathrm{sp}(\mathrm{sp}(\mathrm{Pre \land b}, S),
      \texttt{while } b \texttt{ do } S)
\end{align*}}
\vspace{-2ex}



\noindent The last rule, which handles loops, is recursively defined, making it hard to compute.

%% \paragraph{Sequential Composition} For a sequence of operations \(\alpha ; \beta\), the strongest postcondition is computed by propagating the postcondition of \(\alpha\) through \(\beta\):
%%     \[
%%     \operatorname{sp}(\text{Pre}, (\alpha;\beta)) = \operatorname{sp}(\operatorname{sp}(\text{Pre}, \alpha), \beta).
%%     \]
    
%% \paragraph{Nondeterministic Choice} For a nondeterministic choice \(\alpha \cup \beta\), the strongest postcondition is the disjunction of the strongest postconditions of both branches:
%%     \[
%%     \operatorname{sp}(\text{Pre}, \alpha \cup \beta) = \operatorname{sp}(\text{Pre}, \alpha) \lor \operatorname{sp}(\text{Pre}, \beta).
%%     \]
    
%% \paragraph{Guard} A guard \(\mathcal{?}Q\) enforces \(Q\) while preserving the precondition \(\text{Pre}\), resulting in:
%%     \[
%%     \operatorname{sp}(\text{Pre}, \mathcal{?}Q) = \text{Pre} \land Q.
%%     \]
    
%% \paragraph{Repetition} The strongest postcondition for a loop \(\alpha^*\) is defined recursively, accounting for zero or more iterations:
%%     \[
%%     \operatorname{sp}(\text{Pre}, \alpha^*) = \text{Pre} \lor \operatorname{sp}(\operatorname{sp}(\text{Pre}, \alpha),\alpha^*).
%%     \]
    
% \paragraph{Assignment} For an assignment \(x \gets e(x)\), the strongest postcondition asserts the existence of a prior value \(x'\) satisfying the precondition:
%     \[
%     \operatorname{sp}(P(x), x \gets e(x)) = \exists x'.\left(x = e(x')\right) \land P(x') \quad (x' \not\in e(x), P(x))
%     \]

%% \paragraph{Conditional} 
%% The conditional \(\text{if}\ Q\ \text{then}\ \alpha\ \text{else}\ \beta\) is semantically equivalent to the nondeterministic choice between two guarded paths \((\mathcal{?}Q ; \alpha)\) and \((\mathcal{?}\neg Q ; \beta)\). Formally:
%% \[
%% \text{if } Q\ \text{then}\ \alpha\ \text{else}\ \beta \triangleq (\mathcal{?}Q ; \alpha) \cup (\mathcal{?}\neg Q ; \beta)
%% \]
%% So the strongest postcondition for this conditional statement is derived by combining the strongest postconditions of its branches under their respective guards.  Applying the rules for \emph{nondeterministic choice} (\(\cup\)) and \emph{sequential composition} (\(;\)), the strongest postcondition of the conditional is:
%%     \[
%%     \operatorname{sp}(\text{Pre}, \text{if}\ Q\ \text{then}\ \alpha\ \text{else}\ \beta) = \operatorname{sp}(Q \land \text{Pre}, \alpha) \lor \operatorname{sp}(\neg Q \land \text{Pre}, \beta)
%%     \]

% One of the key properties of strongest postcondition is \textit{compositionality}, which enables the systematic propagation of program states through sequential compositions. Specifically, for two sequential programs \(C_1\) and \(C_2\), the strongest postcondition of the combined program \(C_1; C_2\) with respect to an initial precondition \(P\) can be expressed as:

% \[
% \text{sp}(P, C_1; C_2) = \text{sp}(\text{sp}(P, C_1), C_2).
% \]

% Leveraging the \textit{compositionality} property, we can use strongest postcondition calculation to propagate the program's states forward. By comparing the computed states (or the program's expected behavior) at each step with a specific specification, we can determine whether the program's execution meets the desired requirements.

%% In program verification, loop analysis presents a fundamental challenge due to its potential unbounded iteration space. Researchers have developed analytical methods to represent this structural complexity, primarily through two complementary paradigms: invariant-based \textit{over-approximation} using logical constraints, and bounded verification via \textit{under-approximation} with fixed-bound loop unrollings.

%% While bounded verification requires significantly less manual effort than invariant derivation, this under-approximation approach suffers from inherent limitations: for certain inputs that demand numerous iterations, even thousands of unrollings are still insufficient to establish program correctness. 

The $k$-induction method~\cite{donaldson2011software} applies mathematical-induction principles to practical loop analysis. Let $P(s_i)$ denote the target property at loop state $s_i$ after $i$ iterations. The method has two components: the \emph{base case} (1), which checks that the property holds for the first $k$ iterations, and the \emph{inductive step} (2), which proves that any window of $k$ satisfying iterations implies the next one also satisfies it:
\[
\text{(1) }\bigwedge_{i=0}^{k-1} P(s_i)
\qquad\qquad
\text{(2) }\forall n\ge 0.\,\left( \bigwedge_{i=0}^{k-1} P(s_{n+i}) \right) \rightarrow P(s_{n+k})
\]




% \section{Summary - to help Prof. Mechtaev}

% \projName is a framework designed to enhance the reasoning capabilities of LLms with respect to program correctness. It combines \emph{Hoare logic} principles, tailored code preprocessing, loop unrolling, structured LLM prompting, and multiple verification ``classifiers.'' This section provides an overview of the methodology and the eight classifiers that \projName offers.

% \subsection{Input Preparation and Preprocessing}
% \projName processes two primary inputs:
% \begin{enumerate}
%     \item A \emph{problem description} specifying the intended functionality of the program.
%     \item A \emph{Python script} containing the full implementation, which may include function definitions, global code, and import statements.
% \end{enumerate}
% To ensure unbiased analysis, \projName preprocesses the script by removing import statements, stripping comments, and replacing function names with generic placeholders. If the script contains only global code, it is encapsulated within a function for uniform processing. Each function is then analyzed individually, and the results are combined to assess the correctness of the entire implementation.

% \subsection{Precondition Extraction}
% \projName infers an \emph{initial precondition} for each function by querying the LLM with the problem description and the function signature (excluding its body). The precondition defines the expected input format, initial variable states, and necessary assumptions about the function’s execution. For example, if a function processes a list of ages, the precondition might specify that \texttt{numbers} is a list of non-negative integers. This extracted precondition serves as the foundation for further analysis, ensuring that subsequent reasoning aligns with the intended input constraints.

% \subsection{Hoare Logic Foundations and Code Decomposition}
% \projName follows the principles of Hoare logic to systematically break down a function into smaller segments or statements. Program correctness is modeled using Hoare triples:
% \[
% \{P\} \quad C \quad \{Q\}
% \]
% where \(P\) represents the \emph{precondition} (the program state before executing the code segment \(C\)), and \(Q\) is the \emph{postcondition} (the state after execution).

% Unlike traditional Hoare logic, which applies formal inference rules to derive postconditions, \projName leverages an LLM to infer them dynamically. Given a precondition (expressed in natural language) and a corresponding code segment, \projName queries the LLM to generate the expected postcondition, also in natural language. 

% To enhance Accuracy, few-shot learning examples are incorporated into the LLM prompt, demonstrating how to compute postconditions for similar code constructs. Different templates are used depending on the type of code segment (e.g., assignments, conditionals, loops, exceptions), aligning with the underlying logic of Hoare reasoning. The generated postcondition then serves as the precondition for subsequent statements, enabling structured, recursive analysis of the program’s behavior.

% \subsection{Compound Statements}
% Compound statements, such as blocks within function bodies, are processed recursively, with each sub-command analyzed according to its type. For example, if a sub-command is a return statement, it is handled immediately, and its postcondition is recorded as the final state for that execution path.

% The recursive analysis halts upon encountering an early return, reflecting realistic program behavior where any code following a return statement is unreachable and thus excluded from further processing. This ensures that execution flows are accurately modeled while avoiding redundant computations.

% \subsection{Simple Statements}
% Single-line operations such as assignments and variable updates are directly analyzed by the LLM. Given the current precondition, \projName queries the LLM to infer the resulting postcondition after executing the statement. 

% To optimize efficiency, if the \texttt{concat\_simple} option is enabled, consecutive simple statements are processed together in a single API call. This reduces the number of LLM queries and minimizes token usage while maintaining the logical flow of the analysis.

% \subsection{Handling Print Statements}
% In many programs, critical output is produced through \texttt{print} statements rather than return values. To accurately capture this behavior, \projName processes each \texttt{print} statement separately by querying the LLM to determine the exact output generated at that point in execution. This extracted information is stored and later incorporated into the annotated execution states tree.

% If explicit print tracking is not required, users can disable this feature by setting \texttt{annotate-prints=false} in the configuration. In this case, print statements are treated as simple statements, and their output is not recorded in the execution state analysis.


% \subsection{Return Statements}
% Return statements serve as distinct termination points within a function. When a function has multiple return points, \projName computes and tracks the postcondition for each return statement individually. 

% This approach ensures that all possible execution paths are accurately documented. Additionally, once a return statement is encountered, further analysis of subsequent code is halted, reflecting real execution behavior and preventing unnecessary processing of unreachable statements.


% \subsection{Conditional Statements}
% If-else blocks and related constructs are processed by first incorporating the condition into the precondition. \projName evaluates both the \texttt{if} and \texttt{else} branches independently and then merges their results into a unified postcondition at the join point.

% The process begins by refining the precondition for each branch:
% \begin{enumerate}
%     \item An API call is made to update the precondition with the \texttt{if} condition, forming the refined precondition for the \texttt{if} branch.
%     \item Similarly, a separate API call integrates the \emph{negation} of the condition to establish the precondition for the \texttt{else} branch.
% \end{enumerate}
% Elif constructs (\texttt{elif}) are automatically converted into nested \texttt{if-else} blocks using Python's Abstract Syntax Tree (AST), ensuring a standardized analysis.

% Each branch is then analyzed recursively based on its internal structure, using the corresponding updated precondition. Once the final postconditions for both the \texttt{if} and \texttt{else} branches are determined, another LLM query is issued to merge them into a single postcondition representing the program state after the entire conditional block has executed.


% \subsection{Exception Handling}
% Try-except blocks are analyzed by separately evaluating the execution paths of the \texttt{try} and \texttt{except} clauses. Each path results in a distinct postcondition, which is later merged into a comprehensive postcondition capturing both normal execution and exception-handling behavior.

% The process follows these steps:
% \begin{enumerate}
%     \item The \texttt{try} block is analyzed first, with \projName computing its postcondition under the assumption that no exceptions occur.
%     \item Each \texttt{except} block is then processed independently. Before analyzing an \texttt{except} clause, the precondition is updated to reflect that an exception of the corresponding type has occurred. This ensures that the LLM correctly models the execution state when the exception is caught.
%     \item A separate postcondition is derived for each \texttt{except} clause, capturing how the program state changes when a particular exception is raised and handled.
%     \item Once all individual postconditions are determined, \projName aggregates them into a unified postcondition using the \texttt{TryTriple} structure, ensuring a consistent representation of all possible execution flows.
% \end{enumerate}

% \subsection{Loop Handling via Unrolling} Loops, whether \texttt{for} or \texttt{while}, are analyzed in \projName using the same fundamental approach: loop unrolling. Instead of attempting to derive loop invariants, \projName simulates the execution of the loop by computing the program state after a fixed number of iterations before generalizing the final postcondition. Additionally, \projName offers the option of no unrolling, where the entire loop code is provided to the LLM, which then reasons about the final execution state after the loop completes. This mode can be enabled by setting \texttt{loop-unrolling-count} to 0.

% \paragraph{Loop Unrolling Strategy} The unrolling process follows these steps: \begin{enumerate} \item The loop’s precondition is established based on the state before its execution. \item The loop body is executed up to a configurable number of times (\texttt{loop-unrolling-count}, defaulting to 3). \item Each unrolled iteration updates the precondition for the next iteration by querying the LLM. \item If a return statement is encountered, unrolling stops early, reflecting realistic execution paths. \item After unrolling, the LLM is asked to infer the final postcondition, representing the overall effect of the loop once it has completed execution. \end{enumerate}

% \paragraph{No Unrolling Option} When \texttt{loop-unrolling-count} is set to 0, \projName bypasses the unrolling process entirely. In this mode, the LLM is presented with the entire loop code along with its initial state and is asked to directly infer the final execution state after the loop completes. This approach is useful when loop unrolling is impractical or when a direct reasoning process is preferred, allowing the LLM to consider the loop logic in its entirety rather than through individual iterations. Not using unrolling also reduces the time and number of API calls required but can negatively impact performance as can be seen in the motivational example and results section.

% \paragraph{Handling of While and For Loops} While \texttt{for} and \texttt{while} loops use distinct LLM prompt templates, the core unrolling logic remains the same: \begin{itemize} \item For \texttt{while} loops, execution proceeds as long as the condition holds. The precondition is updated at each step, and unrolling stops either upon reaching the configured limit or if the loop terminates early. \item For \texttt{for} loops, iteration is treated similarly, with the loop variable incrementing across iterations. If the loop iterates over a range or collection, the execution state evolves accordingly, capturing changes in variables and control flow. \end{itemize}

% Unrolling enables \projName to capture how loop variables evolve over the first few iterations, helping the LLM reason about their overall impact. Instead of requiring explicit invariant derivation, which can be complex and unreliable, the unrolled states provide concrete examples from which the LLM can generalize. Alternatively, the no-unrolling option allows the LLM to consider the full loop logic at once, providing flexibility in handling complex loops.

% The choice of a fixed unroll factor of 3 (configurable) balances Accuracy and efficiency. Larger values increase the number of LLM calls, which can become computationally expensive, especially for nested loops. A motivating example for this approach is presented in Section~\ref{unrolling-motivation}.

% By systematically tracking the program state across iterations and allowing the LLM to infer the final loop postcondition, \projName effectively models loops without requiring explicit formal reasoning about invariants, while also offering a no-unrolling mode for direct loop reasoning when needed.
% \subsection{Annotated Code Tree and Final Verification}
% As each statement is processed, the derived postcondition is appended as an \emph{inline annotation} to the code, forming an \emph{annotated code tree}. However, to avoid redundancy and information overload, only key execution states are retained in the final tree, while unimportant or highly repetitive postconditions are omitted.

% \paragraph{Selective State Annotations}
% Since postconditions often include all information present in the precondition alongside updates, they can become highly repetitive. Additionally, excessive information in LLM prompts has been observed to degrade performance rather than improve it. To mitigate this, execution states are only annotated at key program points:
% \begin{itemize}
%     \item \textbf{After conditionals} (\texttt{if-else}): The final postcondition merging both branches is included.
%     \item \textbf{After loops}: Only the overall loop postcondition is retained, not intermediate states from unrolling.
%     \item \textbf{After return statements}: The function's state at each return point is recorded.
%     \item \textbf{After print statements}: If output tracking is enabled (\texttt{annotate-prints=true}), the state reflecting printed values is included.
%     \item \textbf{After try-except blocks}: The merged postcondition encapsulating all exception handling paths is annotated.
% \end{itemize}

% \paragraph{Function Summary and Final Verification}
% Beyond inline annotations, \projName employs multiple verification strategies to assess program correctness. Depending on the classifier used, verification can take different forms:

% \begin{itemize}
%     \item \textbf{Annotated Tree Verification(simple Tree)}: In this approach, only the annotated execution states are provided to the LLM. The correctness decision is based on analyzing the structured execution flow without an explicit function summary.
%     \item \textbf{Function Summary Verification}: Instead of presenting the full annotated tree, this approach condenses the function’s overall behavior into a high-level description, which the LLM then compares against the problem description to determine correctness.
%     \item \textbf{Combined Verification(complex tree)}: This classifier provides both the annotated execution states and the function summary, allowing the LLM to make a correctness judgment with access to both fine-grained execution details and a broader functional overview.
% \end{itemize}

% \subsection{Validator Classifiers}
% The \emph{Validator Classifiers} enhance the LLM’s reasoning by incorporating an iterative verification step. Instead of discarding the initial correctness assessment from the naive approach, \projName stores this original response and prompts the LLM to verify or refine its judgment using additional structured insights. 

% \paragraph{Motivation}
% While the naive approach of directly querying the LLM about correctness can often detect real bugs, adding execution states or function summaries may sometimes introduce confusion. This can happen when:
% \begin{itemize}
%     \item The execution states provide excessive detail, making the LLM overly critical or leading it to overlook real bugs.
%     \item The execution states fail to directly capture a bug that was correctly identified in the original assessment.
%     \item The LLM, when presented with extra information, becomes more cautious and incorrectly flags code as incorrect when it is actually valid.
% \end{itemize}
% The validator classifiers aim to mitigate these issues by asking the LLM to \emph{verify} whether its original assessment still holds when provided with additional insights.

% \paragraph{Types of Validator Classifiers}
% Three validator classifiers extend the naive approach by incorporating structured program analysis in a verification step:
% \begin{itemize}
%     \item \textbf{Annotated Tree Validator(simple verify):} The LLM first determines correctness without additional execution states. It is then provided with the annotated execution tree and asked to reassess whether its initial judgment was accurate.
%     \item \textbf{Function Summary Validator(function summary verify):} Instead of the full execution states, the LLM is given only the high-level function summary in the verification step, refining its decision with a broader perspective.
%     \item \textbf{Combined Validator(complex verify):} The LLM receives both the annotated execution tree and the function summary, ensuring it has the most complete context for verification.
% \end{itemize}

% \paragraph{Validation Process}
% The validation classifiers follow a two-step approach:
% \begin{enumerate}
%     \item \emph{Initial Assessment:} The LLM is queried with a naive prompt, producing an initial correctness judgment.
%     \item \emph{Verification Step:} The LLM is presented with additional insights (execution states, function summary, or both) and asked to reconsider its response in light of this new information.
% \end{enumerate}

% This iterative approach is particularly effective in cases where:
% \begin{itemize}
%     \item The naive approach correctly identifies a bug, ensuring it is not overlooked even if execution states do not explicitly highlight it.
%     \item The execution states introduce excessive scrutiny, preventing false positives.
%     \item Combining structured reasoning with the LLM’s own prior assessment results in a more reliable correctness evaluation.
% \end{itemize}


% % \begin{figure*}[h]
% %     \centering
% %     \begin{minipage}[t]{0.57\textwidth}
% %           \vspace{0pt}
% %         % Vanilla Prompting - Prompt Box
% %          % Vanilla Prompting - Response Box
% %         \begin{tcolorbox}[colframe=darkgreen!100, colback=gray!3, title=Problem Description, fonttitle=\centering, before upper=\raggedright, boxrule=0.8mm, sharp corners=southwest, left=2pt, right=2pt, top=2pt, bottom=2pt, boxsep=1pt]
% %         \footnotesize
% %            Equalize the total number of people (students, administrators, and professors) between Building A (12 classes) and Building B (16 classes) by removing excess individuals from the more populated building. The input consists of five lists: students\_A (12 values), admins\_A (12 values), students\_B (16 values), admins\_B (16 values), and profs\_B (16 values), representing the number of people in each class. Each class has different number of people than the others. Each remaining person in both buildings contributes \$10 to a bonus pool, which is evenly distributed among those removed. Return the per-person bonus for the moved individuals.
% %         \end{tcolorbox}
% %         \begin{tcolorbox}[colframe=nicered!90, colback=gray!3, title=Tester prompt with original function, fonttitle=\centering, before upper=\raggedright, boxrule=0.8mm, sharp corners=southwest, left=2pt, right=2pt, top=2pt, bottom=2pt, boxsep=1pt]
% %             {\fontsize{9pt}{11pt}\selectfont \textbf{Give me test cases(input/ output pairs) for the follwoing program and based on the problem description}}
% %             \begin{minted}[escapeinside=||]{python}
           
% % def func_1(students_A, admins_A, students_B, admins_B, profs_B):
% %     total_A = sum(students_A) + sum(admins_A)
% %     total_B = sum(students_B) + sum(admins_B) + sum(profs_B)
% %     surplus = abs(total_A - total_B)
% %     stayers = total_A + total_B - surplus
% %     bonus_pool = stayers * 10
% %     |\highlight{\textcolor{black}{ bonus\_per\_moved = bonus\_pool / surplus}}|
% %     return bonus_per_moved
   
% %             \end{minted}
% %         \end{tcolorbox}
            
% %         % Vanilla Prompting - Response Box
% %         \begin{tcolorbox}[colframe=nicered!90, colback=gray!3, title=Tester Response, fonttitle=\centering, before upper=\raggedright, boxrule=0.8mm, sharp corners=southwest, left=2pt, right=2pt, top=2pt, bottom=2pt, boxsep=1pt]
% %         \footnotesize
% %             <The rest of the test cases are ommited for brevity>
          
% % Input = 
% % [10,10,10,10,10,10,10,10,10,10,10,10], [2,2,2,2,2,2,2,2,2,2,2,2], [10,10,10,10,10,10,10,10,10,10,10,10,10,10,10,10], [2,2,2,2,2,2,2,2,2,2,2,2,2,2,2,2], [3,3,3,3,3,3,3,3,3,3,3,3,3,3,3,3] \\
% % Output =
% % zero division exception \# \highlight{\textcolor{black}{wrong the 2 building don't have same size}}
% % \\
% % Input =[20,20,20,20,20,20,20,20,20,20,20,20], [5,5,5,5,5,5,5,5,5,5,5,5], [10,10,10,10,10,10,10,10,10,10,10,10,10,10,10,10], [2,2,2,2,2,2,2,2,2,2,2,2,2,2,2,2], [3,3,3,3,3,3,3,3,3,3,3,3,3,3,3,3]\\
% % Output: 
% % 10 \# \highlight{\textcolor{black}{wrong, actual is 80}}


% %         \end{tcolorbox}
% %     \end{minipage}
% %     \hfill
% %     \begin{minipage}[t]{0.42\textwidth}
% %         % Prompted with Execution States - Prompt Box
% %           \vspace{0pt}
% %         \begin{tcolorbox}[colframe=niceblue, colback=gray!3, title=Prompt with Execution States as comments in the code, fonttitle=\centering, before upper=\raggedright, boxrule=0.8mm, sharp corners=southwest, left=2pt, right=2pt, top=2pt, bottom=2pt, boxsep=1pt]
% %         {\fontsize{9pt}{11pt}\selectfont \textbf{Is the following program correct based on the program description?}}
% %             \begin{minted}[escapeinside=||]{python}
% %     def func_1(students_A, admins_A, students_B, admins_B, profs_B):
% %     total_A = sum(students_A) + sum(admins_A)
% %     total_B = sum(students_B) + sum(admins_B) +sum(profs_B)
% %     surplus = abs(total_A - total_B)
% %     stayers = total_A + total_B - surplus
% %     bonus_pool = stayers * 10
% %     bonus_per_moved = bonus_pool / surplus
% %     return bonus_per_moved
% %     #The program returns `bonus\_per\_moved` which is 
% %     #`bonus\_pool / surplus`. Given that 'total\_A' is 
% %     #a positive integer equal to the number of people
% %     # originally at building A, and 'total\_B' is a 
% %     # positive integer equal to the number
% %     # of people originally at building B  
% %     #`bonus\_pool is a positive integer equal 
% %     #to 10*(total\_A+  total\_B - surplus) 
% %     #|\highlightblue{\textcolor{commentcolor}{and `surplus` is a positive integer or 0}}|
% %     # equal to the  absolute difference between `total\_A`
% %     #and `total_B`, the program returns 
% %     #`(10*(total\_A+  total\_B - surplus)) / abs(total\_A - total\_B)`.
    
% %             \end{minted}
% %         \end{tcolorbox}

% %         % Prompted with Execution States - Response Box
% %         \begin{tcolorbox}[colframe=niceblue, colback=gray!3, title= Response for Prompt with Execution States, fonttitle=\centering, before upper=\raggedright, boxrule=0.8mm, sharp corners=southwest, left=2pt, right=2pt, top=2pt, bottom=2pt, boxsep=1pt]
% %         \footnotesize
% %             The provided solution is \textbf{\textcolor{blue}{incorrect}} as it could lead to a potential division by zero exception since \highlightblue{\textcolor{black}{the variable \texttt{surplus}} could be 0 } in the case where the number of people in the 2 buildings is equal.
% %         \end{tcolorbox}
% %     \end{minipage}
% %     \caption{Comparison of Tester Prompting and Prompting with Execution States. }
% %     \label{fig:prompt_test_comparison}
% % \end{figure*}


 
% \subsection{Nine Classifiers summary}
% \projName provides eight distinct ``classifiers,'' representing different levels of analysis detail and verification steps:

% \begin{enumerate}
%     \item \textbf{Vanilla No-CoT}:  
%     A direct, single-step query that asks, ``Is the code correct?'' with no chain-of-thought or additional annotations.

%     \item \textbf{Vanilla (CoT)}:  
%     Similar to Vanilla (No-CoT), but instructs the LLM to provide an explicit chain-of-thought explaining how it reaches its correctness judgment.

%     \item \textbf{Simple Tree}:  
%     Uses \emph{\projName}'s basic execution-state annotations (the “simple” code tree) to help the LLM see intermediate states. The LLM then decides correctness with this tree as context.

%     \item \textbf{Complex Tree}:  
%     An extension of \emph{Simple Tree} that also integrates a higher-level function summary into the code tree. This additional summary can clarify global function behavior.

%     \item \textbf{Function Summary}:  
%     Rather than providing detailed inline execution-state annotations, this classifier gives the LLM a short textual summary of the function’s actions, prompting it to compare the summary against the problem description.

%     \item \textbf{Simple Verify}:  
%     A two-step approach. First, the LLM gives an answer via \emph{Vanilla} prompts. Then, in a second query, we provide the \emph{Simple Tree} annotations and ask it to \emph{verify or update} its initial correctness judgment.

%     \item \textbf{Complex Verify}:  
%     Similar to \emph{Simple Verify}, but uses the \emph{Complex Tree} annotations (including the function summary) during the second verification step.

%     \item \textbf{Function Summary Verify}:  
%     The LLM first answers via \emph{Vanilla} prompting. Then, in a follow-up verification prompt, the classifier supplies the function summary to refine or correct the initial response.

%    \item \textbf{Tester}:
%     Executes LLM-generated Python test cases to determine correctness. Passing all assertions means correct; failures mean incorrect. Serves as a baseline to compare testing-based verification with \projName’s execution-state reasoning.
% \end{enumerate}

\section{\projName}
\label{sec:hoareprompt}

\begin{figure}
    \centering
    \includegraphics[width=0.75\linewidth]{workflow_new.png}
    \vspace{-2mm}
    \caption{\projName's correctness classification workflow.\label{fig:workflow}}
    \vspace{-2mm}
\end{figure}

We introduce the \emph{natural Hoare triple}, a Hoare triple where pre- and postconditions are expressed in natural language (NL). To differentiate natural Hoare triples from traditional Hoare triples, we highlight sentences from NL with blue: $$\mathcolor{softblue}{\{\text{Pre}\}}  \, C \, \mathcolor{softblue}{\text{\{Post\}}}.$$

In analogy with the strongest postcondition calculus, we introduce the \emph{natural strongest postcondition (NSP)}. The traditional algorithm \( \text{sp}(\text{Pre}, C) \) cannot be applied in the context of NL. Instead, we create $\mathcal{SP}$, an LLM-based natural Hoare triple completion model defining a conditional probability distribution $$\mathcal{SP}(\mathcolor{softblue}{\mathrm{Post}} \mid\,\mathcolor{softblue}{\mathrm{Pre}}, \mathrm{Stmt}).$$ Let $\llbracket \cdot \rrbracket: \mathrm{NL} \to 2^S$ capture the human interpretation of NL state descriptions as a mapping from sentences to sets of program states $S$. Then, the goal to achieve by $\mathcal{SP}$ to ensure plausible reasoning can be measured using the following objective function:

\begin{definition}[NSP Objective Function] \label{def:objective} The objective function for $\mathcal{SP}$ is the cross-entropy loss between predicted postconditions and the traditional strongest postconditions computed over the interpretations of natural language preconditions: $$\mathcal{L} = -\sum_{\mathcolor{softblue}{\mathrm{Post}} \in \mathrm{NL}} \mathbb{I}(\llbracket \mathcolor{softblue}{\mathrm{Post}} \rrbracket \equiv \mathrm{sp}(\llbracket \mathcolor{softblue}{\mathrm{Pre}} \rrbracket, C)) \log \mathcal{SP}(\mathcolor{softblue}{\mathrm{Post}} \mid \mathcolor{softblue}{\mathrm{Pre}}, C),$$ where \( \mathbb{I}(\cdot) \) is the indicator function.
\end{definition}

\projName applies $\mathcal{SP}$ to classify program correctness w.r.t. natural language requirements using the workflow depicted in \Cref{fig:workflow}. First, it prompts an LLM to extract a precondition (a description of valid program inputs) from the requirements. Second, it traverses the program's control-flow graph (CFG), while sampling natural strongest postconditions from $\mathcal{SP}$, starting from the precondition, for each code segment. Finally, it annotates key program points with the inferred reachable state descriptions, and prompts an LLM to check whether the code meets the requirements.

\subsection{Precondition Extraction}

\projName prompts an LLM to extract a precondition from the given requirements in order to identify the set of valid program inputs. For example, for the problem description
\[
\text{``given $N$, write a function that finds the $N$th Fibonacci number''}
\]
and the function signature \texttt{def func(num)}, the LLM produces the precondition:
\vspace{-4mm}
\[
\mathcolor{softblue}{\{\text{num is a non\mbox{-}negative integer}\}}.
\]
We implement precondition extraction using few-shot learning (see supplementary materials, \Cref{pre_prompt,pre_multi_prompt}).


%% \begin{figure}
%%     \centering
%%     \includegraphics[width=1\linewidth]{example_def_long.png}
%%     \caption{An example of precondition extraction}
%%     \label{fig:def-example-label}
%% \end{figure}

\begin{figure}[t]
    \centering
    \includegraphics[width=\linewidth]{example_if_blue.png}
    \caption{Computing NSP for an \texttt{if-else} statement.\label{fig:ifelse-example-label}}
    \vspace{-2mm}
\end{figure}

\subsection{NSP Computation}
\label{sec:nsp}

In principle, $\mathcal{SP}$ can be realized by fine-tuning an LLM to optimize the NSP objective function (\Cref{def:objective}). Due to budget constraints, we opted instead to adapt an LLM through prompt engineering and in-context learning. Specifically, we curated a pilot dataset, distinct from our evaluation set, and carefully designed prompts and few-shot examples for Python programs by eyeballing the LLM's output on this dataset.

\paragraph*{Simple Statements} Atomic operations such as assignments are processed by directly prompting an LLM to compute the postcondition. We use a dedicated prompt with few-shot examples, presented in supplementary materials \Cref{simple_statement_prompt}, that tasks the model to infer the updated state based on variable modifications while preserving relevant information from the precondition. Apart from that, we use dedicated prompts (in supplementary materials \Cref{return_prompt,print_prompt}) to capture effects of return and print commands, since program correctness often depends on them.

\paragraph*{Compound Statements} Similar to the traditional sp, $\mathcal{SP}$ is defined compositionally, which enables it to precisely reason about complex programs by breaking them down into simpler fragments. For example, given a sequence of statements $S_1; S_2$, $\mathcal{SP}$ is computed by first completing natural Hoare triple $\mathcolor{softblue}{\{\text{Pre}\}}S_1\mathcolor{softblue}{\{\text{Post}\}}$, and then its postcondition $\mathcolor{softblue}{\text{Post}}$ serves as the precondition for $S_2$. To reduce the number of LLM queries, we group a configurable number of atomic statements together using a dedicated prompt (in supplementary materials \Cref{compund statement prompt}). 

%% \begin{figure}
%%     \centering
%%     \includegraphics[width=1\linewidth]{example-while.png}
%%     \caption{Computing NSP for a \texttt{while-do} statement with few-shot driven k-induction (k=3)}
%%     \label{fig:while-example-label}
%% \end{figure}

\paragraph*{If Statements} \Cref{fig:ifelse-example-label} shows how \texttt{if-else} statements are processed by \projName. Given a precondition $\mathcolor{softblue}{\{\text{Pre}\}}$ and a statement $\mathtt{if}\ b\!:S_1\ \mathtt{else}\!:\ S_2$, it first samples preconditions for the branches:
$$
\mathcolor{softblue}{\mathrm{Pre_{if}}}\sim \mathcal{SP}(\,\cdot\mid\,\mathcolor{softblue}{\mathrm{Pre}}, \mathtt{assert}\ b),\ \mathcolor{softblue}{\mathrm{Pre_{else}}}\sim \mathcal{SP}(\,\cdot\mid\,\mathcolor{softblue}{\mathrm{Pre}},  \mathtt{assert}\ \neg b)
$$
Then, it samples NSPs for each branch:
$$\mathcolor{softblue}{\mathrm{Post_{if}}}\sim \mathcal {SP}(\,\cdot\mid\,\mathcolor{softblue}{\mathrm{Pre_{if}}}, S_1),\
\mathcolor{softblue}{\mathrm{Post_{else}}}\sim \mathcal {SP}(\,\cdot\mid\,\mathcolor{softblue}{\mathrm{Pre_{else}}}, \mathrm{S_2})
$$
Finally, it uses a dedicated prompt to merge $\mathcolor{softblue}{\mathrm{Post_{if}}}$ and $\mathcolor{softblue}{\mathrm{Post_{else}}}$ into a single postcondition $\mathcolor{softblue}{\mathrm{Post}}$, which will hold after the loop. Prompts used for these operations are given in supplementary materials \Cref{if_pre_prompt,else_pre_prompt,if_else_prompt}.

\paragraph*{Loop Statements} \projName reasons about loops using our few-shot-driven $k$-induction method. For a while statement $\mathtt{while}\ b\!: S$ and a precondition $\mathcolor{softblue}{\mathrm{Pre}}$, it first samples a precondition for the loop body, and a postcondition for the first iterations: $$\mathcolor{softblue}{\mathrm{Pre_{while}}}\sim \mathcal{SP}(\,\cdot\mid\,\mathcolor{softblue}{\mathrm{Pre}}, \mathtt{assert}\ b),\ \mathcolor{softblue}{\mathrm{Post_1}}\sim \mathcal {SP}(\,\cdot\mid\,\mathcolor{softblue}{\mathrm{Pre_{while}}}, S).$$ Then, it continues to unwind the loop $k-1$ times, sequentially computing the following natural pre- and postconditions for each iteration:
$$\mathcolor{softblue}{\mathrm{Pre}_i}\sim \mathcal{SP}(\,\cdot\mid\,\mathcolor{softblue}{\mathrm{Post}_{i-1}}, \mathtt{assert}\ b),\ \mathcolor{softblue}{\mathrm{Post}_i}\sim \mathcal {SP}(\,\cdot\mid\,\mathcolor{softblue}{\mathrm{Pre}_i}, S),$$ where $i\in[2,3,...,k]$.

Finally, the NSP for the whole while statement is computed using our inductive prompt that includes the natural Hoare triples $$\Bigl\{\,\mathcolor{softblue}{\{\text{Pre}_i\}}  \, S \, \mathcolor{softblue}{\{\text{Post}_i\}}\,\Bigr\}_{i\in[1,2,...,k]}$$ as few-shot examples. This process is illustrated in \Cref{fig:kinduction_example} for a \texttt{for} statement, which is handled accordingly.

% The challenge of handling \textt{for} loops in Python is that they are implemented using iterators, and iterators store their internal state, which is updated at each iteration. Therefore, when dealing with programs containing iterators, in addition to the states of other variables, we pay extra attention to changes in the iterator itself to ensure an accurate description of program states. For convenience, beyond the existing iterator, we explicitly create an index to accompany it, which increments by 1 with each call to \texttt{next()} on the iterator, representing the iteration count. 
% \textcolor{green}{However, while the iterator state evolves through execution, the loop index is not a program statement but a counter tracking the iteration progress. Since HoarePrompt reasons about execution state changes caused by program statements, we do not treat the index as part of the SP. Instead, we maintain it as an independent counter and integrate it into the natural language description of state.}
% \textcolor{green}{Rather than embedding the index update inside SP, we maintain it as an independent counter:}
% $$
% \textcolor{green}{\mathrm{Index}_i = \mathrm{Index}_{i-1} + 1}
% $$
% \textcolor{green}{This ensures that SP remains focused on reasoning about program state changes due to statements, while index progression is explicitly tracked to describe iteration behavior.}
% This also means that in our natural language descriptions of program states, we can refer to elements of the iterator by these indices. \textcolor{green}{By combining iterator state tracking via SP and independent index tracking, we construct a complete execution state description that aligns with natural language reasoning.} Ultimately, this index, which always coexists with the iterator, will be summarized via k-induction together with other parts of the state for a loop.

Many programming languages allow iterations over a collection, e.g., the {\tt for} loop in Python. 
The challenge of handling such iteration constructs is that they are implemented using implicit iterators, which store internal states updated at each iteration. Since iterators are not explicit variables, LLM tends to not include them when producing NSPs. Therefore, when dealing with iterators, we also explicitly model the state of the iterator.

Our modeling of Python's iterators involves maintaining the index of the next element of the iterator, which increments by 1 with each call to \texttt{next()}. Our natural language descriptions of program states refer to elements from the iterator as ``the \{index\}th elements of the iterator'', where ``\{index\}'' represents our modeled index, and the description of the iterator itself is always accompanied with the description of the current index representing its state. If an iterator is used across loops, its index also needs to be summarized via k-induction together with other parts of the state.

For a loop $\texttt{for}\ \texttt{item} \ \texttt{in I}$, to accurately reflect the additional changes of both the iterator and the index for  at the start of each iteration, the precondition for the i-th iteration is sampled as:
$$
\mathcolor{softblue}{\mathrm{Pre}_i}\sim \mathcal{SP}(\,\cdot\mid\,\mathcolor{softblue}{\mathrm{Post}_i}, \text{stmt}_\mathrm{iter}),
$$
where $i\in[1,2,...,k]$, $\mathcolor{softblue}{\mathrm{Post_0}}=\mathcolor{softblue}{\mathrm{Pre}}$, $\text{stmt}_\mathrm{iter}$ is 
\[
\begin{array}{l}
\texttt{try:}\\
\quad\texttt{item = next(I\_iter})\\
\quad\texttt{index += 1}\\
\texttt{except StopIteration:}\\
\quad\texttt{break}
\end{array}
\]
where \texttt{I\_iter} is an iterator obtained from \texttt{I} and \texttt{index} is initialized to \texttt{0}. For commonly used loops over lists or ranges of elements such as \texttt{range(0,n)}, we directly refer to the index of list without modeling the semantics of \texttt{next()} for simplicity.

The supplementary materials include our prompts for computing preconditions of the loop body (\Cref{while_first_prompt,for_first_prompt}), for handling \texttt{while} and \texttt{for} loops transitions (\Cref{while_next_prompt,for_next_prompt}), and for inductive inference (\Cref{total_while_prompt,total_for_prompt}).

% \begin{figure}[t]
%     \centering
%     \includegraphics[width=0.9\linewidth]{example_iter_blue.png}
%     \caption{Computing NSP for the iterator in a \texttt{for} statement.\label{fig:iter-example-label}}
%     \vspace{-2mm}
% \end{figure}

% \textcolor{yellow}{HERE NEW INFO ON ITERATORS}
% \textbf{Iterator State and Side-Effects with \texttt{next()}}

% Consider the following Python iterator loop example:

% \begin{lstlisting}[language=python, basicstyle=\ttfamily\footnotesize]
% iterator = iter(lst) 
% while True:
% num = next(iterator, None) 
% if num is None:
% break
% # loop body omitted
% \end{lstlisting}

% This example illustrates the explicit tracking of an iterator's internal state through each iteration, highlighting the irreversible consumption of elements by the \texttt{next()} function:

% \begin{itemize}
% \item \textbf{Precondition before the first \texttt{next()} execution:} \texttt{lst} is a list of integers with at least one element, \texttt{iterator} points at the first element of \texttt{lst}.

% \item \textbf{Postcondition after the first \texttt{next()} execution:} \texttt{lst} is unchanged, \texttt{num} contains the first element of \texttt{lst}, and \texttt{iterator} points at the second element of \texttt{lst} if it exists, otherwise points to \texttt{None}.

% \item \textbf{Precondition before the second \texttt{next()} execution:} \texttt{lst} is unchanged, \texttt{iterator} points at the second element of \texttt{lst}, list has at least two elements.

% \item \textbf{Postcondition after the second \texttt{next()} execution:} \texttt{lst} is unchanged, \texttt{num} contains the second element of \texttt{lst}, and \texttt{iterator} points at the third element of \texttt{lst} if it exists, otherwise points to \texttt{None}.


% \item \textbf{Final state after loop exit condition:} \texttt{lst} is unchanged, \texttt{num} is \texttt{None}, \texttt{iterator} points to \texttt{None}, indicating that no more elements remain in the iterator.

% \end{itemize}

% This explicit state tracking highlights the side effects of \texttt{next()} calls, ensuring that preconditions and postconditions accurately reflect changes to the iterator state across iterations.

\begin{figure}[t]
    \centering
    \includegraphics[width=\linewidth]{example_stdin_blue.png}
    \vspace{-6mm}
    \caption{Computing NSP for the call of \texttt{input()} for \texttt{stdin}.\label{fig:stdin-example-label}}
    \vspace{-4mm}
\end{figure}

\paragraph{Exceptions Handling} Given $\mathcolor{softblue}{\mathrm{Pre}}$, and a \texttt{try-except} statement $$\text{try}\!:\alpha\ \text{except}\ E_1\!: \beta_1\ \text{except}\ E_2\!: \beta_2,\ldots\,\text{except}\ E_m\!: \beta_m,$$ we first sample the NSP for the \texttt{try}-block: $$\mathcolor{softblue}{\mathrm{Post_{try}}}\sim \mathcal{SP}(\,\cdot\mid\,\mathcolor{softblue}{\mathrm{Pre}}, \alpha),$$ and then the NSPs for different exception handling blocks: $$\mathcolor{softblue}{\mathrm{Post_i}}\sim \mathcal {SP}(\,\cdot\mid \mathcolor{softblue}{\text{``variables may take arbitrary values"}}, \beta_i).$$ 

Finally, we merge $\mathcolor{softblue}{\mathrm{Post_{try}}},\mathcolor{softblue}{\mathrm{Post_1}},\mathcolor{softblue}{\mathrm{Post_2}},\ldots,\mathcolor{softblue}{\mathrm{Post_m}}$ using a prompt provided in the supplementary materials (\Cref{try_except_prompt}). Note that, due to the complexity of handling nonlinear control flow, the NSP computation for exception handling blocks does not inherit any preconditions here, assuming that the state may be arbitrary.

% \begin{align}
% \mathrm{sp}&\left(, \text{stmt}_{try}\right) \notag\\
% &= \underbrace{\operatorname{HP}(P_{\text{NL}}, \alpha)}_{\text{Normal Path}} \ \lor\ \underbrace{\bigvee_{i} \operatorname{HP}\left(P_{\text{NL}} \land \text{Raised}(E_i), \beta_i\right)}_{\text{Exception Paths}}
% \end{align}

\paragraph*{Input Stream} \Cref{fig:stdin-example-label} shows how data from the standard input (\texttt{stdin}) is processed by \projName. When programs read from \texttt{stdin}, the precondition extraction explicitly describes the expected structure and content of the \texttt{stdin} based on the given problem description. So, the precondition would clearly indicate that \texttt{stdin} contains these inputs in that specific order. Every time a postcondition for a statement reading $\texttt{stdin}$ is sampled as $$\mathcolor{softblue}{\mathrm{Post}}\sim \mathcal{SP}(\,\cdot\mid\,\mathcolor{softblue}{\text{Pre}}, \texttt{x\ = input()}),$$ the description of the stream state in $\mathcolor{softblue}{\mathrm{Post}}$ will be updated accordingly. After all data is read, the state of the stream is described as ``stdin is empty''. Note that this modeling applies only to non-interactive programs, where \texttt{stdin} is predetermined and does not depend on the execution of the program.

% \textcolor{green}{Initial Problem Description:
% The program first reads two numbers from the user: the first number represents the user's age, and the second represents their weight.
% Initial Precondition:
% \texttt{stdin} contains two inputs: the user's age (on top), followed by their weight.
% Postcondition after executing x = input():
% Variable x now contains the user's age. \texttt{stdin} contains one input: the user's weight.
% Postcondition after executing y = input():
% Variable x contains the user's age, variable y contains the user's weight, and \texttt{stdin} is empty.}

%% \textcolor{purple}{Overall, \projName\ does not just replace a symbolic strongest postcondition  function with a single LLM query in natural language. Rather, it leverages \emph{multiple, targeted prompts} to mirror every step of Hoare-style reasoning while capitalizing on the expressive power of NL. Specifically, each syntactic construct—assignments, compound statements, branching, loops—has a \emph{dedicated prompt} that refines the state accordingly.  We generate \emph{intermediate NL postconditions} at each control-flow edge by updating only the pieces of state actually changed by the code snippet in question. In an \texttt{if-else}, for example, \projName\ uses one prompt to strengthen the precondition by “assuming the condition is true” (thus describing that branch’s postcondition) and another prompt to handle the negation for the else branch; these are later merged into a single updated state. For loops, it applies a series of prompts that iteratively specify “what must hold for the loop to continue,” effectively unrolling or applying 
%% k-induction in natural language. This pipeline preserves compositional structure, so that each code fragment updates only the relevant variables in the “NL postcondition” while reusing the rest of the prior description. Staying in natural language lets us capture \emph{complex real-world semantics} in simpler terms than raw logical formulas. For instance, whereas code might check multiple divisibility conditions to confirm a “leap year,” \projName\ can just incorporate the notion “year is a leap year” as a single textual constraint. Because we keep these constraints in plain text, we neither lose valuable domain-specific nuances nor need to encode them manually as complex first-order predicates. In short, \projName\ leverages step-by-step prompts \emph{and} the flexibility of NL to yield a thorough, human-readable record of every control-flow path—one that can seamlessly incorporate advanced verification logic and higher-level user-defined concepts.}

\subsection{Correctness Classification}
\label{sec:classification}

Let \( R \) be the natural language requirement, \( P \) be the program to be evaluated, \( D = \{(R_1, P_1, y_1), (R_2, P_2, y_2), \dots, (R_n, P_n, y_n)\} \) be a labeled dataset, where \( y_i \in \{ \textrm{Correct}, \textrm{Incorrect} \} \) represents whether program \( P_i \) satisfies requirement \( R_i \). The problem of classifying correctness w.r.t. natural language requirements can be formalized as the task of learning a function \( f \colon (R, P) \to \{\textrm{Correct}, \textrm{Incorrect}\} \).

In this work, we investigate several classification functions based on NSP, as well as alternative ideas.

{\textbf{\projName}:} Given $R$ and $P$, \projName first annotates key program points with descriptions of reachable states computed using $\mathcal{SP}$. Specifically, it adds annotations as comments after top-level if and loop statements. Apart from that, it adds descriptions of returned and printed values after return and print statements, as these values are often used to assess correctness. An example of an annotation program is given in \Cref{fig:motivating_example}. The annotated program is passed alongside $R$ to an LLM via our classification prompt. The supplementary materials contain the classification prompts for single- and multi-function programs (\Cref{correctness-single-func-prompt,correctness-multi-func-prompt}).

{\textbf{\projName (No Unroll)}: } We also implemented a variant of \projName that employs a non-compositional approach to computing the NSP for loops by prompting an LLM to infer the postcondition of the entire loop directly within a single LLM query. The prompts applied in this approach are given in the supplementary materials (\Cref{while_NO_UNROLL_PROMPT,for_NO_UNROLL_PROMPT}).

{\textbf{Vanilla}: } Given $R$ and $P$, it directly asks an LLM to classify program correctness using the prompt given in \Cref{vanilla-prompt}.

\textbf{\textbf{Zero-shot-CoT}:} Our Zero-shot-CoT~\cite{kojima2022large} prompt extends the Vanilla by triggering chain-of-thought reasoning~\cite{wei2022chain} by adding the ``think step-by-step'' instruction (\Cref{0-shot-COT-prompt}).

\textbf{\textbf{Tester}:} As an alternative to static reasoning, we also employed testing for classifying program correctness. Specifically, this method first generates tests based on $R$ with an LLM using the testing prompt adapted from AgentCoder~\cite{huang2023agentcoder}, executes $P$ on the generated tests, and classifies it as correct iff it passes the tests.

\section{Evaluation}
\label{sec:evaluation}

Our experiments address these research questions:
\begin{description}
\item[RQ1:] How does the classification performance of \projName compare to other techniques?
\item[RQ2:] How does few-shot-driven $k$-induction impact HoarePrompt’s effectiveness?
\item[RQ3:] How does the LLM token use of \projName compare to other techniques?  
\changes{\item[RQ4:] How effective is \projName as a feedback mechanism in a downstream code generation task?}
\end{description}

We selected four representative open-weight LLMs. Three come from the  Qwen~\cite{qwen} family: Qwen2.5-7B-Instruct, a popular smaller model, Qwen2.5-Coder-32B-Instruct, a model specialized for coding tasks, and Qwen2.5-72B-Instruct, demonstrating performance comparable to GPT-4o-mini~\cite{llm-stats}. \changes{Another one is from the Llama family: Llama3.1-70b-Instruct.} \changes{To measure classification quality, we adopted the Matthews Correlation Coefficient (MCC) as it considers all four components of the confusion matrix, and produces more reliable statistical results than F1 score~\cite{mcc}. Further information on MCC is available in supplementary materials (Appendix~\ref{appendix:mcc}.)} We also report other standard metrics such as Balanced Accuracy, True Positive Rate (TPR), False Negative Rate (FNR), and False Positive Rate (FPR). \changes{To ensure robustness against LLM non-determinism, we repeated each experiment multiple times. Following best practices~\cite{ouyang2025empirical,xiong2024llmsexpressuncertaintyempirical}, we calibrated the rerun count based on model confidence and consistency. This yielded 3 reruns per experiment. Details of this derivation are included in supplementary materials (Appendix~\ref{appendix:reruns})}. All experiments were conducted in a virtual machine with a AMD EPYC 7543 64-bit CPU, two cores and 16GB of RAM.

\begin{table}[t]
\centering
\caption{\changes{Effect of unroll bound $k$ on MCC and token usage (normalized over $k{=}0$).}}
\vspace{-2mm}
\label{tab:k-ablation}
\begin{tabular}{@{}lcc@{}}
\toprule
\changes{\textbf{Unroll Bound $k$}} & \changes{\textbf{MCC}}  & \changes{\textbf{Token Increase}} \\
\midrule
\changes{0 (no unroll)} &  \changes{0.211}  & \changes{1.00$\times$} \\
\changes{1} & \changes{0.228} & \changes{3.02$\times$} \\
\changes{3} & \changes{\textbf{0.262}}  & \changes{11.29$\times$} \\
\changes{5} & \changes{0.231} & \changes{24.53$\times$} \\
\bottomrule
\end{tabular}
\vspace{-2mm}
\end{table}


\begin{table*}[t]
    \centering
    \caption{Performance of classifiers of program correctness w.r.t. natural language requirements and their token use across four LLMs evaluated on \benchName. \projName significantly improves the MCC in comparison with baselines at the expense of a higher token use. \projName without loop unrolling represents a trade-off between classification quality and token cost. }
    \label{tab:results}
   \begin{tabular}{|c|l|p{0.9cm}p{1.2cm}p{0.9cm}p{0.9cm}p{0.9cm}p{0.9cm}|ccc|}

        \hline
        \multirow{2}{*}{\textbf{Model}} & \multirow{2}{*}{\textbf{Classifier}} & \multicolumn{6}{c|}{\textbf{Correctness Classification Quality Metrics}} & \multicolumn{3}{c|}{\textbf{Tokens (Millions)}} \\
        \cline{3-11}
        & & \textbf{MCC} & $\Delta\mathbf{MCC}$ & \textbf{BA} & \textbf{TPR} & \textbf{FNR} & \textbf{FPR} & \textbf{IT} & \textbf{OT} & \textbf{TT} \\
        \hline
       \multirow{5}{*}{\parbox{1.8cm}{\centering Qwen2.5-7B}}
         & Vanilla       & 0.055 & \textcolor{red}{  } & 0.521 & 0.193 & 0.807 & 0.151 & 1.8 & 0.8 & 2.6 \\
        & Zero-shot-CoT  & 0.085 & \textcolor{darkgreen}{+54.5\%}  & 0.543 & 0.563 & 0.437 & 0.478 & 2.1 & 1.3 & 3.4 \\
       
        & Tester             & 0.109 & \textcolor{darkgreen}{+98.18\%} & 0.519 & 0.051 & 0.949 & 0.012 & 1.7 & 3.8 & 5.5 \\
        \cline{2-11}
        & \textbf{HoarePrompt}         & \textbf{0.159} & \textcolor{darkgreen}{\textbf{+189.1\%}} & 0.577 & 0.700 & 0.300 & 0.546 & 205.1 (93.9) & 31.9 & 237.1 (125.9) \\
        & HoarePrompt no unroll & 0.119 & \textcolor{darkgreen}{+116.36\%} & 0.558 & 0.666 & 0.334 & 0.550 & 19.9 (12.3) & 3.4 & 23.4 (15.7) \\
        \hline
        \hline
        \multirow{5}{*}{\parbox{1.8cm}{\centering Qwen2.5-Coder-32B}}
        & Vanilla       & 0.137 & \textcolor{red}{ } & 0.553 & 0.233 & 0.767 & 0.128 & 1.9 & 0.9 & 2.8 \\
        & Zero-shot-CoT  & 0.123  & \textcolor{red}{-10.22\%} & 0.560 & 0.438 & 0.562 & 0.318 & 2.4 & 1.5 & 3.9 \\
        & Tester  & 0.096 &\textcolor{red}{-29.93\%}  & 0.518 & 0.055 & 0.945 & 0.019 & 1.9 & 3.3 & 5.2 \\
        \cline{2-11}
        & \textbf{HoarePrompt}         & \textbf{0.231} & \textcolor{darkgreen}{\textbf{+68.61\%}}  &0.616 & 0.593 & 0.407 & 0.362 & 186.5 (88.2) & 42.6 & 229.1 (130.8) \\
        & HoarePrompt no unroll & 0.158 &\textcolor{darkgreen}{+15.04\%}  & 0.579 & 0.571 & 0.429 & 0.412 & 19.2 (11.6) & 4.4 & 23.6 (15.9) \\
        \hline
        \hline
        \multirow{5}{*}{\parbox{1.8cm}{\centering Qwen2.5-72B}} 
        & Vanilla       & 0.169  & \textcolor{red}{ }  & 0.578 & 0.382 & 0.618 & 0.226 & 1.7 & 0.9 & 2.6 \\
        & Zero-shot-CoT  & 0.193 & \textcolor{darkgreen}{+14.20\%}  & 0.593 & 0.714 & 0.286 & 0.527 & 2.0 & 1.3 & 3.2 \\
        & Tester             & 0.131  &\textcolor{red}{-22.49\%}  & 0.522 & 0.052 & 0.948 & 0.007 & 1.9 & 3.8 & 5.6 \\
        \cline{2-11}
        & \textbf{HoarePrompt}         & \textbf{0.259} &  \textcolor{darkgreen}{\textbf{+53.25\%}} & 0.629 & 0.667 & 0.333 & 0.409 & 163.5 (78.7) & 32.5 & 196.0 (111.2) \\
        & HoarePrompt no unroll & 0.231 & \textcolor{darkgreen}{+36.69\%}  & 0.615 & 0.618 & 0.382 & 0.387 & 19.3 (11.8)  & 3.8 &  23.1 (15.6) \\
        
         \hline
        % \hline
        % \multirow{5}{*}{\parbox{1.8cm}{\centering Llama3.1-70b}} 
        % & Vanilla       & 0.161  & \textcolor{red}{ }  & 0.579 &  0.671 & 0.329 & 0.513 & 1.7 & 0.01 & 1.7 \\
        % & Zero-shot-CoT  & 0.157 & \textcolor{red}{-2.48\%}  & 0.575 & 0.728 & 0.329 & 0.514 & 1.8 & 0.7 & 2.5 \\
        % & Tester             & 0.099  &\textcolor{red}{-38.51\%}  & 0.527 & 0.109 & 0.891 & 0.055 & 1.7 & 2.4 & 4.1 \\
        % \cline{2-11}
        % & \textbf{HoarePrompt}         & \textbf{0.250} &  \textcolor{darkgreen}{\textbf{+55.28\%}} & 0.601 & 0.895 & 0.105 & 0.692 & 144.3 (106.8) & 27.3 & 171.6 (134.1) \\
        % & HoarePrompt no unroll & 0.209 & \textcolor{darkgreen}{+29.81\%}  & 0.588 & 0.858 & 0.142 & 0.681 & 12.0 (9.5)  & 1.7 &  13.7 (11.2) \\
        % \hline
        \multirow{5}{*}{\parbox{1.8cm}{\centering \changes{Llama3.1-70b}}} 
    & \changes{Vanilla}       & \changes{0.161}  & \textcolor{red}{ }  & \changes{0.579} &  \changes{0.671} & \changes{0.329} & \changes{0.513} & \changes{1.7} & \changes{0.01} & \changes{1.7} \\
    & \changes{Zero-shot-CoT}  & \changes{0.157} & \textcolor{red}{-2.48\%}  & \changes{0.575} & \changes{0.728} & \changes{0.329} & \changes{0.514} & \changes{1.8} & \changes{0.7} & \changes{2.5} \\
    & \changes{Tester}             & \changes{0.099}  & \textcolor{red}{-38.51\%}  & \changes{0.527} & \changes{0.109} & \changes{0.891} & \changes{0.055} & \changes{1.7} & \changes{2.4} & \changes{4.1} \\
    \cline{2-11}
    & \textbf{\changes{HoarePrompt}}         & \textbf{\changes{0.250}} &  \textcolor{darkgreen}{\textbf{+55.28\%}} & \changes{0.601} & \changes{0.895} & \changes{0.105} & \changes{0.692} & \changes{144.3 (106.8)} & \changes{27.3} & \changes{171.6 (134.1)} \\
    & \changes{HoarePrompt no unroll} & \changes{0.209} & \textcolor{darkgreen}{+29.81\%}  & \changes{0.588} & \changes{0.858} & \changes{0.142} & \changes{0.681} & \changes{12.0 (9.5)}  & \changes{1.7} &  \changes{13.7 (11.2)} \\
\hline

    \end{tabular}
   \caption*{\small \textbf{Abbreviations:} MCC = Matthews Correlation Coefficient, $\Delta\mathbf{MCC}$: Relative MCC Improvement over Vanilla , BA = Balanced Accuracy, TPR = True Positive Rate, FNR = False Negative Rate, FPR = False Positive Rate, IT = Input Tokens, OT = Output Tokens, TT = Total Tokens. Parentheses in IT and TT for HoarePrompt classifiers indicate token counts excluding few-shot learning examples, which can be eliminated via fine-tuning.}
   \vspace{-5mm}
\end{table*}


\subsection{\benchName Dataset}

Although there exist numerous datasets of natural language requirements~\cite{apps,mbpp}, some including incorrect implementations~\cite{tan2017codeflaws}, we could not use most of them due to potential data leakage, since memorization of correct solutions will significantly impact LLM reasoning~\cite{xie2024memorization,deng-etal-2024-unveiling}. Apart from that, we required non-trivial requirements and subtle bugs to thoroughly evaluate the ability of LLMs to reason about program correctness. We constructed a benchmark \benchName (\underline{Co}de \underline{Co}rrectness \underline{Cla}ssification w.r.t. \underline{N}atural \underline{L}anguage requirements). It consists of 645 problem-solution pairs from Codeforces programming contests that took place in the first half of 2024, with solutions labelled as correct or incorrect by the Codeforces' internal testing system. Totally, 322 of the solutions (49.9\%) are incorrect. All problems and solutions are released since January 2024, after the cut-off data of many SOTA LLMs, and the problems are challenging, have unambiguous descriptions, and solutions to these problems are rigorously tested to find subtle bugs~\cite{huang-etal-2024-competition}. To capture subtle and challenging coding errors, we specifically included incorrect submissions that immediately preceded correct ones, hypothesizing that these would likely contain nuanced mistakes. \changes{Further details about the dataset are provided in supplementary materials (\Cref{dataset_details})}.


\subsection{Hyperparameters}
\label{sec:hyperparameters}
\changes{We tuned the unrolling bound $k \in \{0,1,3,5\}$ using \texttt{Llama3.1-70B} on a 140-entry pilot set randomly sampled from the same source as \benchName, but disjoint from the 645 \benchName programs. As shown in \Cref{tab:k-ablation}, $k{=}3$ achieves the best MCC (0.262), balancing generalization and token cost. Higher $k$ increase token usage while degrading the quality. Notably, no matter the choice of $k$ \projName outperforms the baseline classifiers by at least 23.7\%.}

\changes{Another key design choice is how many and which annotations to include in the final prompt. We mark only major program points—after loops, branches, returns, prints, and try-catch blocks—to reduce prompt size while maintaining clarity. In pilots, full-state annotation improved MCC over Zero-shot-CoT by 19.7\%, whereas our sparse strategy achieved 62\%. We hypothesize that excessive detail dilutes focus, while key transitions aid reasoning.}



\subsection{RQ1: Classification Performance}

In the first research question, we examine whether structural reasoning with descriptions of reachable states helps \projName improve  LLMs' ability to evaluate program correctness compared to other techniques. By systematically  tracking program states through \projName, we hypothesize that structured reasoning will enhance correctness classification.

% To measure classification quality, we adopted the Matthews Correlation Coefficient (MCC) as it considers all four components of the confusion matrix, and produces more reliable statistical results than F1 score~\cite{mcc}. For information on the advantages of MCC and its calculation refer to Appendix \ref{appendix:mcc}. Besides MCC, we also report other standard classification metrics such as Balanced Accuracy, True Positive Rate (TPR), False Negative Rate (FNR), and False Positive Rate (FPR). We say that a classification result is a true positive if a correct program is classified as correct. (moved this earlier)

\Cref{tab:results} shows the results for each model. In comparison with LLM-based baselines, Vanilla and Zero-shot-CoT, \projName consistently demonstrated higher classification quality. On average, \projName improves the MCC by 61\% over Zero-shot-COT and 72\% over Vanilla. The impact of \projName is most pronounced in smaller models. The most significant improvement of MCC is in Qwen2.5-7B, being 87\% over Zero-shot-COT and 189\% over Vanilla. For Qwen2.5-72B, MCC improves by 34\%  and 53\% over the Zero-shot-CoT and Vanilla respectively. \changes{The results of the Llama model follow similar trends to those exhibited by the Qwen2.5 models, being especially close to Qwen2.5-72b, which is of a similar size.}

Similar improvement is observed across other classification metrics. \projName's TPR is 15.5\% higher than Zero-shot-COT and 107.1\% higher than Vanilla. \projName also reduces the FNR compared to Zero-shot-COT and Vanilla, The FPR for \projName (0.502), is also just slightly higher than Zero-shot-COT (0.459) but significantly higher than Vanilla (0.254).

% \projName also outperformed the Tester classifier. On average, its MCC is 106\% higher than that of Tester. Moreover, Tester exhibits an extremely low True Positive Rate of only 0.067, meaning it correctly identifies just 6.7\% of correct programs. This suggests that the expected output oracles inferred from the requirements are incorrect, leading to misclassifications of correct programs as incorrect.  and analyzed in more detail in Subsection \ref{tester_analysis}.

\projName also significantly outperforms the \texttt{Tester} classifier, increasing MCC by 106\% on average. Notably, Tester achieves a very low average TPR of only 6.7\%, indicating frequent misclassifications of correct programs. This issue stems from test generation inaccuracies: LLMs often fail to produce correct output oracles, leading to false negatives. \changes{\Cref{tab:tester-stats} shows the breakdown of Tester behavior. On average, each program was evaluated with 13–28 test cases, but 15--25\% of examples had no passing test cases, and only 3--8\% had no failing test. Thus, the program was marked incorrect in over 90\% of cases, including the vast majority of correct programs, due to at least one failing test case. This affected all model sizes, showcasing a limitation of test-based correctness classification. A qualitative analysis of the misclassifications of \projName and Tester are presented in \Cref{qualitative}.}

\begin{table}[t]
\centering
\caption{\changes{Tester statistics across four LLMs. “Avg. Total” and “Avg. Passed” --- the average number of generated/passed tests per program. “0 Pass \%” and “0 Fail \%” --- the percentage of programs for which all tests failed or all tests passed.}}
\label{tab:tester-stats}
\begin{tabular}{@{}lcccc@{}}
\toprule
\changes{\textbf{Model}} & \changes{\textbf{Avg. Total}} & \changes{\textbf{Avg. Passed}} & \changes{\textbf{0 Pass \%}} & \changes{\textbf{0 Fail \%}} \\
\midrule
\changes{Qwen2.5-7B} & \changes{27.81} & \changes{9.99} & \changes{23.17\%} & \changes{3.32\%} \\
\changes{Qwen2.5-32B} & \changes{17.30} & \changes{7.86} & \changes{15.14\%} & \changes{3.63\%} \\
\changes{Qwen2.5-72B} & \changes{17.47} & \changes{6.50} & \changes{16.90\%} & \changes{3.13\%} \\
\changes{LLaMA3.1-70B} & \changes{12.72} & \changes{5.02} & \changes{25.34\%} & \changes{8.19\%} \\
\bottomrule
\end{tabular}
\vspace{-5mm}
\end{table}

\begin{tcolorbox}[colback=gray!5, colframe=black!20, arc=2pt, boxrule=0.3mm, breakable, sharp corners, left=2pt, right=2pt, top=2pt, bottom=2pt]
\textbf{RQ1:} \textit{\projName} consistently outperforms the baselines in terms of classification quality across all LLMs, achieving an average improvement of \textbf{61\%} in MCC over Zero-shot-COT, \textbf{72\%} over Vanilla and \textbf{106\%} over Tester.
\end{tcolorbox}

%These results confirm that incorporating structured reasoning into program correctness assessment substantially enhances an LLM’s ability to assess correctness, particularly benefiting smaller models with weaker inherent reasoning capabilities.



% \subsection{Discussion of Tester classifier limitations}
% \label{tester_analysis}
% To further understand the limitations of using test-based classification, we analyzed the behavior of the \texttt{Tester} classifier across four models. Table~\ref{tab:tester-stats} presents key statistics per evaluation instance, including average total tests, average passed tests, and the percentage of evaluations with zero passing or zero failing tests. The “0 pass” percentage reflects how often no passing test was produced, while “0 fail” indicates how often no failing test was generated.


% These results show that the percentage of evaluations with no passing test case is relatively low (15--25\%), suggesting that even for challenging problems, LLMs are often capable of generating some simple test cases that reflect the core functionality of the task. Many incorrect programs pass these simple tests, indicating that their failures typically stem from more complex edge cases.

% Conversely, the low “0 fail” rates (3--8\%) reveal that over 90\% of the time at least one test case is marked as failing, even for correct code. This underscores the difficulty LLMs face in predicting accurately the output for complex inputs. The trend holds for both smaller and larger models, demonstrating a fundamental limitation in test generation for nuanced semantic correctness.



\subsection{RQ2: Impact of Inductive Reasoning}

Loops introduce complexity for reasoning about programs. \projName mitigates this complexity by inductively reasoning about loops using our novel few-shot-driven k-induction technique. This technique involves unrolling loops k times, with $k=3$ in our experiments, before generalizing the final state.

This question investigates whether loop unrolling  enhances correctness classification quality. To measure the impact of our few-shot-driven k-induction algorithm, we compare \projName's classification performance with that of \projName (no unroll) that summarizes loops within a single step.

As shown in \Cref{tab:results}, the few-shot-driven $k$-induction technique provides significant additional performance gains, improving the MCC by an average of 25.7\% over the non-inductive version (\textit{\projName} no unroll). Specifically, it improves MCC by 33.6\% in Qwen2.5-7B, 45.3\% in Qwen2.5-Coder-32B, 12.1\% in Qwen2.5-32B  \changes{and 19.6\% for LLama3.1-70b.} Even without unrolling, \projName demonstrates the second-best result among the studied classifiers.

\begin{tcolorbox}[colback=gray!5, colframe=black!20, arc=2pt, boxrule=0.3mm, breakable, sharp corners, left=2pt, right=2pt, top=2pt, bottom=2pt]
\textbf{RQ2:} Few-shot-driven k-induction enhances classification performance across all models, improving the MCC by an average of \textbf{25.7\%} over the non-inductive version of \projName.
\end{tcolorbox}

%These results suggest that reasoning about loops is a challenge for LLMs, and few-shot-driven k-induction effectively alleviates this challenge by inferring more precise loop postconditions.

\subsection{RQ3: Token Usage}

Since \projName involves multiple LLM queries, it consumes more tokens than direct LLM prompts~\cite{liu2024groupdebate}. We quantified this increase and investigated the trade-off between token use and classification quality. Table~\ref{tab:results} provides the total token usage across classifiers and models for the 645 problem-solution pairs of the \benchName. For each classifier, we measured the input tokens (IT), the output tokens (OT), and total tokens (TT). As explained in \Cref{sec:nsp}, we chose to implement our natural strongest postcondition computation using prompt engineering with few-shot learning examples due to budget constraints. Tokens representing these hard-coded examples are repeatedly sent to an LLM during the algorithm iterations; in principle, they can be eliminated via fine-tuning. Because of that,  for the \projName variants, we also report the number of ``essential'' input and total token, in parentheses, that excludes hard-coded few-shot example tokens.

On average, \projName consumes 85.9 times more tokens than Vanilla (51 times when excluding non-essential tokens), 64.1 times more tokens than Zero-shot-CoT (38.6 times more when excluding non-essential tokens), and  40.8  times more tokens than Tester (24.6 times more, when excluding non-essential tokens).

A key reason for the increased token usage is our inductive reasoning algorithm (few-shot-driven $k$-induction) that performs loop unrolling. On average, it requires 10.9 times extra tokens (8.6 times extra when excluding non-essential tokens) than the non-inductive variant. Thus, the few-shot k-induction is more suitable for situations when classification accuracy is prioritized over cost. The non-inductive version, \projName no unroll, represents a practical trade-off by retaining a substantial performance boost (27.8\% better than the top performing approach, Zero-shot-COT) while consuming 90\% fewer tokens than full unrolling. Overall, it consumes only 6.4 times more tokens (4.5 times more, when excluding non-essential tokens) than Zero-shot-COT.

\begin{tcolorbox}[colback=gray!5, colframe=black!20, arc=2pt, boxrule=0.3mm, breakable, sharp corners, left=2pt, right=2pt, top=2pt, bottom=2pt]
\textbf{RQ3:} \projName's significant performance boost comes at the expense of an average token usage increase of 64 times (39 times when excluding non-essential tokens that can be eliminated via fine-tuning) compared to the most accurate baseline, Zero-shot-COT. The version of \projName without unrolling represents a practical trade-off between classification quality and cost, demonstrating the second-best performance, while consuming 90\% fewer tokens than the main variant.
\end{tcolorbox}

\subsection{RQ4: Code generation}
\label{sec:codegen}

\changes{To assess \projName as a feedback mechanism in the context of code generation, we adapted the AgentCoder~\cite{huang2023agentcoder} pipeline and implemented two variants: one using test-based feedback (as in the original), and one replacing test failures with correctness judgments from \projName. The third baseline generates a final solution in a single step without refinement. We evaluated all three methods on 100 problems randomly sampled from \benchName. We excluded problems with multiple valid outputs --- cases where different programs produce acceptable but non-identical outputs --- since correctness becomes ambiguous and not reliably testable~\cite{quan2025codeelobenchmarkingcompetitionlevelcode}.}

\changes{For our main experiments, we used \texttt{Llama3.1-70b-instruct}. Due to the complexity of \benchName and to evaluate generalization on more capable models, we additionally employed the stronger \texttt{GPT-4.1} model.} \changes{Results for both models are shown in \Cref{tab:agentcoder_results_combined}.}


\begin{table}[t]
\centering
\caption{\changes{Code generation results on 100 problems from \benchName. 
\textbf{Correct\%}: percentage of programs passing all hidden tests; 
\textbf{ATP}: average fraction of test cases passed.}}
\label{tab:agentcoder_results_combined}
\begin{tabular}{lcccc}
\toprule
\multirow{2}{*}{\changes{\textbf{Method}}} & 
\multicolumn{2}{c}{\changes{\textbf{Llama3.1}}} & 
\multicolumn{2}{c}{\changes{\textbf{GPT-4.1}}} \\
\cmidrule(lr){2-3} \cmidrule(lr){4-5}
 & \changes{\textbf{Correct\%}} & \changes{\textbf{ATP}} 
 & \changes{\textbf{Correct\%}} & \changes{\textbf{ATP}} \\
\midrule
\changes{\projName} & \changes{\textbf{33\%}} & \changes{\textbf{57.98\%}} & \changes{\textbf{47\%}} & \changes{\textbf{69.55\%}}\\
\changes{Testing}   & \changes{29\%} & \changes{54.89\%} & \changes{42\%} & \changes{69.53\%}\\
\changes{Baseline}  & \changes{30\%} & \changes{56.34\%} & \changes{37\%} & \changes{62.03\%}\\
\bottomrule
\end{tabular}
\end{table}



%  With the smaller Llama3.1-70b model, the testing-based feedback sometimes introduced regressions by generating incorrect oracles that misled the refinement process. This slightly reduced both the percentage of fully correct programs and the average tests passed compared to the baseline. In contrast, HoarePrompt’s structured correctness analysis provided modest but consistent improvements over both baseline direct prompting and test-driven refinement, highlighting its promise as a more reliable feedback mechanism. We provide an example of such a regression in \ref{ex_agent_coder_test} and an improvement of Hoareprompt over the baseline in \ref{ex_hoareprompt_beautiful_pairs}.
% With GPT-4.1, which is significantly stronger, the picture changes: test-based feedback now improves over the baseline on both metrics, showing that given a sufficiently capable model, iterative test refinement becomes more robust. However, \projName continues to yield further gains—achieving the highest fraction of fully correct programs (47\%) and the best average test pass rate (69.55\%).
% \
% Interestingly, the relatively small difference in average tests passed between \projName and the testing-based approach with GPT-4.1 (69.55\% vs.~69.53\%), alongside a clearer jump in the percentage of completely correct solutions, suggests that \projName helps catch subtle edge cases that might not drastically change the bulk of tests passed but still determine full correctness.
\changes{With Llama3.1-70B, the testing-based variant occasionally regressed due to incorrect oracles, reducing both fully correct programs and test pass rates compared to the baseline. In contrast, \projName consistently improved both metrics, highlighting its robustness. The supplementary materials include a regression example (\Cref{ex_agent_coder_test}) and \projName's improvement (\Cref{ex_hoareprompt_beautiful_pairs}). With the stronger GPT-4.1, the test-based variant improves over the baseline --- suggesting that test-driven refinement benefits from stronger models. However, \projName still outperforms both, achieving the highest correct rate (47\%) and best ATP (69.55\%).}

%\changes{Although the ATP of \projName and test-based refinement are nearly identical (69.55\% vs.~69.53\%), the larger gain in fully correct programs suggests that \projName captures subtle edge cases missed by test oracles, improving end-to-end correctness.}

\begin{tcolorbox}[colback=gray!5, colframe=black!20, arc=2pt, boxrule=0.3mm, breakable, sharp corners, left=2pt, right=2pt, top=2pt, bottom=2pt]
\changes{\textbf{RQ4:} \projName is effective as a feedback mechanism for iterative code generation, increasing the number of solved problems by 12--14\% compared to using generated tests as feedback.}
\end{tcolorbox}


% \subsection{Qualitative Error Analysis}
% \label{qualitative}

% \begin{figure*}[t]
%   \begin{center}
%     \includegraphics[width=\textwidth]{tester_example.png}
%   \end{center} 
%   \caption{\changes{A problem and its correct solution from \benchName. Vanilla consistently misclassifies the program as incorrect, as it fails to correctly interpret the algorithm. Tester generates tests involving incorrect outputs, as LLMs struggle to predict results of non-trivial computations. \projName's annotations enable it correctly interpret the algorithm to judge its correctness.\label{fig:tester_example}}}
%   \vspace{-2mm}
% \end{figure*}

  
% \changes{To better understand the limitations of \projName and Tester, we manually analyzed a subset of misclassifications. For \projName, the most common sources of misclassifications are}

% \begin{itemize}[leftmargin=*]
%     \item \changes{\textbf{Missing Edge Case:}} \changes{The NSP chain failed to capture critical edge behavior, leading to overly optimistic judgments.}

%     \item \changes{\textbf{Oversimplified Annotations:}} \changes{The annotations glossed over subtle state distinctions, reinforcing a false sense of correctness.}

%     \item \changes{\textbf{Incorrect Loop Generalization:}} \changes{The model incorrectly extrapolated loop behavior, producing faulty postconditions.}

% %    \item \changes{\textbf{(H4) Semantic Trust Bias:}} \changes{When code structure closely resembled the problem description even if the annotations described different behavior,], the LLM sometimes ignored those contradictory annotations and falsely assumed correctness.}

%     \item \changes{\textbf{Interprocedural Failures:}} \changes{\projName made reasoning mistakes across multiple functions, leading to misjudgments.}

%     \item \changes{\textbf{Variable Scope/Shadowing Errors:}} \changes{In loop-heavy code, scope bleed or reused loop indices caused postcondition inaccuracies.}

% %    \item \changes{\textbf{(H7) Description-Code Misalignment (rare):}} \changes{In functionally correct programs using logic patterns different from the description (e.g., DP instead of greedy) the annotations described by \projName sometimes failed to capture the real purpose of the code.}
% \end{itemize}

%   \changes{For Tester, the primary reason for misclassification are false negatives, which we categorize into the following classes:}

% \begin{itemize}[leftmargin=*]
%     \item \changes{\textbf{Misunderstood Problem Semantics (19\%):}} \changes{The LLM did not properly capture the relationship between inputs and expected outputs due to a misunderstanding of the problem.}

%     \item \changes{\textbf{Output Computation Failure (58\%):}} \changes{The LLM correctly described the logic of the problem requirements, but failed to execute this logic accurately when computing expected outputs.}

%     \item \changes{\textbf{Multiple Valid Solutions (21\%):}} \changes{For problems allowing multiple correct outputs for the same input, LLM-generated test asserted one output while the program produced another one. An example is presented in supplementary materials \Cref{tester_multiple_correct}.}

%     \item \changes{\textbf{Incomplete Test Sets (2\%):}} \changes{The LLM generated tests that did not sufficiently cover corner cases.}
% \end{itemize}

% \changes{To illustrate Tester's output computation failures, consider the example in \Cref{fig:tester_example}, which complements the example in \Cref{fig:motivating_example}, since it involves a correct solution. The program implements a typical greedy algorithm using a suffix array to compute the minimum cost for Kirill to reach the front of the queue according to the problem description. The Vanilla classifier fails to correctly interpret the algorithm, claiming that the use of \texttt{min(a[i], b[i])} did not properly account for bribes across all intermediate positions. The tester classifier also misclassifies it as incorrect as it incorrectly predicts the expected output. This phenomenon was also observed in previous research that showed that for non-trivial computations, it is easier for an LLM to generate a program to solve the problem than directly predict the intended output~\cite{gao2023pal}. In contrast, \projName's state annotations clarified the role of the list \texttt{c} as encoding the minimum cost to pass or swap per position, helping the LLM to arrive at the correct conclusion.}

\subsection{Qualitative Error Analysis}
\label{qualitative}

\begin{figure*}[t]
  \centering
  \includegraphics[width=\textwidth]{tester_example.png}
  \caption{\changes{A \benchName problem with a correct solution. \textsc{Vanilla} misclassifies by misreading the algorithm; \textsc{Tester} fails by predicting wrong outputs; \projName’s state annotations enable a correct judgment.}\label{fig:tester_example}}
  \vspace{-2mm}
\end{figure*}

\changes{We manually audited misclassifications to understand failure modes of \projName and Tester.}First for \projName:
\begin{itemize}[leftmargin=*,nosep]
  \item \changes{\textbf{Missing Edge Case:}} \changes{The NSP chain omits critical boundary conditions, leading to overly optimistic judgments.}

  \item \changes{\textbf{Oversimplified Annotations:}} \changes{Annotations blur fine-grained state differences, producing misleading correctness cues.}

  \item \changes{\textbf{Incorrect Loop Generalization:}} \changes{Loop effects are overgeneralized, resulting in faulty postconditions.}

  \item \changes{\textbf{Semantic Trust Bias:}} \changes{When code structure resembles the description, the LLM sometimes ignores contradictory annotations and falsely assumes correctness.}

  \item \changes{\textbf{Interprocedural Failures:}} \changes{Reasoning errors occur across function boundaries, corrupting inference consistency.}

  \item \changes{\textbf{Variable Scope or Shadowing Errors:}} \changes{In loop-heavy code, reused indices or scope bleed distort postconditions.}

  \item \changes{\textbf{Description–Code Misalignment (rare):}} \changes{For correct programs using alternative logic patterns (e.g., DP vs.\ greedy), annotations fail to reflect intent.}
\end{itemize}
and for Tester, false negatives fell in four categories:
\begin{itemize}[leftmargin=*,nosep]
  \item \changes{\textbf{Misread semantics (19\%):} misunderstood IO relation.}
  \item \changes{\textbf{Output computation failure (58\%):} logic is described correctly but numerical results are wrong.}
  \item \changes{\textbf{Multiple valid answers (21\%):} tests assert one of several correct outputs (\Cref{tester_multiple_correct}).}
  \item \changes{\textbf{Insufficient coverage (2\%):} corner cases not tested.}
\end{itemize}

\changes{In \Cref{fig:tester_example}, the program (a greedy method with a suffix array) is correct. \textsc{Vanilla} misreads the role of \texttt{min(a[i],b[i])}; Tester predicts an incorrect output—consistent with findings that LLMs struggle to directly compute nontrivial results~\cite{gao2023pal}. \projName’s annotations make the purpose of \texttt{c} explicit (per-position minimum cost), enabling the correct verdict.}

% Threats to Validity Section
% \section{Threats to Validity}
% \label{sec:threats_to_validity}

% While our experiments demonstrate the effectiveness of \projName, several factors may affect the validity and generalizability of our findings.

% LLMs are inherently nondeterministic, meaning that the same prompt can yield different responses across multiple runs. To mitigate this, we conducted multiple reruns per sample and used a confidence-based sampling approach to determine the necessary number of reruns. This ensures a high-confidence evaluation of classifier performance. While additional reruns might further refine our results, our methodology already provides a robust estimate of performance.
% \sloppy
% We evaluated \projName across models of varying sizes and specializations, ranging from 7b to 72b, as well as \texttt{Qwen2.5-Coder-32B-instruct}, which is optimized for code-related tasks. This selection covers a broad spectrum of model capabilities, from general-purpose LLMs to specialized coding models. While results might vary for other LLM families or architectures, the observed trends suggest that the reasoning improvements from \projName are widely applicable. Furthermore, even if \projName’s benefits were limited to this model family, the findings remain valuable given the increasing industry focus on proprietary small to mid sized proprietary models instead of relying solely on external state-of-the-art systems.
% \fussy


% To ensure the integrity of our evaluation, we avoided datasets known to have leaked into LLM training corpora (e.g., APPS, MBPP) and instead constructed a new dataset using Codeforces competition problems from 2024. These problems are human-written and highly challenging, ensuring a robust evaluation of LLM capabilities. In future work, we aim to extend our analysis to additional unseen datasets with varying levels of complexity. Additionally, we acknowledge that dataset selection can introduce biases in problem difficulty and error types. While we attempted to mitigate this by selecting a diverse set of submissions from different users with different degrees of correctness, further studies on alternative datasets could provide additional insights.

% \projName relies on structured prompting techniques to guide the LLM in reasoning about execution states. While prompt engineering can impact model performance, we mitigate this by applying the same structured prompt framework across all classifiers. This ensures that comparisons between classifiers remain fair, even if prompt phrasing influences absolute performance. 

% All experiments were conducted using the Alibaba Cloud API and Fireworks AI API, which provides a stable inference environment. However, LLM behavior can vary across different API providers and hardware configurations due to caching mechanisms, optimizations, or fine-tuning differences. While results may shift slightly under different conditions, we expect the overall performance trends of \projName relative to baseline classifiers to remain consistent. 

% A crucial consideration when evaluating an LLM-based approach is ensuring that the reported results are based on actual API responses rather than manually curated outputs. To ensure full transparency and reproducibility, we provide all API calls, including prompts and responses, in our experimental artifacts. This allows independent verification that our findings are derived from real model interactions rather than selective sampling or cherry-picking of favorable results.

% \section{Limitations \& Future Work}
% \label{sec:limitations}


% \subsection{High Token Usage}
% \projName, particularly with loop unrolling enabled, significantly increases token consumption due to the iterative nature of execution-state analysis. Specifically, the unrolling version requires \textbf{7.8$\times$} the tokens of the non-unrolled variant and \textbf{34.9$\times$} the tokens of a simple vanilla classifier. However, our results indicate that even the no-unroll version provides significant performance improvements while consuming \textbf{87\% fewer tokens} than the fully unrolled version. Notably, \projName (No Unroll) still outperforms the top-performing vanilla classifier (zero-shot-COT) by \textbf{28\%}, making it a cost-effective alternative for many practical use cases.  

% Step-by-step execution-state reasoning techniques, such as \projName, are inherently more computationally expensive than their direct classification counterparts. This trend aligns with well-known trade-offs in software verification, where symbolic execution is more costly than static analysis but provides deeper insights. Similarly, \projName's structured reasoning provides stronger correctness classification but at the expense of increased computational overhead.  

% Future work will explore optimizations to reduce unnecessary API calls. Specifically, we aim to develop an adaptive unrolling approach that selectively unrolls only high-level loops while summarizing lower-level control structures such as \texttt{if}, \texttt{try-except}, and basic operations in a single API call. This hybrid strategy could maintain much of the accuracy boost while significantly reducing token consumption.






% \subsection{Lack of Formal Guarantees and Symbolic Execution Limitations}

% \subsubsection{Absence of Formal Guarantees}
% A key limitation of \projName is that it does not provide formal guarantees of correctness. While structured execution-state reasoning improves classification accuracy, the model ultimately relies on learned heuristics rather than symbolic or formal verification methods. This means that even high-confidence predictions can be incorrect, as there is no guarantee of soundness or completeness in the classification.

% In contrast, formal verification techniques, such as symbolic execution and Hoare logic-based verification, provide provable correctness guarantees under specific conditions. However, these methods are typically infeasible at scale due to high computational cost, making LLM-based methods attractive as a scalable yet heuristic-driven alternative. Future work could explore hybrid approaches that integrate LLM reasoning with symbolic verification methods—for example, using \projName to propose candidate invariants that are formally checked using SMT solvers or incorporating lightweight static analysis to detect likely misclassifications.

% \subsubsection{Limitations in Symbolic Reasoning}
% While \projName applies Hoare logic-inspired execution-state tracking, it does not engage in true symbolic reasoning. The LLM reasons about concrete execution states based on natural language descriptions rather than symbolically manipulating variables as in traditional theorem proving or constraint solving. This limits its ability to generalize across all possible program executions, particularly when dealing with complex loop invariants or recursive functions.
% Integrating external constraint solvers with \projName could enable symbolic refinement, allowing the model to verify certain program properties formally rather than relying solely on natural language inference.

% \subsection{Implementation Limitations }

% \subsubsection{Handling of Variable Scope, Side Effects, and Data Structures}
% \projName does not explicitly track variable scope, side effects, or complex data structures, which can lead to misclassifications in cases involving variable shadowing, mutable state, or operations on nested data structures. Future improvements could enhance state tracking to better reason about variable mutations and data flow.

% \subsubsection{Handling of Multiple Functions and Function Calls}
% Currently, \projName analyzes functions in isolation without propagating execution states across calls, limiting its effectiveness for multi-function programs. A potential improvement is using function postconditions as preconditions for subsequent analyses, enabling more accurate reasoning about interdependent functions.
% \subsection{Generlisability}
% In future work we also want to explore a broader range of datasets with varying complexity to further assess the generalizability of \projName. Expanding our evaluation to larger and more diverse LLMs, including models beyond the Qwen family, will provide deeper insights into the scalability of structured reasoning techniques. Additionally, incorporating multiple API providers and hardware environments will help quantify the impact of provider-specific optimizations on model behavior, ensuring more robust and reproducible findings across different inference settings

% \section{Discussion}
% \label{sec:discussion}

% \projName introduces a new paradigm of structurally reasoning about programs in natural language. Although it does not guarantee correctness, it demonstrates a significant increase in the correctness classification accuracy w.r.t. natural language requirements by bridging the gap between formal semantics and informal requirements. Incorporating techniques like symbolic CoT~\cite{xu2024faithful} may help further enhance reasoning and provide formal guarantees.

% Previous research showed that formal proofs can be effectively constructed by following steps of informal proofs~\cite{jiang2022draft}. Thus, \projName can serve as the first step in the process of automating formal software verification. Similarly, state description computed by \projName may enhance testing by increasing test coverage and producing more precise oracles.

% Shi et al.~\cite{shi2024natural} proposed using natural language outlines to improve human-AI interaction. \projName can enhance these outlines by more accurately capturing program semantics.


% A limitation of \projName, especially when employing the few-shot-driven $k$-induction, is a high token consumption. Although a comparable balance between cost and performance is shared by many agentic systems~\cite{liu2024groupdebate}, and we show that the token use can be significantly reduced via potential fine-tuning, or by disabling unrolling, we believe more research is required to reduce the cost. Applying prompt compression~\cite{jiang2023longllmlingua} or grouping several statements during compositional reasoning are promising directions.

% Our implementation has limitations in handling program semantics, include handling variable scope, side effects, complex data structures, and multi-function analyses. Enhancements in state tracking for mutations and data flow, as well as using function postconditions as preconditions in interdependent function analyses, represent directions for future improvement.

% Natural language is an effective medium to reason about programs. First, it facilitates the incorporation of informal specifications such as comments and documentation, which are more ubiquitous in real-world settings than formal specifications. Second, the rich vocabulary of natural language permits succinct abstractions of program behaviour, e.g. the natural triple $$\langle \text{a is array} \rangle \,\mathtt{sort(a)}\, \langle \text{a is sorted in non-decreasing order} \rangle$$ relies on the English terms ``sorted'' and ``non-decreasing order''. In contrast, formal logic based reasoning often grapples with a trade-off between expressivity and computational tractability, constrained by a more limited vocabulary. To express an equivalent triple, existing techniques resort to the theory of arrays, which not only leads to a less concise depiction but also renders the reasoning problem undecidable~\cite{bradley2006s}: $$\{ a: \mathrm{array} \} \,\mathtt{sort(a)}\, \{ \forall i.\ \forall j.\ i>0 \wedge j<\mathit{len} \wedge i\leq j \rightarrow \mathrm{select}(a,i)\leq \mathrm{select}(a,j) \}.$$ Third, reasoning in natural language leverages the natural language channel of code~\cite{casalnuovo2020theory}, such as variable names, to enhance analysis precision. One notable limitation of this approach is its inability to offer formal correctness guarantees, as it acts as an fuzzy correctness classifier rather than a prover. Nonetheless, we demonstrate its practical utility in tasks like bug finding and code generation.

\section{Threats of Validity}
\label{sec:threats}

LLMs are inherently nondeterministic~\cite{ouyang2025empirical}, leading to varied responses. To address this, we employed multiple reruns per sample with confidence-based sampling to ensure robust performance estimates. Although further reruns could slightly refine our results, the current methodology provides high-confidence evaluations.

Main threats to external validity are that the selected dataset may not generalize to other programs, or that performance enhancements will not generalize to other LLMs. To avoid data leakage, we constructed a new dataset from recent Codeforces competition problems, ensuring challenging, human-authored, and varied submissions. Despite efforts to minimize biases, future studies on alternative datasets could further validate generalizability and performance trends. Apart from that, real-world requirements contain ambiguities, which are typically not present in programming competition problems. Techniques like ClarifyGPT~\cite{mu2023clarifygpt} that detects ambiguities in informal requirements may address this concern.

\changes{We evaluated \projName across models of diverse sizes (7B to 72B) from 2 families (Qwen, Llama), including Qwen-32B-Coder, a model specialized for coding tasks.} This broad spectrum enhances generalizability; however, different LLM architectures may yield varying results. Nevertheless, observed trends suggest broad applicability, especially relevant given industry shifts toward proprietary small and mid-sized models~\cite{chen2024role}. Evaluating \projName on larger and more powerful models may help assess its scalability and effectiveness in leveraging more advanced reasoning capabilities.

\vspace{-2mm}

\section{Related Work}
\label{sec:related_work}
% This section should provide an overview of existing approaches to program verification, particularly those that use natural language or LLM-based reasoning. Compare and contrast \projName with other methods used for verifying correctness in code generation tasks. Also, mention traditional Hoare logic tools and how they have evolved with AI.

% the natural language entailment operator~\cite{androutsopoulos2010survey}

% Recent work applied Large Language Models (LLMs) to assess the correctness of code w.r.t. natural language specification~\cite{chen2023teaching,olausson2023self}, but their accuracy remains low on this task.

% MBPP~\cite{austin2021program} and APPS~\cite{hendrycks2021measuring} benchmarks

% discuss work on code summarization

% LLMs have demonstrated strong capabilities in code generation and reasoning, yet their ability to systematically verify correctness remains an open challenge. Several works have sought to enhance LLM-based program reasoning through execution traces, reranking techniques, and automated repair, but they fundamentally differ from HoarePrompt in methodology and objectives.

% NExT~\cite{ni2024next} improves LLM-based program repair by training on execution traces, enhancing reasoning about concrete program behaviors. However, it relies on observed executions rather than abstract execution-state reasoning, making it dependent on specific test cases. In contrast, HoarePrompt generalizes correctness verification by modeling all reachable execution states via Hoare Logic-inspired strongest postconditions.

% CodeMind~\cite{liu2024codemind} and REval~\cite{reval2024} serve as evaluation frameworks, assessing the reasoning capabilities of LLMs in tracking execution and predicting runtime behavior. While CodeMind identifies limitations in execution tracking and logical reasoning, and REval evaluates consistency in intermediate execution-state predictions, neither introduces a structured correctness verification method. HoarePrompt moves beyond evaluation by enforcing structured execution-state tracking to systematically verify correctness.

% Several works have explored improving LLM-based program synthesis and repair. CodeRSA~\cite{pragmatic} optimizes code generation through pragmatic reasoning, reranking multiple candidate solutions probabilistically. D4C~\cite{alligning} refines LLM-generated patches for automated program repair by aligning output formats with pretraining objectives.NIER~\cite{evaluating} evaluates the generalizability of LLM-based repair techniques under syntactic transformations. While these approaches enhance code generation and repair, they do not ensure correctness through structured execution-state reasoning. HoarePrompt, in contrast, provides a verification methodology independent of dataset-driven validation, making it robust to program transformations and variations.

% As LLMs are increasingly used in software development, ensuring correctness in generated code is critical. One approach involves detecting inconsistencies through defect detection methods, such as metamorphic prompt testing combined with cross-validation~\cite{wang2024validating}.

%% Other research has focused on improving LLM-driven program repair. CodeRSA~\cite{pragmatic} optimizes code synthesis via pragmatic reasoning, reranking multiple generated solutions probabilistically. D4C~\cite{alligning} enhances LLM-based patch generation by structuring output formats for better test-based validation, while NIER~\cite{evaluating} examines the generalizability of LLM-generated fixes across program transformations. While these works refine code generation and repair, they do not ensure correctness through structured execution-state reasoning.

\projName introduces a new paradigm for reasoning about programs in NL. Its contributions are relevant to works on reasoning about code with LLMs, LLM-assisted testing and verification.

%LLMs have demonstrated strong capabilities in code generation and reasoning, yet their ability to systematically verify correctness remains an open challenge. Several works have sought to enhance LLM-based program reasoning through execution traces, reranking techniques, and automated repair, but they fundamentally differ from HoarePrompt in methodology and objectives.


\textit{Reasoning About Code with LLMs.} Many generic approaches have been introduced to realize effective chain-of-though (CoT) reasoning~\cite{wei2022chain}. We employed Zero-shot-CoT~\cite{kojima2022large}, a popular automation of CoT, as a baseline in our experiments. CJ-Eval~\cite{zhao2024codejudge} uses LLMs as judges of program correctness applying an approach equivalent to our baseline Vanilla prompt (\Cref{sec:classification}). Since CJ-Eval is based on a widely-used Apps dataset~\cite{apps}, which is subject to data leakage, we constructed a new dataset \benchName. A large body of work focused on reasoning about concrete execution states. Zhang et al.~\cite{grounding} fine-tuned LLMs with concrete execution states to improve their ability to track execution faithfully. Similarly, NExT~\cite{ni2024next} enhances debugging using programs annotated with concrete execution traces. CodeMind~\cite{liu2024codemind} includes a ``specification reasoning'' task that measures how including test data in the specification helps the model generate correct code. REval~\cite{runtime} evaluates runtime behaviour reasoning capabilities, specifically, consistency in intermediate state prediction by analyzing variable states at specific execution points. The key difference of \projName is that it reasons about not individual executions but, potentially, infinitely many executions, without relying of specific inputs.

% Several works have explored LLMs for generating formal specifications and correctness proofs that can be verified by automated tools. Marmaragan~\cite{spark} employs LLMs to synthesize Ada/SPARK contracts, automatically inserting formal preconditions, postconditions, and invariants into programs, successfully generating verifiable specifications in approximately half of the cases. Similarly, SynVer~\cite{symver} biases LLM-generated C code towards verifiability by generating assertions and invariants, which a backend verifier, such as Coq or the Verified Software Toolchain (VST), attempts to formally prove correct.

% Beyond specification inference, recent work has integrated LLMs into the verified compilation pipeline. LLMLIFT~\cite{llmlift} translates source code into a high-level intermediate representation while simultaneously generating proof sketches, including loop invariants and semantic equivalences, which are then formally verified using SMT solvers. Unlike traditional verification workflows, LLMLIFT enables LLMs to participate in the proof process by guiding automated solvers.

% Extending this paradigm, Stanford et al.~\cite{pc3} propose \textit{Proof-Carrying Code Completions (PC³)}, envisioning AI-generated code that is always accompanied by a machine-checkable proof of correctness. Their prototype leverages the Dafny verifier, prompting an LLM to generate Dafny \cite{dafny} code with pre/postconditions, allowing end-users to formally verify AI-generated snippets without manual intervention.

% These approaches advance the integration of LLMs into formal methods by embedding verification properties directly into code generation, enabling correctness proofs via tools such as Dafny, Coq\cite{Coq}, and Z3\cite{z3}. While HoarePrompt shares the goal of leveraging LLMs for correctness reasoning, it differs in that it does not rely on external solvers. Instead, HoarePrompt enforces correctness verification through structured execution-state tracking via Natural Strongest Postconditions (NSP) and k-induction, allowing reasoning within the LLM itself rather than outsourcing verification to formal tools.

\textit{LLM-Assisted Formal Verification.} Recent works on LLM-assisted verification focused on translating natural language specifications into formal ones, and rely on external theorem provers or SMT solvers~\cite{endres2024can}. Marmaragan~\cite{spark} synthesizes formal Ada/SPARK contracts from natural language. Similarly, SynVer~\cite{symver} biases LLM-generated C code towards verifiability by producing assertions and invariants, subsequently verified by Coq or Verified Software Toolchain (VST). LLMLIFT~\cite{llmlift} integrates LLMs into verified compilation, generating proof sketches alongside code translation, with formal verification performed via SMT solvers. Stanford et al.~\cite{pc3} introduce Proof-Carrying Code Completions (PC3), prompting LLMs to produce Dafny~\cite{dafny} programs accompanied by machine-checkable correctness proofs, automating verification through Dafny. While these approaches embed formal verification into LLM-driven code generation pipelines, they fundamentally rely on external verification systems such as Dafny, Coq~\cite{Coq}, or Z3~\cite{z3}. Despite the promise of these methods, they have shown limited scalability, with Marmaragan achieving only 50.7\% correctness in its generated Ada/SPARK contracts \cite{spark} and Endres et al.~\cite{endres2024can} reporting that even with optimized prompting, only 77\% of LLM-generated postconditions were verifiable. \projName employs an alternative approach of constructing structured informal proofs in natural language, leveraging its expressive power and flexibility, for classifying program correctness. Following previous work~\cite{jiang2022draft} such informal proofs might be translated into formal ones to automate formal verification. Closest to our loop reasoning technique, Liu et al.~\cite{liu2024generalloopinvariantgeneration} propose LLM-SE, which combines symbolic execution with LLMs to generate loop invariants for memory-manipulating programs. While \projName does not target formal invariant synthesis, instead focusing in textual inductive reasoning. \changes{Although \projName does not provide formal guarantees, prior work shows that formal proofs mirror informal reasoning~\cite{jiang2022draft}, suggesting \projName could serve as a foundation for automated verification. Integrating \projName with techniques like symbolic CoT~\cite{xu2024faithful} may further enhance reasoning and provide guarantees. A promising future direction is to develop a hybrid (formal and natural) language to promote more grounding of the LLM's output}.


\textit{LLM-Based Automated Test Generation} Recent works~\cite{li2024large,schafer2023empirical,yuan2023no,chen2024chatunitest} have employed LLMs for automated test generation. Apart from that, the analysis of LLM-generated test outcomes has been used to facilitate program synthesis tasks~\cite{xiong2023program}, exemplified by CODET~\cite{chen2022codet}, AgentCoder~\cite{huang2023agentcoder}, CodeCoT~\cite{huang2023codecot}, and CoCoST~\cite{he2024cocost}. \Cref{sec:evaluation} shows that using LLM-generated tests for correctness classification, embodied in our \textit{Tester} baseline adopting AgentCoder's approach~\cite{huang2023agentcoder}, suffers from high false negative rate, as LLMs fail to infer precise assertions for complex requirements. Ideas of \projName could be synergized with testing to improve accuracy and cost. \newline
\hspace*{0.5em}\textit{Competition-Level Benchmarks} Recent studies reveal significant data leakage in widely used LLM benchmarks, inflating model performance due to test-training overlaps~\cite{training_on_bench,deng-etal-2024-unveiling}. New benchmarks, such as LiveCodeBench~\cite{jain2024livecodebenchholisticcontaminationfree}, CodeElo~\cite{quan2025codeelobenchmarkingcompetitionlevelcode}, ProBench~\cite{yang2025probenchbenchmarkinglargelanguage}, and USACO-based evaluations~\cite{shi2024language}, mitigate leakage by incorporating recent programming competition problems. However, they primarily assess code generation capabilities and do not provide solution submissions with correctness labels. In contrast, \benchName specifically targets correctness classification by offering labeled problem-solution pairs sourced from Codeforces 2024 contests, enabling evaluation of LLMs on their execution-state reasoning capabilities rather than memorization or test-based validation alone.

\vspace{-3mm}

\section{Conclusion}

This work introduces \projName, the first LLM-based program analysis approach to assess program correctness against natural language requirements. Unlike traditional verification and testing that depend on external tools or concrete test executions, \projName systematically applies natural language-based structural reasoning inspired by the strongest postcondition calculus and k-induction. Our evaluation shows significant performance improvements in classification quality, with \projName increasing the MCC by 61\% over Zero-shot-CoT, 72\% over Vanilla, and 106\% over an LLM testing-based classifier. Our novel k-induction further enhances performance, yielding an average MCC improvement of 26\% compared to a non-inductive approach. Although \projName consumes more tokens, the no-unroll variant provides a practical trade-off, achieving a 28\% MCC improvement over Zero-shot-CoT while using 90\% fewer tokens than the standard \projName version.

%Future research will focus on exploring adaptive loop-unrolling strategies to mitigate token consumption, like selectively summarizing low-level control structures. Evaluating \projName across larger, more advanced LLM architectures beyond the Qwen family and more diverse datasets will further validate scalability and generalizability.
% \section{Conclusions and Future Work}
% \label{sec:conclusion}
% Summarize the key takeaways from the study, particularly regarding the performance of \projName on handling functions with multiple return points. Highlight any significant contributions to the field. In future work, suggest potential improvements to the system, such as better handling of edge cases or expanding to new datasets.
\begin{acks}
This work was sponsored by the National Key Research and Development Program of China under Grant No.\ 2022YFB4501902 and the National Natural Science Foundation of China (NSFC) under Grant No.\ W2411051.

\end{acks}

\bibliographystyle{ACM-Reference-Format}
\bibliography{references}

\onecolumn
\pagebreak
\twocolumn
\appendix

\section{Supplementary Materials}

All code, documentation, results and complimentary material for this work is available in our repository: \url{https://github.com/msv-lab/HoarePrompt}.\looseness=-1

\subsection{Details of Motivating Example}
\label{appendix:motivating_example}

For evaluating the motivating example presented in \Cref{sec:motivating_states}, we ran a series of experiments using three different \texttt{qwen2.5-instruct} models (\texttt{7b}, \texttt{coder-32b}, \texttt{72b}) \changes{and one Llama model \\(\texttt{Llama3.1-70b-instruct})} at three temperatures \((0.5, 1.0, 1.5)\).\\ For each temperature, we tested:
\begin{itemize}
    \item \textbf{Not Annotated}: 6 different prompts with no reachable state annotations.
    \item \textbf{Annotated}: 6 different prompts that included reachable state descriptions.
\end{itemize}
Each model/prompt combination was run 10 times, for a total of 1080 calls across all setups. Table~\ref{appendix:tab:motivating_example_results} shows the success rates (the percentage of calls that correctly identified the code as buggy).


\begin{table}[ht]
\centering
\caption{Success rates (\%) for detecting the indentation bug under different prompts (``Not Annotated'' vs.~``Annotated''), temperatures, and models. Each cell shows \emph{successful runs} / \emph{total runs} (percentage).}
\label{appendix:tab:motivating_example_results}
\begin{tabular}{l c c}
\toprule
\multicolumn{3}{c}{\textbf{Temperature = 0.5}} \\
\midrule
\textbf{Model} & \textbf{Not Annotated} & \textbf{Annotated} \\
\midrule
\texttt{7b-instruct}          & 0/60 (0\%)   & 7/60 (11.7\%) \\
\texttt{coder-32b-instruct}   & 9/60 (15\%)  & 51/60 (85\%)  \\
\texttt{72b-instruct}         & 21/60 (35\%) & 51/60 (85\%)  \\
% \texttt{llama3.1-70b-instruct}  & 30/60 (50\%) & 60/60 (100\%)  \\
\changes{\texttt{llama3.1-70b-instruct}}  & \changes{30/60 (50\%)} & \changes{60/60 (100\%)}  \\
\midrule
\multicolumn{3}{c}{\textbf{Temperature = 1.0}} \\
\midrule
\textbf{Model} & \textbf{Not Annotated} & \textbf{Annotated} \\
\midrule
\texttt{7b-instruct}          & 0/60 (0\%)   & 2/60 (3.3\%)  \\
\texttt{coder-32b-instruct}   & 9/60 (15\%)  & 50/60 (83.3\%)\\
\texttt{72b-instruct}         & 29/60 (48.3\%) & 49/60 (81.7\%)\\
% \texttt{llama3.1-70b-instruct}  & 30/60 (50\%) & 60/60 (100\%)  \\
\changes{\texttt{llama3.1-70b-instruct}}  & \changes{30/60 (50\%)} & \changes{60/60 (100\%)}  \\

\midrule
\multicolumn{3}{c}{\textbf{Temperature = 1.5}} \\
\midrule
\textbf{Model} & \textbf{Not Annotated} & \textbf{Annotated} \\
\midrule
\texttt{7b-instruct}          & 4/60 (6.7\%)  & 5/60 (8.3\%)  \\
\texttt{coder-32b-instruct}   & 7/60 (11.7\%) & 51/60 (85\%)  \\
\texttt{72b-instruct}         & 31/60 (51.7\%) & 42/60 (70\%)  \\
% \texttt{llama3.1-70b-instruct}  & 40/60 (67\%) & 60/60 (100\%)  \\
\changes{\texttt{llama3.1-70b-instruct}}  & \changes{40/60 (67\%)} & \changes{60/60 (100\%)}  \\
\bottomrule
\end{tabular}
\end{table}


%% \subsubsection{Comparison with a Dedicated Reasoning Model}
%% One could argue that instead of explicitly annotating execution states, a reasoning model should be able to detect the bug independently. To investigate this, we tested the QwQ-32B-Preview model, which specializes in logical reasoning, using the same vanilla prompting approach.

%% A key challenge observed in this setup was that when the model detected a potential issue, it often attempted to generate test cases to validate its hypothesis. However, due to the complexity of reasoning about multiple execution paths, the model frequently exceeded the output token limit before reaching a final conclusion, leaving the bug undetected within a single inference pass. To be fair, we assume that if the model were allowed additional calls to continue the conversation, it would eventually detect the bug. The results are shown in Table~\ref{tab:reasoning-results}.
%%  \begin{table}[H]
%%     \centering
%%     \caption{Success rates (\%) for detecting the indentation bug using Qwen2.5-32B-Preview (Reasoning Model) at different temperatures.}
%%     \label{tab:reasoning-results}
%%     \begin{tabular}{|c|c|}
%%         \hline
%%         \textbf{Temperature} & \textbf{Success Rate} \\
%%         \hline
%%         0.5  & 15/60 (25\%) \\
%%         \hline
%%         1.0  & 16/60 (26.7\%) \\
%%         \hline
%%         1.5  & 8/60 (13.3\%) \\
%%         \hline
%%     \end{tabular}
%% \end{table}

%% Even with a model explicitly trained for reasoning, the success rates remain relatively low. More importantly, the QwQ 32B-Preview model required significantly longer responses and, in many cases, failed to provide an answer at all due to the reasoning overhead.

%% By contrast, our \projName approach, even when applied to a relatively small model like Qwen2.5-7B, significantly outperforms the reasoning model by leveraging structured execution-state annotations. This demonstrates the potential of \projName to enhance LLM reasoning by guiding the model through critical intermediate states, rather than relying on a model’s internal, unconstrained reasoning mechanisms.

% \subsection{Motivational Example for Tester Classifier from MBPP Dataset}
% \label{sec:appendix_tester_example_Mbpp}
% Despite recent advances, large language models still struggle to reliably generate high-quality test cases for software. In a systematic evaluation of algorithmic problems, \textit{TestCase-Eval} reveals that even the top-performing model, Qwen3-32B, achieves only 43.8\% on the Fault Exposure task—far below the 93.3\% performance of human experts~\cite{yang2024testcaseeval}. This gap highlights LLMs' difficulty in identifying inputs that expose subtle or complex faults. Moreover, \textit{VALTEST} shows that even the strongest model evaluated, GPT-4o, produces invalid test cases in at least 26\% of attempts~\cite{taherkhani2024valtest}, necessitating automated validation mechanisms. Similar limitations appear in unit testing domains: \textit{Evaluating and Improving ChatGPT for Unit Test Generation} reports that only 24.8\% of tests generated by ChatGPT execute successfully, with 57.9\% failing to compile and another 17.3\% crashing at runtime due to incorrect assumptions~\cite{liu2024chatgptunit}. Additionally, a comparative study~\cite{guzu2025chatgptsbst} shows that ChatGPT often generates invalid or ineffective assertions that fail to uncover the intended bugs, further undermining the trustworthiness of its test outputs.

% This section presents a concrete failure case from the MBPP (Mostly Basic Python Programming) dataset \cite{mbpp}, a widely-used benchmark designed to evaluate basic Python programming skills. MBPP is considered relatively easy for current models: for instance, GPT-4o achieves 87.5\% accuracy on this benchmark \cite{Wang_2024} for code generation, and newer models are reported to achieve near-perfect accuracy on code generation tasks\footnote{\url{https://paperswithcode.com/sota/code-generation-on-mbpp}}.

% Despite its simplicity, this dataset still exposes fundamental limitations of the LLM based testing, using the exact prompt from \cite{huang2023agentcoder}. Specifically, we demonstrate how a correct implementation is misclassified as incorrect due to an inaccurate expected output in an LLM-generated test case. This illustrates that even when test generation succeeds syntactically, it may fail semantically, especially for problems involving nontrivial computation or external knowledge (e.g., prime factorization).



% \begin{figure}[h]

%     \centering
%  \begin{minipage}[t]{0.48\textwidth}
% \begin{tcolorbox}[colframe=grey!100, colback=gray!3, title=Problem Description, fonttitle=\centering, before upper=\raggedright, boxrule=0.8mm, sharp corners=southwest, left=2pt, right=2pt, top=2pt, bottom=2pt, boxsep=1pt]
% \footnotesize
% Write a Python function that returns the largest prime factor of a given number. 
% \end{tcolorbox}
% \begin{tcolorbox}[colframe=darkgreen!80, colback=gray!3, title=Tester prompt with problem description, fonttitle=\centering, before upper=\raggedright, boxrule=0.8mm, sharp corners=southwest, left=2pt, right=2pt, top=2pt, bottom=2pt, boxsep=1pt]
% {\fontsize{7pt}{11pt}\selectfont
% \textbf{Role}: As a tester, your task is to create comprehensive test cases for the incomplete \texttt{largest\_prime} function. These test cases should encompass \textbf{Basic}, \textbf{Edge}, and \textbf{Large Scale} scenarios to ensure the code's robustness, reliability, and scalability.

% \textbf{Input Code Snippet}:
% \begin{lstlisting}[language=Python]
% def largest_prime(number: int) -> int:
%     """
%     Write a Python function that returns the largest prime factor of a given number.
%     >>> largest_prime(4)
%     2
%     >>> largest_prime(9)
%     3
%     """
% \end{lstlisting}

% \textbf{1. Basic Test Cases}:\\
% - \textbf{Objective}: To verify the fundamental functionality of the \texttt{largest\_prime} function under normal conditions.

% \textbf{2. Edge Test Cases}:\\
% - \textbf{Objective}: To evaluate the function's behavior under extreme or unusual conditions.

% \textbf{3. Large Scale Test Cases}:\\
% - \textbf{Objective}: To assess the function’s performance and scalability with large data samples.

% \textbf{Instructions}:
% \begin{itemize}
% \item Implement a comprehensive set of test cases following the guidelines above.
% \item Ensure each test case is well-documented with comments explaining the scenario it covers.
% \item Pay special attention to edge cases as they often reveal hidden bugs.
% \item For large-scale tests, focus on the function's efficiency and performance under heavy loads.
% \end{itemize}
% }
% \end{tcolorbox}


% \begin{tcolorbox}[colframe=nicered!90, colback=gray!3, title=Tester Response, fonttitle=\centering, before upper=\raggedright, boxrule=0.8mm, sharp corners=southwest, left=2pt, right=2pt, top=2pt, bottom=2pt, boxsep=1pt]
% \footnotesize
% <Other test cases  omitted for brevity>\\

% \textbf{Test 10} \\
%  assert largest\_prime(1234567890123456789) == 953051  \#Largest prime factor of 1234567890123456789 should be 953051
%  \#\highlight{\textcolor{black}{wrong, actual output is 27961}} \\
% \textbf{Test 11} \\
%  assert largest\_prime(999999999) == 3 \#Largest prime factor of 999999999 should be 3
%  \#\highlight{\textcolor{black}{wrong, actual output is 333667}} \\

% \end{tcolorbox}
% \end{minipage}
%  \caption{Tester Prompting Example for MBPP problem }
%     \label{fig:prompt_test_comparison}
% \end{figure}
% \textbf{Explanation:} This test is intended to evaluate the function’s performance with a large number. The tester incorrectly assumes 3 is the largest prime factor of 999999999, and that 953051 is the largest prime factor of 1234567890123456789 but in fact, it is 333667 and 27961 respectively. As a result, even a correct implementation that outputs the correct largest prime factors will be misclassified as incorrect by the Tester classifier due to the flawed test oracle.
% \subsection{Motivational Example for Tester Classifier}
% \label{sec:appendix_tester_example}

% This section provides a detailed example highlighting the limitations of LLM-generated test cases in program correctness assessment.  It demonstrates that LLM-based testing struggles with generating effective test cases and computing expected outputs accurately for complex requirements.

% \begin{figure*}[h]
%     \centering
%     \begin{minipage}[t]{0.57\textwidth}
%           \vspace{0pt}
%         % Vanilla Prompting - Prompt Box
%          % Vanilla Prompting - Response Box
%         \begin{tcolorbox}[colframe=darkgreen!100, colback=gray!3, title=Problem Description, fonttitle=\centering, before upper=\raggedright, boxrule=0.8mm, sharp corners=southwest, left=2pt, right=2pt, top=2pt, bottom=2pt, boxsep=1pt]
%         \footnotesize
%            Equalize the total number of people (students, administrators, and professors) between Building A (12 classes) and Building B (16 classes) by removing excess individuals from the more populated building. The input consists of five lists: students\_A (12 values), admins\_A (12 values), students\_B (16 values), admins\_B (16 values), and profs\_B (16 values), representing the number of people in each class. Each class has different number of people than the others. Each remaining person in both buildings contributes \$10 to a bonus pool, which is evenly distributed among those removed. Return the per-person bonus for the moved individuals.
%         \end{tcolorbox}
%         \begin{tcolorbox}[colframe=nicered!90, colback=gray!3, title=Tester prompt with problem description, fonttitle=\centering, before upper=\raggedright, boxrule=0.8mm, sharp corners=southwest, left=2pt, right=2pt, top=2pt, bottom=2pt, boxsep=1pt]
% \fontsize{7pt}{11pt}\selectfont
% \textbf{Role}: As a tester, your task is to create comprehensive test cases for the following coding problem. These test cases should encompass \textbf{Basic} and \textbf{Edge} scenarios to ensure the code's robustness, reliability, and scalability.


% \textbf{Problem Description}: Fill with Problem Description\\
% \textbf{1. Basic Test Cases}: To verify the fundamental functionality of the function under normal conditions.\\
% \textbf{2. Edge Test Cases}: To evaluate the function's behavior under extreme or unusual conditions.\\

% \textbf{Instructions}:
% \begin{itemize}
% \item Implement a comprehensive set of test cases following the guidelines above.
% \item Ensure each test case is complete (no omission) and well-documented with comments explaining the scenario it covers.
% \item Pay special attention to edge cases as they often reveal hidden bugs.
% \item Do not repeat, do not summarize.
% \end{itemize}

% \end{tcolorbox}
            
%         % Vanilla Prompting - Response Box
%         \begin{tcolorbox}[colframe=nicered!90, colback=gray!3, title=Tester Response, fonttitle=\centering, before upper=\raggedright, boxrule=0.8mm, sharp corners=southwest, left=2pt, right=2pt, top=2pt, bottom=2pt, boxsep=1pt]
%         \footnotesize
%             <The rest of the test cases are ommited for brevity>
          
% Input = 
% [10,10,10,10,10,10,10,10,10,10,10,10], [2,2,2,2,2,2,2,2,2,2,2,2], [10,10,10,10,10,10,10,10,10,10,10,10,10,10,10,10], [2,2,2,2,2,2,2,2,2,2,2,2,2,2,2,2], [3,3,3,3,3,3,3,3,3,3,3,3,3,3,3,3] \\
% Output =
% zero division exception \# \highlight{\textcolor{black}{wrong the 2 building don't have same size}}
% \\
% Input =[20,20,20,20,20,20,20,20,20,20,20,20], [5,5,5,5,5,5,5,5,5,5,5,5], [10,10,10,10,10,10,10,10,10,10,10,10,10,10,10,10], [2,2,2,2,2,2,2,2,2,2,2,2,2,2,2,2], [3,3,3,3,3,3,3,3,3,3,3,3,3,3,3,3]\\
% Output: 
% 10 \# \highlight{\textcolor{black}{wrong, actual is 80}}


%         \end{tcolorbox}
%     \end{minipage}
%     \hfill
%     \begin{minipage}[t]{0.42\textwidth}
%         % Prompted with Execution States - Prompt Box
%           \vspace{0pt}
%         \begin{tcolorbox}[colframe=niceblue, colback=gray!3, title=Prompt with Reachable States as comments in the code, fonttitle=\centering, before upper=\raggedright, boxrule=0.8mm, sharp corners=southwest, left=2pt, right=2pt, top=2pt, bottom=2pt, boxsep=1pt]
%         {\fontsize{9pt}{11pt}\selectfont \textbf{Is the following program correct based on the program description?}}
%             \begin{minted}[escapeinside=||]{python}
%     def func_1(students_A, admins_A, students_B, admins_B, profs_B):
%     total_A = sum(students_A) + sum(admins_A)
%     total_B = sum(students_B) + sum(admins_B) +sum(profs_B)
%     surplus = abs(total_A - total_B)
%     stayers = total_A + total_B - surplus
%     bonus_pool = stayers * 10
%     bonus_per_moved = bonus_pool / surplus
%     return bonus_per_moved
%     #The program returns `bonus\_per\_moved` which is 
%     #`bonus\_pool / surplus`. Given that 'total\_A' is 
%     #a positive integer equal to the number of people
%     # originally at building A, and 'total\_B' is a 
%     # positive integer equal to the number
%     # of people originally at building B  
%     #`bonus\_pool is a positive integer equal 
%     #to 10*(total\_A+  total\_B - surplus) 
%     #|\highlightblue{\textcolor{commentcolor}{and `surplus` is a positive integer or 0}}|
%     # equal to the  absolute difference between `total\_A`
%     #and `total_B`, the program returns 
%     #`(10*(total\_A+  total\_B - surplus)) / abs(total\_A - total\_B)`.
    
%             \end{minted}
%         \end{tcolorbox}

%         % Prompted with Execution States - Response Box
%         \begin{tcolorbox}[colframe=niceblue, colback=gray!3, title= Response for Prompt with Reachable States, fonttitle=\centering, before upper=\raggedright, boxrule=0.8mm, sharp corners=southwest, left=2pt, right=2pt, top=2pt, bottom=2pt, boxsep=1pt]
%         \footnotesize
%             The provided solution is \textbf{\textcolor{blue}{incorrect}} as it could lead to a potential division by zero exception since \highlightblue{\textcolor{black}{the variable \texttt{surplus}} could be 0 } in the case where the number of people in the 2 buildings is equal.
%         \end{tcolorbox}
%     \end{minipage}
%     \caption{Comparison of Tester Prompting and Prompting with Execution States. }
%     \label{fig:prompt_test_comparison}
% \end{figure*}

% In this example, the program contains a division-by-zero error that only occurs when a specific computed condition holds. However, the LLM-generated test cases fail to provide such inputs, preventing the bug from manifesting. Furthermore, even when valid inputs are provided, the LLM incorrectly predicts the expected output, leading to incorrect classification.

% \begin{figure}[t]
%     \centering
%     \includegraphics[width=0.8\linewidth]{example2-sec5.png}
%     \caption{Example 2 - Unique Email Domains \label{fig:unique_email}}
%     \vspace{-2mm}
% \end{figure}
% \begin{figure}[t]
%     \centering
%     \includegraphics[width=0.8\linewidth]{example3-sec5.png}
%     \caption{Example 3 - First name extraction \label{fig:first_name}}
%     \vspace{-2mm}
% \end{figure}

% \changes{\subsection{Motivational Example from \benchName}}
% \label{example_appendix}
% \changes{In Figure~\ref{motivating_from_dataset}, we show a representative example from our dataset where an LLM misclassified a correct Python program as incorrect when no execution annotations were provided. The program implements an efficient greedy strategy using a suffix array to compute the minimum cost for Kirill to reach the front of the queue. However, the LLM initially judged the logic to be flawed, claiming that the use of \texttt{min(a[i], b[i])} did not properly account for bribes across all intermediate positions.}

% \changes{The tester classifier also misclassifies the implementation as incorrect as the expected output for the given test case, does not account for some of the problem description requirements.}

% \changes{In contrast, when lightweight execution state annotations were added to the code—specifically clarifying the role of the list \texttt{c} as encoding the minimum cost to pass or swap per position—the LLM revised its judgment and correctly recognized the solution as valid. These annotations served to bridge the semantic gap between the code's compact logic and the problem's complex narrative, revealing how even simple metadata can unlock correct semantic interpretation for large language models.}

% \begin{figure}[t]
% \caption{\changes{Motivating example for \benchName dataset}}
% \label{motivating_from_dataset}
% \centering
% \begin{minipage}[t]{0.48\textwidth}
% \begin{tcolorbox}[colframe=grey!100, colback=gray!3, title=Problem Description, fonttitle=\centering, before upper=\raggedright, boxrule=0.8mm, sharp corners=southwest, left=2pt, right=2pt, top=2pt, bottom=2pt, boxsep=1pt]
% \footnotesize
% There are $n$ people standing in a queue, labeled from $1$ to $n$, waiting to ask Serafim the Owl about the meaning of life. Kirill arrives late and joins the end of the queue at position $n+1$. He is unhappy with his position and wants to move ahead in the queue.

% To do so, Kirill may swap positions with any person who stands before him. Specifically, if Kirill is at position $i$ and he chooses to swap with person $j$ such that $j < i$, he must:
% \begin{itemize}
%     \item Pay $a_j$ coins to the person at position $j$ for the swap.
%     \item Additionally, for every person $k$ between $j$ and $i$ (i.e., $j < k < i$), he must pay $b_k$ coins as a bribe to let him pass.
% \end{itemize}
% Kirill can perform this type of move repeatedly, choosing different positions to jump to. He wants to move up the queue to end up among the first $m$ positions (i.e., final position $\leq m$), and he wants to minimize the total number of coins spent.
% Kirill’s goal is to reach a position between $1$ and $m$ (inclusive) in the queue while spending as few coins as possible. You must compute the minimum total cost he needs to achieve this.
% \end{tcolorbox}


% \begin{tcolorbox}[colframe=black!85!blue, colback=gray!3, title=Original Program, fonttitle=\centering, before upper=\raggedright, boxrule=0.8mm, sharp corners=southwest, left=2pt, right=2pt, top=2pt, bottom=2pt, boxsep=1pt]
% \footnotesize
% \begin{lstlisting}[language=Python]
% def solve(n, k, a, b):
%     c = [min(a[i], b[i]) for i in range(n)]
%     suff = [0] * (n + 1)
%     for i in range(n - 1, -1, -1):
%         suff[i] = suff[i + 1] + c[i]
%     ans = float('inf')
%     for i in range(k):
%         ans = min(ans, a[i] + suff[i + 1])
%     print(ans)

% t = int(input())
% for _ in range(t):
%     n, k = map(int, input().split())
%     a = list(map(int, input().split()))
%     b = list(map(int, input().split()))
%     solve(n, k, a, b)
% \end{lstlisting}
% \end{tcolorbox}


% \begin{tcolorbox}[colframe=nicered!90, colback=gray!3, title=Zero-Shot LLM Response Without Annotations, fonttitle=\centering, before upper=\raggedright, boxrule=0.8mm, sharp corners=southwest, left=2pt, right=2pt, top=2pt, bottom=2pt, boxsep=1pt]
% \footnotesize
% \textbf{Correctness: False} \\
% The provided program does not correctly solve the problem as described. The approach taken in the program does not account for the fact that when Kirill swaps with a person, he also has to pay additional coins to all the people between the two positions. The program only considers the minimum of $a[i]$ and $b[i]$ for each position, which is not sufficient to determine the minimum cost for Kirill to move to one of the first $m$ positions.
% \end{tcolorbox}


% \begin{tcolorbox}[colframe=nicered!90, colback=gray!3, title=Tester Response, fonttitle=\centering, before upper=\raggedright, boxrule=0.8mm, sharp corners=southwest, left=2pt, right=2pt, top=2pt, bottom=2pt, boxsep=1pt]
%         \footnotesize
%             <The rest of the test cases are ommited for brevity>
          
% \textbf{Input} = 
% 10 1 \\
% 10 9 8 7 6 5 4 3 2 1\\
% 1 1 1 1 1 1 1 1 1 1\\
% \textbf{Output} = 10 \\
% \# \highlight{\textcolor{black}{This test case is wrong. Krill must pay 10 to swap with the first }}\\
% \# \highlight{\textcolor{black}{and 1 at each of the other 9 people in front of him so 19 not 10}}
% \end{tcolorbox}
% \begin{tcolorbox}[colframe=softblue, colback= softblue!5, title=Annotation with Purpose of \texttt{c} list , fonttitle=\centering, before upper=\raggedright, boxrule=0.8mm, sharp corners=southwest, left=2pt, right=2pt, top=2pt, bottom=2pt, boxsep=1pt]
% \footnotesize
% \# for each position `i`, `c[i]` is the minimum cost to move past the person at position `i` (the minimum of either `a[i]` or `b[i]`)
% \end{tcolorbox}


% \begin{tcolorbox}[colframe=nicegreen!90, colback=gray!3, title=LLM Response With Execution State Annotations, fonttitle=\centering, before upper=\raggedright, boxrule=0.8mm, sharp corners=southwest, left=2pt, right=2pt, top=2pt, bottom=2pt, boxsep=1pt]
% \footnotesize
% \textbf{Correctness: True} \\
% Let's break down the steps: \\
% 1. \textbf{Compute \texttt{c}}: For each position $i$, $c[i]$ is the minimum cost to move past the person at position $i$ (either $a[i]$ or $b[i]$). \\
% 2. \textbf{Compute \texttt{suff}}: \texttt{suff[i]} is the suffix sum of $c$ from position $i$ to the end of the queue. This helps in quickly calculating the total cost to move past all people from any position $i$ to the end. \\
% 3. \textbf{Determine the minimum cost}: For each position $i$ from $0$ to $k-1$, the cost to move to position $i$ is $a[i] + \texttt{suff[i + 1]}$. The function then selects the minimum of these costs.
% \end{tcolorbox}

% \changes{\textit{\textbf{Note}}: Problem description, prompt, and response content have been truncated for brevity.}
% \end{minipage}
% \end{figure}

\subsection{Natural Language and LLMs for State Tracking}
\label{appendix:NL}

We explore the effectiveness of Qwen2.5-72B-Instruct in inferring the final state (postcondition) of a program segment by leveraging semantic context and real-world knowledge. Unlike formal verification methods that rely on explicit rules and symbolic execution, Qwen2.5-72B-Instruct can generate intuitive and human-readable explanations of code behavior. This approach enhances analysis precision by utilizing variable names and structure to improve reasoning. While it does not provide formal correctness guarantees, it acts as a practical fuzzy correctness classifier, making it useful for tasks like bug detection and code generation.


In Figures \ref{fig:unique_email},\ref{fig:first_name} we present two more examples to complement that in \Cref{sec:motivating_NL}. These examples illustrate how an LLM can concisely infer express program states by leveraging natural language. We believe this ability of LLM to succinctly capture program behavior using user-centric terms helps \projName to accurately classify program correctness by bridging the gap between program semantics and informal requirements.
\begin{figure}[t]
    \centering
    \includegraphics[width=0.8\linewidth]{example2-sec5_blue.png}
    \caption{Example 2 - Unique Email Domains \label{fig:unique_email}}
    \vspace{-2mm}
\end{figure}
\begin{figure}[t]
    \centering
    \includegraphics[width=0.8\linewidth]{example3-sec5_blue.png}
    \caption{Example 3 - First name extraction \label{fig:first_name}}
    \vspace{-2mm}
\end{figure}


\subsection{Dataset Details}

\begin{figure}
    \centering
    \includegraphics[width=0.9\columnwidth]{rating_dis.png}  
    \caption{The Frequency Histogram of Rating for \benchName with x-axis representing the rating intervals of the problems and y-axis representing the number of problem-code pairs within each rating interval in the dataset.}
    \label{fig:rating_distribution}
    \vspace{-3mm}
\end{figure}



\label{dataset_details}

To better illustrate the difficulty of problems in \benchName, we use the \emph{rating} tags for problems provided by Codeforces as an indicator. Initially, rating was solely an indicator of user rankings on programming competition platforms. Subsequently, to better assist users in selecting appropriate problems for practice, Codeforces began assigning each problem a rating to represent its difficulty, which is determined by the solving statistics of participants for this problem~\cite{codeforces2019problem}. Formally, a problem's rating $x$ quantifies that participants with a rating of $x$ have a 50\% probability of solving it on their first encounter.

Rating for problems ranges from 800 to 3500 and serves as a difficulty coefficient indicator, where a higher rating indicates greater problem difficulty. During the problem collection process, we also gathered the ratings of each problem and present the rating distribution of our dataset. As shown in Figure~\ref{fig:rating_distribution}, our dataset has an average rating of 1206. According to the statistics and evaluations from CodeElo\cite{quan2025codeelobenchmarkingcompetitionlevelcode}, a problem with a rating of 1206 indicates that approximately 60\% of all participants on Codeforces find it challenging (with less than a 50\% probability of solving it on their first encounter), and only 2 of the 33 tested LLMs achieved a higher rating than 1206, which is difficult for LLMs to simply retrieve. 


Additionally, as can be seen from the figure, the data presents a skewed distribution of "higher on the left and lower on the right". The frequency in low rating intervals (such as 800–1000) is significantly higher, and as the rating increases, the frequency generally shows a downward trend. The main reason is that as the difficulty of the problem increases (the rise in rating), the number of people who solve this problem (or even the number of people who attempt to solve it) decreases. Therefore, there are fewer submission records available for sampling, and even fewer submission records that meet our sampling strategy.

The scraped problems were originally in raw HTML format. We parsed out all the useful content w.r.t. the problem and converted it into Markdown format (to be consistent with our prompt), which may include the problem description, input format, output format, examples, notes, and so on, to ensure the sufficiency of the problem information and details.

For a problem, we filtered out all its Python submissions. The specific strategy for selecting submissions is as follows:
\begin{itemize}
    \item Include three consecutive incorrect/correct code pairs (each pair comes from the same user). If there are less than three, include all that are available. If the problem has not been solved by anyone, skip it. 
    \item If available, include one code with a test passed rate of 50\% (or close to 50\% due to an odd number of test cases), and one code with only a single failing test case. 
\end{itemize}

\changes{Furthermore, we provide a detailed comparison in Table~\ref{tab:dataset_comparison} of our \benchName~dataset to two prominent datasets: LiveCodeBench~\cite{jain2024livecodebenchholisticcontaminationfree} and APPS~\cite{apps}, based on key metrics such as solution complexity, code length, loop depth, and description verbosity.}

\setlength{\tabcolsep}{4pt}
\begin{table}[h]
\centering
\begin{tabular}{|l|c|c|c|c|c|c|}
\hline
\changes{\textbf{Dataset}} & \changes{\textbf{LOC}} & \changes{\textbf{LD}} & \changes{\textbf{DL}} & \changes{\textbf{CC}} & \changes{\textbf{HD}} & \changes{\textbf{HV}} \\
\hline
\changes{\textbf{\benchName}}      & \changes{25.1}  & \changes{1.29} & \changes{2156.83} & \changes{7.77}  & \changes{28.8}  & \changes{1063.2} \\
\changes{\textbf{LiveCodeBench}}   & \changes{9.95}  & \changes{1.01} & \changes{42.35}   & \changes{4.18}  & \changes{21.81} & \changes{416.2}  \\
\changes{\textbf{APPS}}            & \changes{24.03} & \changes{1.29} & \changes{1416.06} & \changes{7.63}  & \changes{34.12} & \changes{1103.6} \\
\hline
\end{tabular}
\caption{\changes{Comparison of dataset complexity based on Lines of Code (LOC), Loop Depth (LD), Description Length (DL), Cyclomatic Complexity (CC), Halstead Difficulty (HD), and Halstead Volume (HV). Higher CC, HD, and HV values indicate more complex programs.}}
\label{tab:dataset_comparison}
\end{table}
\setlength{\tabcolsep}{6pt}

\subsection{Rerun Strategy and Variance Control}
\label{appendix:reruns}

\changes{To ensure statistical reliability given LLM response variability, we computed the number of reruns required to achieve a high-confidence estimate with low error. Our approach follows the framework of Ouyang et al.~\cite{ouyang2025empirical} and the confidence-weighted aggregation method of Xiong et al.~\cite{xiong2024llmsexpressuncertaintyempirical}.}

\paragraph{\changes{Confidence Estimation.}}  
\changes{Given a binary prediction (Correct or Incorrect), the LLM assigns a log-probability \( \log P(Y) \), which is converted to confidence via \( \exp(\log P(Y)) \). We aggregate confidences across \( M \) reruns using the Avg-Conf method:}
\[
\changes{
C_{\text{avg-conf}} = \frac{\sum_{i=1}^{M} \mathbb{I}(Y_i = Y) \cdot C_i}{\sum_{i=1}^{M} C_i}
}
\]
\changes{where \( Y_i \) is the $i$-th response, \( C_i \) its confidence score, and \( \mathbb{I}(Y_i = Y) \) is an indicator for agreement with the majority response.}

\paragraph{\changes{Rerun Count Calculation.}}  
\changes{Let \( P \) be the estimated consistency of the LLM across reruns, \( X \) the number of evaluation instances, and \( R \) the number of reruns. We aim to bound the margin of error \( \epsilon \) for a desired confidence level \( C \), modeled with a binomial distribution:}
\[
\changes{
\epsilon = Z \sqrt{\frac{P(1-P)}{N}}, \quad N = R \cdot X, \quad R = \frac{Z^2 P(1-P)}{X \cdot \epsilon^2}
}
\]
\changes{where \( Z \) is the critical value for the desired confidence level (e.g., \( Z = 2.33 \) for 98\%).}

\paragraph{\changes{Final Parameters.}}  
\changes{
Using:
\begin{itemize}
    \item Dataset size: \( X = 645 \),
    \item Observed average confidence: \( C = 0.963 \),
    \item Desired confidence level: 98\% (\( Z = 2.33 \)),
    \item Error margin: \( \epsilon = 0.01 \),
\end{itemize}
we compute \( R \approx 3 \), meaning all experiments are repeated three times to achieve statistically reliable results.
}

\subsection{Example: Tester Misclassification due to multiple correct solutions}
\label{tester_multiple_correct}

\changes{This example from our dataset \benchName{}, as seen on Figure \ref{multiple_correct}, highlights a common limitation of the LLM-generated tester. Specifically, it shows how problems with multiple valid outputs can lead the tester to incorrectly reject a correct implementation.}
\begin{figure}[h]


\caption{\changes{Example of problem with multiple correct solutions}}
    \centering
    \label{multiple_correct}
    \begin{minipage}[t]{0.31\textwidth}
        \vspace{0pt}
        \begin{tcolorbox}[colframe=darkgreen!100, colback=gray!3, title=Problem Description, fonttitle=\centering, boxrule=0.7mm, sharp corners=southwest, boxsep=1pt]
        \footnotesize
Given integer $x$, find any $y$ ($1 \le y < x$) such that
\[
\gcd(x,y) + y
\]
is maximized.

Multiple $y$ may be equally valid.

        \end{tcolorbox}
    \end{minipage}
    \hfill
    \begin{minipage}[t]{0.31\textwidth}
        \vspace{0pt}
        \begin{tcolorbox}[colframe= softblue, colback=gray!3, title=Tester Generated Test Case, fonttitle=\centering, boxrule=0.7mm, sharp corners=southwest, boxsep=1pt]
        \footnotesize

        
\textbf{Input:}
\begin{verbatim}
1
10
\end{verbatim}
\textbf{Output:}
\begin{verbatim}
5
\end{verbatim}

\[
\gcd(10,5)=5
\quad 5+5=10
\]

Tester selected $y=5$ as the solution.
        \end{tcolorbox}
    \end{minipage}
    \hfill
    \begin{minipage}[t]{0.31\textwidth}
        \vspace{0pt}
        \begin{tcolorbox}[colframe=nicered!90, colback=gray!3, title=Tester Classification, fonttitle=\centering, boxrule=0.7mm, sharp corners=southwest, boxsep=1pt]
        \footnotesize
        \textbf{Incorrect}\\
Our implementation produced:

\begin{verbatim}
9
\end{verbatim}

\[
\gcd(10,9)=1,\;1+9=10
\]

Also optimal, but tester flagged as \textbf{incorrect} due to strict expected match.
        \end{tcolorbox}
    \end{minipage}
\end{figure}



\subsection{Example : Impact of Incorrect Test Cases on LLM-Driven Code Refinement based on Agent Coder}
\label{ex_agent_coder_test}

\changes{In our experiments with Llama3.1, we observed a striking example of how incorrect test cases can mislead LLM-based program synthesis or correction systems, resulting in previously correct implementations being erroneously altered.}

\paragraph{\changes{Context.}} \changes{The underlying programming problem involves a little boy, Nikita, who builds a tower using cubes. The process follows two simple rules:}
\begin{itemize}
    \item \changes{In each of $n$ moves, Nikita must either \emph{add exactly one cube} on top of the tower or \emph{remove exactly one cube} from the top.}
    \item \changes{The tower starts with $0$ cubes, and it is never allowed to have a negative height—removing a cube from an empty tower is invalid.}
\end{itemize}

\changes{Given this setup, the task is to determine whether it is possible, after exactly $n$ moves, to end up with a tower that has exactly $m$ cubes.}

\changes{Two simple conditions fully characterize the problem:}
\begin{itemize}
    \item \changes{$m \le n$, since each move changes the height by at most $1$.}
    \item \changes{$(n - m)$ is even, ensuring moves can be paired into balanced ``add and remove'' operations to reach exactly $m$.}
\end{itemize}
\changes{Thus, the tower ends with $m$ cubes after $n$ moves if and only if:}
\[
\changes{
\text{possible} \Longleftrightarrow \left(m \le n \;\;\wedge\;\; (n - m) \bmod 2 = 0\right)
}
\]
\changes{The LLM produced an implementation which faithfully reflects this reasoning:}
\begin{verbatim}
def solve(n, m):
    if m > n:
        return "No"
    if (n - m) % 2 == 0:
        return "Yes"
    else:
        return "No"
\end{verbatim}

\changes{\paragraph{Faulty Test Cases.} LLaMA3.1-70B generated test cases that the expected output was incorrectly calculated by the model and they were subsequently used as feedback to refine the code.Two such test cases are:}
\begin{quote}
\begin{itemize}
    \item \changes{\texttt{Input: 1 0}, \texttt{Expected Output: Yes}}
    \item \changes{\texttt{Input: 5 0}, \texttt{Expected Output: Yes}}
\end{itemize}
\end{quote}
\changes{which contradict the problem’s constraints: for example, with \texttt{1 0}, it is impossible to end with 0 cubes after a single move since removing a cube is invalid on an empty tower.}
\paragraph{\changes{Misguided Corrections.}} 
\changes{The faulty test cases were executed on the original correct implementation, resulting in output mismatches. Interpreting these mismatches as genuine errors, the LLM attempted to fix the code to satisfy the flawed expectations—introducing ad-hoc conditions that deviated from the correct logic:}

\begin{verbatim}
if (n - m) % 2 == 0 or ((n - m) % 2 == 1 and m != 0):
    return "Yes"
\end{verbatim}
\changes{which breaks the parity logic and admits invalid solutions, thereby turning correct code into incorrect code.}

\subsection{Example: HoarePrompt-Guided Correction for Code Generation based on Agent Coder}
\label{ex_hoareprompt_beautiful_pairs}

\paragraph{\changes{Context.}}  
\changes{Count all pairs \(\langle i,j\rangle\) in an array \(a\) so that}
\[
\changes{
a_i + a_j \equiv 0 \pmod{x}
\quad\text{and}\quad
a_i - a_j \equiv 0 \pmod{y}.
}
\]

\paragraph{\changes{Faulty Implementation.}}
\begin{verbatim}
for num in a:
    rx = num % x
    ry = num % y
    remainder = (rx + ry) % y
    complement = ((x - rx) % x + ry) % y
    beautiful_pairs += cnt[complement]
    cnt[remainder] += 1
\end{verbatim}

\paragraph{\changes{HoarePrompt Annotation.}}  
\changes{HoarePromptwith Llama3.1,highlights the line}
\begin{verbatim}
    remainder = (rx + ry) % y
\end{verbatim}
\changes{with the note}  
\changes{\emph{the remainder is the remainder of the sum of num divided by x plus the remainder of num divided by y, divided by y}}

\paragraph{\changes{Feedback.}}  
\changes{“The program contains a bug. This bug stems from merging \(\texttt{num}\bmod x\) and \(\texttt{num}\bmod y\) into one \(\bmod y\) expression, dropping the independent \(\bmod x\) check.”}

\paragraph{\changes{Corrected Implementation.}}
\begin{verbatim}
cnt = defaultdict(int)
for num in a:
    r_x = num % x
    r_y = num % y
    need_x = (-r_x) % x
    beautiful_pairs += cnt[(need_x, r_y)]
    cnt[(r_x, r_y)] += 1
\end{verbatim}
\changes{Here each element is bucketed by the pair \((r_x,r_y)\), so “sum $\equiv$ 0 mod x” and “difference $\equiv$ 0 mod y” are enforced independently.}

\subsection{MCC Calculation and Justification}
\label{appendix:mcc}

To evaluate the quality of binary correctness classification, we adopt the \textbf{Matthews Correlation Coefficient (MCC)} as our primary metric. Given the confusion matrix entries—true positives (TP), true negatives (TN), false positives (FP), and false negatives (FN)—the MCC is computed as:

\[
\text{MCC} = \frac{(\text{TP} \times \text{TN}) - (\text{FP} \times \text{FN})}{\sqrt{(\text{TP} + \text{FP})(\text{TP} + \text{FN})(\text{TN} + \text{FP})(\text{TN} + \text{FN})}}
\]

The MCC yields a score in \([-1, 1]\), where:
\begin{itemize}
    \item \(1\) indicates perfect prediction,
    \item \(0\) indicates random prediction,
    \item \(-1\) indicates total disagreement.
\end{itemize}

We chose MCC over commonly used alternatives such as ROC AUC or accuracy because MCC takes into account \emph{all four} confusion matrix categories, and it remains informative even under class imbalance. This makes it particularly suitable for correctness classification, where false negatives (missing bugs) and false positives (flagging correct code as incorrect) are both costly.  As shown by Chicco and Jurman~\cite{mcc}, ROC AUC can yield misleadingly high scores even for classifiers with low precision or poor practical performance. Moreover, ROC AUC considers performance over all thresholds, not the concrete operating point used for binary decisions. In contrast, MCC directly reflects classifier behavior at the applied threshold and ensures high values only when \emph{sensitivity, specificity, precision, and NPV} are all high.
Therefore, we follow their recommendation that MCC should replace ROC AUC as the standard metric in binary classification settings~\cite{mcc}, especially in sensitive domains like program verification, where balanced correctness evaluation is critical.



\onecolumn  % Switch to single-column format

\section{LLM Prompts}

In this section of the Appendix, we include all prompts provided to the LLMs during the evaluation and experiments in our study. Where FSL was employed, the FSL examples are also included.

\subsection{Zero-Shot Chain-of-Thought (CoT) Prompt}
\label{0-shot-COT-prompt}
\begin{tcolorbox}[ colback=gray!10, colframe=black, arc=5pt]
\begin{lstlisting}[language=, basicstyle=\ttfamily\footnotesize]
Your task is to determine if a given Python program is correct based on the provided problem description. Assume valid inputs as described in the problem description.

First, explain your reasoning step by step, then reply with:  
- Correctness: `True` if the given program is correct.  
- Correctness: `False` if the given program is incorrect.


Problem:  
{decription}

Program:  
{code}

# Your response:  
Reasoning:
Correctness: **True** or **False**

\end{lstlisting}
\end{tcolorbox}


\subsection{Vanilla Prompt (No Chain-of-Thought)}
\label{vanilla-prompt}
\begin{tcolorbox}[ colback=gray!10, colframe=black, arc=5pt]
\begin{lstlisting}[language=, basicstyle=\ttfamily\footnotesize]
Your task is to determine if a given Python program is correct based on the provided problem description. Assume valid inputs as described in the problem description.

Reply with:  
- Correctness: `True` if the given program is correct.  
- Correctness: `False` if the given program is incorrect.


Problem:  
{description}

Program:  
{code}

# Your response:  
Correctness: **True** or **False**

\end{lstlisting}
\end{tcolorbox}


\subsection{Prompt for Input-Output Test Case Generation}
\begin{tcolorbox}[ colback=gray!10, colframe=black, arc=5pt]
\begin{lstlisting}[language=, basicstyle=\ttfamily\footnotesize]
**Role**: As a tester, your task is to create comprehensive test cases for the following coding problem. These test cases should encompass Basic and Edge scenarios to ensure the code's robustness, reliability, and scalability.

**Problem Description**:  
{description}

**1. Basic Test Cases**:  
- **Objective**: To verify the fundamental functionality of the `has_close_elements` function under normal conditions.

**2. Edge Test Cases**:  
- **Objective**: To evaluate the function's behavior under extreme or unusual conditions.

**Instructions**:  
- Implement a comprehensive set of test cases following the guidelines above.  
- Ensure each test case is complete (no omission) and well-documented with comments explaining the scenario it covers.  
- Pay special attention to edge cases as they often reveal hidden bugs.  
- Do not repeat, do not summarize.  

All test cases you give need to strictly follow the problem description and format like this:

# Test 1  
**Input**:  
```

```
**Output**:  
```

```

\end{lstlisting}
\end{tcolorbox}

\subsection{HoarePrompt Correctness Classification Prompt}
\label{correctness-single-func-prompt}
\begin{tcolorbox}[ colback=gray!10, colframe=black, arc=5pt]
\begin{lstlisting}[language=, basicstyle=\ttfamily\footnotesize]
Your task is to determine if a given Python program is correct based on the problem description and the execution states of the program provided as comments. Assume valid inputs as described in the problem description.

First, explain your reasoning, then reply with:  
- Correctness: `True` if the given program is correct.  
- Correctness: `False` if the given program is incorrect.


**# Problem:**  
{description}

**# Annotated Program:**  
{annotated_program}

**# Your response:**  
Reasoning:  
Correctness: **True** or **False**

\end{lstlisting}
\end{tcolorbox}

\subsection{HoarePrompt Correctness Classification for Multi-Function Programs}
\label{correctness-multi-func-prompt}
\begin{tcolorbox}[ colback=gray!10, colframe=black, arc=5pt]
\begin{lstlisting}[language=, basicstyle=\ttfamily\footnotesize]
Your task is to determine if a given Python program is correct based on the problem description and the execution states of the program provided as comments. Assume valid inputs as described in the problem. The program is made of multiple functions and the program is **correct** only if all its functions together meet the problem description.

First, explain your reasoning, then reply with:  
- Correctness: `True` if the given program is correct.  
- Correctness: `False` if the given program is incorrect.


**# Problem:**  
{description}

**# Annotated Functions:**  
{functions}

**# Your response:**  
Reasoning:  
Correctness: **True** or **False**

\end{lstlisting}
\end{tcolorbox}
\flushbottom

\subsection{Precondition Extraction Prompt}
\label{pre_prompt}
\begin{tcolorbox}[ colback=gray!10, colframe=black, arc=5pt]
\begin{lstlisting}[language=, basicstyle=\ttfamily\footnotesize]
You are given a programming problem description and a function definition for a function that solves this problem. From the problem description, extract a description of the values of the program's input variables and the relationships between these variables. We refer to this description as the **precondition**. Print the precondition following the word **"Precondition"**, and surround it with double asterisks (**). Follow these examples:

### Example 1  
**Problem description:** Write a function to find the minimum cost path to reach (m, n) from (0, 0) for the given cost matrix cost[][] and a position (m, n) in cost[][].  
**Function definition:**
```python
def min_cost(cost, m, n):
```
**Precondition:** **cost is a 2D list of non-negative integers, m and n are non-negative integers such that 0 <= m < len(cost) and 0 <= n < len(cost[0]).**

### Example 2  
**Problem description:** Write a function to find the similar elements from the given two tuple lists.  
**Function definition:**
```python
def similar_elements(test_tup1, test_tup2):
```
**Precondition:** **test_tup1 and test_tup2 are tuples.**

### Example 3  
**Problem description:** Write a Python function to identify non-prime numbers.  
**Function definition:**
```python
def is_not_prime(n):
```
**Precondition:** **n is an integer greater than 1.**

### Example 4  
**Problem description:** Write a function to find the largest integers from a given list of numbers using the heap queue algorithm.  
**Function definition:**
```python
def heap_queue_largest(nums, n):
```
**Precondition:** **nums is a list of integers, and n is a non-negative integer such that 0 <= n <= len(nums).**

### Example 5  
**Problem description:** Write a function to find the number of ways to fill it with 2 x 1 dominoes for the given 3 x n board.  
**Function definition:**
```python
def count_ways(n):
```
**Precondition:** **n is a non-negative integer.**

### Example 6  
**Problem description:** Find the average of the positive integers in a list.  
**Function definition:**
```python
def func_1(numbers):
```
**Precondition:** **numbers is a list of integers.**

### **Your Task**

**Problem description:**  
{description}

**Function definition:**
```python
{program}
```

\end{lstlisting}
\end{tcolorbox}
\clearpage
\subsection{Precondition Extraction Prompt for Multi-Function Programs}
\label{pre_multi_prompt}
\begin{tcolorbox}[ colback=gray!10, colframe=black, arc=5pt]
\begin{lstlisting}[language=, basicstyle=\ttfamily\footnotesize]
You are given a programming problem description and a function that contributes to the solution of this problem. The total solution comprises multiple functions, and this is just one of them. 
From the problem description, and based on the variables used in the signature of this specific function, extract a description of the values of the variables in the function signature and the relationship between them. We refer to this description as the **precondition**. Print the precondition following the word **Precondition**, and surround it with double asterisks (**). Follow these examples:
Remember, the function given may not solve the problem directly but may perform a side functionality that contributes to the total solution. Include information only about the variables in the function signature.

---

### Example 1  
**Problem description:** Write a function to find the minimum cost path to reach (m, n) from (0, 0) for the given cost matrix cost[][] and a position (m, n) in cost[][].  
**Program:**
```python
def min_cost(cost, m, n):
    tc = [[0 for x in range(C)] for x in range(R)]
    tc[0][0] = cost[0][0]
    for i in range(1, m+1):
        tc[i][0] = tc[i-1][0] + cost[i][0]
    for j in range(1, n+1):
        tc[0][j] = tc[0][j-1] + cost[0][j]
    for i in range(1, m+1):
        for j in range(1, n+1):
            tc[i][j] = min(tc[i-1][j-1], tc[i-1][j], tc[i][j-1]) + cost[i][j]
    return tc[m][n]
```
**Precondition:** **cost is a 2D list of non-negative integers, m and n are non-negative integers such that 0 <= m < len(cost) and 0 <= n < len(cost[0]).**

---

### Example 2  
**Problem description:** Write a function to find the similar elements from the given two tuple lists.  
**Program:**
```python
def are_similar(elem, elem1):
    if elem == elem1:
        return True
    else:
        return False
```
**Precondition:** **elem1 and elem are values of any type and value.**

---

### Example 3  
**Problem description:** Write a Python function to identify if 2 consecutive integers in a list are not prime.  
**Program:**
```python
import math
def is_not_prime(n):
    result = False
    for i in range(2, int(math.sqrt(n)) + 1):
        if n % i == 0:
            result = True
    return result
```
**Precondition:** **n is an integer greater than 1.**

\end{lstlisting}
\end{tcolorbox}

\begin{tcolorbox}[ colback=gray!10, colframe=black, arc=5pt]
\begin{lstlisting}[language=, basicstyle=\ttfamily\footnotesize]
### Example 4 
**Problem description:** Write a function to find the largest integers from a given list of numbers using the heap queue algorithm.  
**Program:**
```python
import heapq as hq
def heap_queue_largest(nums, n):
    largest_nums = hq.nlargest(n, nums)
    return largest_nums
```
**Precondition:** **nums is a list of integers, and n is a non-negative integer such that 0 <= n <= len(nums).**

---

### **Your Task**

**Problem description:**  
{description}

**Program:**
```python
{program}
```

\end{lstlisting}
\end{tcolorbox}
\clearpage
\flushbottom
\subsection{Simple Statement Execution Prompt}
\label{simple_statement_prompt}
\begin{tcolorbox}[ colback=gray!10, colframe=black, arc=5pt]
\begin{lstlisting}[language=, basicstyle=\ttfamily\footnotesize]
You have been assigned the role of a program executor, responsible for simulating the execution of Python code. You will be provided with an initial state and a Python code snippet. You need to provide the output state after running the Python code based on the initial state. Please avoid describing how the program runs. When a variable has a specific value, use that specific value directly for calculations. If a return takes place, make sure to always mention that a value or variable has been returned in the output state. You must adhere to the text format: **Output State: `output state`**.

Include all the information of the precondition that is still valid after the code has been executed. Just update the values of the variables that have been changed by the code.

I am giving you some examples to understand the task better. Then I am giving you your task.

---

### Example 1  
**Initial State:** `str` is a string  
```python
n = int(input())
```
**Output State:** **`str` is a string, `n` is an input integer**

---

### Example 2  
**Initial State:** Variables can hold any values  
```python
i += 1
```
**Output State:** **variable `i` is increased by 1**

---

### Example 3  
**Initial State:** `n` is either 3 or 5  
```python
m = n + 1
```
**Output State:** **`n` is either 3 or 5; `m` is either 4 or 6**

---

### Example 4  
**Initial State:** `x` is a positive integer; if `x` is greater than 10, then `z = 0`, else `z = 1`.  
```python
y = -z
```
**Output State:** **`x` is a positive integer, if `x` is greater than 10 then `z=0` and `y=0`, else `z=1` and `y=-1`**

---

### Example 5  
**Initial State:** `total` is 0, `i` is 1  
```python
break
```
**Output State:** **`total` is 0, `i` is 1 and we break out of the most internal loop or if statement.**

\end{lstlisting}
\end{tcolorbox}

\begin{tcolorbox}[ colback=gray!10, colframe=black, arc=5pt]
\begin{lstlisting}[language=, basicstyle=\ttfamily\footnotesize]
### Example 6  
**Initial State:** `total` is positive, `num` is negative, `x` is 0  
```python
x = total - num
```
**Output State:** **`total` is positive, `num` is negative, `x` is a positive value equal to `total - num`.**

### **Your Task**

**Initial State:**  
{pre}

```python
{program}
```

\end{lstlisting}
\end{tcolorbox}


\subsection{Compound Simple Statements Execution Prompt}
\label{compund statement prompt}
\begin{tcolorbox}[ colback=gray!10, colframe=black, arc=5pt]
\begin{lstlisting}[language=, basicstyle=\ttfamily\footnotesize]
You have been assigned the role of a program executor, responsible for simulating the execution of Python code. You will be provided with an initial state and a Python code snippet consisting of multiple lines. Your task is to execute all the lines in sequence and provide the output state after the entire code block has been run. Avoid describing how the program runs step-by-step for individual lines but instead focus on the combined effect of all lines. When a variable has a specific value, use that specific value directly for calculations.

Include all the information from the precondition that remains valid after the code execution and update the values of any variables that are modified by the code. Provide the final state, including the state of all the variables after the execution of the code snippet. Use the text format: **Output State: `output state`**.

Here are some examples to help you understand the task:

---

### Example 1  
**Initial State:** `a` is 5, `b` is 3  
```python
c = a + b
d = c * 2
```
**Output State:** **a is 5, b is 3, c is 8, d is 16**

---

### Example 2  
**Initial State:** `x` is a positive integer  
```python
n = int(input())
x += n
```
**Output State:** **x is a positive integer equal to its original value plus `n`, `n` is an integer**

---

### Example 3  
**Initial State:** `n` is 3, `total` is 1  
```python
total += n
pass
n = max(total, n)
```
**Output State:** **n is 4, total is 4**

---

### **Your Task**

**Initial State:**  
{pre}

```python
{program}
```

Now, please analyze the entire block of code and provide the final output state. List the impact of all lines on the program, check the previous values of affected variables, and calculate the states of the variables after the code executes. Be as specific as possible, combining changes from all lines into a single coherent final state. Include all valid information from the precondition and update only what is modified by the code.

In your response, strictly use the format: **Output State: `the output state you calculate`.**, and describe this output state in natural language easily understandable by humans.

\end{lstlisting}
\end{tcolorbox}



\subsection{Print Commands Execution Prompt}
\label{print_prompt}
\begin{tcolorbox}[ colback=gray!10, colframe=black, arc=5pt]
\begin{lstlisting}[language=, basicstyle=\ttfamily\footnotesize]
You will be given an **initial state** (precondition) and a **Python code snippet** containing a `print` statement. Your task is to **determine exactly what will be printed** when the statement executes.

- If a variable or object has a known **explicit value**, use that value in the output.  
- If a variable or object is defined by a **formula or condition**, describe its value using the given information.  
- Always provide the most **precise** description possible based on the precondition.
- Format the final output as: **Output: `what is printed`**.

I am giving you some examples to understand the task better. Then I am giving you your task.

---

### Example 1  
**Initial State:** `arr` is a list containing 1, 2, 3, 4, 5, and `sum` is the sum of all elements in the list `arr`.  
```python
print(arr[2], sum)
```
**Output:** **3, sum (where sum is the sum of all elements in list `arr`)**

---

### Example 2  
**Initial State:** The `points` list is a list of points. The `shoelace_sum` is the sum of all terms calculated as \(x_1 * y_2 - y_1 * x_2\) for each consecutive pair of points in the `points` list. The `area` is the absolute value of `shoelace_sum` divided by 2.0. `i` is equal to `len(points) - 2`, and `x1` is the first element of `points[i]`, `y1` is the second element of `points[i]`, while `x2` is the first element of `points[i + 1]`, and `y2` is the second element of `points[i + 1]`.
```python
print(area)
```
**Output:** **area (where area is the area of the polygon formed by the points in the `points` list)**

---

### Example 3  
**Initial State:** `balances` is a list of integers, `A` is the first element of the `balances` list, `B` is the second element of the `balances` list, and `amount` is an integer less than or equal to `A`.
```python
print(f"The amount {amount} is less than or equal to {A}")
```
**Output:** **The amount [amount] is less than or equal to [A] (where `amount` is the value of `amount` and `A` is the first element of the `balances` list)**

---

### **Your Task**

**Initial State:**  
{pre}

```python
{program}
```

Now, please think step by step. Based on the precondition that describes the state of the program, variables, objects, etc., before the code is executed, calculate what will be printed when the `print` statement is executed. Explain the values of the variables, objects, etc., that are printed.

Use natural language to describe the output, ensuring it is easily understandable by a human. Strictly adhere to the format **Output: `what is printed`**.

\end{lstlisting}
\end{tcolorbox}

\clearpage
\flushbottom
\subsection{Return Command Execution Prompt}
\label{return_prompt}
\begin{tcolorbox}[ colback=gray!10, colframe=black, arc=5pt]
\begin{lstlisting}[language=, basicstyle=\ttfamily\footnotesize]
You have been assigned the role of a program executor, responsible for simulating the execution of Python code. You will be provided with an **initial state** and a **Python code snippet**. Your task is to determine what the program **returns** based on the given initial state.
- Avoid describing how the program runs step by step.
- If a variable has a **specific value**, use that value directly.
- If a return statement is executed, **always mention the returned value explicitly** in the output state.
- Adhere to the text format: **Output State: `output state`**.

I am giving you some examples to understand the task better. Then I am giving you your task.
### Example 1  
**Initial State:** `str` is a string, 'str' has 3 or more characters  
```python
return str
```
**Output State:** **The program returns string `str` that has 3 or more characters**

### Example 2  
**Initial State:** `numbers` is an empty list, `total` is the sum of all positive integers that were in the original `numbers` list, `count` is the number of positive integers that were in the original `numbers` list, `average` is equal to `total / count`  
```python
return average
```
**Output State:** **The program returns `average`, which is equal to `total/count`, where `total` is the sum of all positive integers in the original `numbers` list, and `count` is the number of positive integers in the original `numbers` list.**

### Example 3  
**Initial State:** `n` is either 3 or 5  
```python
return n + 1
```
**Output State:** **The program returns 4 or 6**

### Example 4  
**Initial State:** `x` is 1, `y` is greater than 3, `z` is 0  
```python
return y + x
```
**Output State:** **The program returns `y + 1`, where `y` is greater than 3**

### Example 5  
**Initial State:** `count` contains the number of numbers greater than 1 in the list `numbers`, `numbers` is a list of integers, `total` is 0  
```python
return count
```
**Output State:** **The program returns the number of integers greater than 1 in list `numbers`**

### **Your Task**

**Initial State:**  
{pre}

```python
{program}
```

Now, please think step by step: List the impact of the code on the program, check the previous values of the affected variables, and then calculate what the program returns. Any variable or value that is included in the return should be fully described with all available information.

Strictly adhere to the format: **Output State: `the output state you calculate`**, and describe this output state in **natural language easily understandable by humans**.

\end{lstlisting}
\end{tcolorbox}

\subsection{Try-Except Statement Execution Prompt}
\label{try_except_prompt}
\begin{tcolorbox}[ colback=gray!10, colframe=black, arc=5pt]
\begin{lstlisting}[language=, basicstyle=\ttfamily\footnotesize]
You have been assigned the role of a program verifier, responsible for summarizing the state of the function after executing a Python `try` statement. You will be provided with the **final state of the program** after executing the `try` block, along with the **changes in the program after executing one or more `except` blocks** in any situation. Your task is to **combine this information to summarize the program's state after the complete execution of the `try` statement**.

- If a `return` statement is executed, **always include it in the output state**.
- Adhere to the text format: **Output State: `output state`**.

I am giving you some examples to understand the task better. Then I am giving you your task.

---

### Example 1  
**Initial State:** `a` is an integer, `b` is an integer.  
```python
try:
    result = a / b
    return result
except ZeroDivisionError:
    return None
```
**Output state after executing the `try` block:** `a` is an integer, `b` is an integer, `result` is `a` divided by `b`, and the function returns `result`.

**Output state after executing the `except` block:** If a `ZeroDivisionError` occurs (i.e., `b` is 0), the function returns `None`.

**Final Output State:**  
A `ZeroDivisionError` might be triggered at `result = a / b`. If `b` is 0, the function returns `None`. Otherwise, the function returns the value of `a` divided by `b`.  
**Output State:** **`a` and `b` are integers. If `b` is zero, the function returns `None`, otherwise the function returns the value of `a` divided by `b`.**

---

### Example 2  
**Initial State:** `file_path` is a string that's supposed to be a path to a file.  
```python
def read_file(file_path):
    try:
        with open(file_path, 'r') as file:
            data = file.read()
            print("File content successfully read.")
            return data
    except FileNotFoundError:
        print("Error: The file was not found. Please check the file path.")
        return None
    except PermissionError:
        print("Error: You do not have permission to read this file.")
        return None
```
**Output state after executing the `try` block:** `file_path` is a string representing a file path, `data` contains the file content, and the function returns `data`.

**Output state after executing the `except` block(s):**
- If a `FileNotFoundError` occurs, the function prints an error message and returns `None`.
- If a `PermissionError` occurs, the function prints a different error message and returns `None`.

**Final Output State:**  
The program could raise a `FileNotFoundError` if the file is not found at the specified path or a `PermissionError` if the user does not have permission to read the file. If the file is found and the user has permission, the function reads the file content and returns it.  
**Output State:** **`file_path` is a string representing a file path. If the file exists and the user has permission, the function returns the file content; otherwise, it prints an error message and returns `None`.**

\end{lstlisting}
\end{tcolorbox}

\begin{tcolorbox}[ colback=gray!10, colframe=black, arc=5pt]
\begin{lstlisting}[language=, basicstyle=\ttfamily\footnotesize]
### **Your Task**

**Initial State:**  
{pre}

**Code for the try-except block:**
```python
{code}
```

**Output state after executing the `try` block:**  
{try_post}

**Output state after executing the `except` block(s):**  
{except_post}

Now, please think step by step: At which point in the program could such an exception occur? Summarize what the `try-except` statement accomplishes and what the program output state is after execution.

Strictly adhere to the format: **Output State: `the output state you calculate`**, and describe this output state in **natural language easily understandable by humans**.

\end{lstlisting}
\end{tcolorbox}

\subsection{Function Functionality and Return Summary Prompt}
\label{functionality_prompt}
\begin{tcolorbox}[ colback=gray!10, colframe=black, arc=5pt]
\begin{lstlisting}[language=, basicstyle=\ttfamily\footnotesize]
You have been assigned the task of a program verifier, responsible for describing the functionality of a Python function. You will be provided with the **constraints and relationships** between the input parameters, as well as the **resulting output** of the function based on these inputs. Your task is to **organize this information and describe the functionality of the function**.

- Avoid describing how the function operates (e.g., local variable initialization or step-by-step execution).
- Focus only on **what parameters the function accepts** and **what it returns**.
- If there are multiple return points in the function, we will split them into **cases** labeled sequentially (Case 1, Case 2, etc.).
- If one case is fulfilled, **none of the others are executed**.
- Adhere strictly to the format: **Functionality: `functionality description`**.

I am giving you some examples to understand the task better. Then I am giving you your task.


### Example 1  
**Parameter constraints and relationships:** `number` is an integer.  
```python
def func(number):
```
**Output:**  
- Case 1: If `number` is even, the function returns `True`.
- Case 2: If `number` is not even, the function returns `False`.

**Final Answer:**  
The function `func` accepts a parameter `number` and returns `True` if `number` is even. If `number` is not even, it returns `False`.  
**Functionality:** **The function accepts a parameter `number`, returns `True` if `number` is even. If `number` is not even, it returns `False`.**

### Example 2  
**Parameter constraints and relationships:** `age` is an integer.  
```python
def func(age):
```
**Output:**  
- Case 1: If `age` is less than 0, the function returns an error message stating that ages can't be negative.
- Case 2: If `age` is between 0 and 18, the function returns "minor".
- Case 3: Otherwise, the function returns "adult".

**Final Answer:**  
The function `func` accepts a parameter `age`. If `age` is below 0, the function returns an error message. If `age` is between 0 and 18, it returns "minor". Otherwise, it returns "adult".  
**Functionality:** **The function accepts an integer `age`, returns an error if `age` is negative, "minor" if `age` is between 0 and 18, and "adult" otherwise.**

### **Your Task**

**Parameter constraints:**  
{pre}

```python
{head}
```

**Output:**  
{body_post}

Now, please think step by step: What are the parameters the function accepts, and what does it return?

Strictly adhere to the format: **Functionality: `the functionality you calculate`**, and describe this functionality in **natural language easily understandable by humans**.

\end{lstlisting}
\end{tcolorbox}


\subsection{If Condition Precondition Extraction Prompt}
\label{if_pre_prompt}
\begin{tcolorbox}[ colback=gray!10, colframe=black, arc=5pt]
\begin{lstlisting}[language=, basicstyle=\ttfamily\footnotesize]
You are assigned the role of a program verifier, responsible for finding the **postcondition** of an `if` statement based on its condition. You will be given:
- A **precondition**, describing the initial state of the program variables before entering the `if` statement.
- An **if condition**, which determines whether the program enters the `if` block.

Your task is to determine the **postcondition**, which describes the state of the program variables after entering the `if` block. The postcondition must extend the precondition by ensuring that the `if` condition is true. This description should include both the values of the variables and the relationships between them.

- **Do not explain how the program runs**; only focus on variable values and relationships.
- Ensure that the postcondition retains all valid information from the precondition while incorporating the truth of the `if` condition.
- Follow the format: **Postcondition: `calculated postcondition`**.

I am giving you some examples to understand the task better. Then I am giving you your task.

---

### Example 1  
**Precondition:** `str` is a string  
**If condition:**  
```python
if len(str) < 3:
```
**Postcondition:** **`str` is a string, and the length of `str` is less than 3**

---

### Example 2  
**Precondition:** `n` can have any value  
**If condition:**  
```python
if isinstance(n, int):
```
**Postcondition:** **`n` is an integer of any value**

---

### Example 3  
**Precondition:** `x` is a positive integer  
**If condition:**  
```python
if x < 2:
```
**Postcondition:** **`x` is a positive integer less than 2**

---

### Example 4  
**Precondition:** `m` is an integer. If `m` is higher than 10, then `n` equals `-m`. `a` is a list of integers.  
**If condition:**  
```python
if n < 0:
```
**Postcondition:** **`m` is an integer, if `m` is higher than 10, `n` equals `-m`. `a` is a list of integers. The current value of `n` is less than 0**

\end{lstlisting}
\end{tcolorbox}


\begin{tcolorbox}[ colback=gray!10, colframe=black, arc=5pt]
\begin{lstlisting}[language=, basicstyle=\ttfamily\footnotesize]
### Example 5  
**Precondition:** `x` is an integer, `a` is a list of integers.  
**If condition:**  
```python
if a[0] != 0:
```
**Postcondition:** **`x` is an integer, `a` is a list of integers. The first element of `a` is not 0**

### **Your Task**

**Precondition:**  
{precondition}

**If condition:**  
```python
{program_fragment}
```

Your task is to complete the **postcondition**. Ensure that:
- All valid information from the **precondition** is retained.
- The **if condition** is now assumed to be **true**.
- If variables have values related to previous conditions, put that early in the postcondition and state the **current** value of the variable clearly.
- The final state of the program **after entering the if block** is fully described.

Strictly adhere to the format: **Postcondition: `the postcondition you calculate`**, and describe this postcondition in **natural language easily understandable by humans**.

\end{lstlisting}
\end{tcolorbox}

\subsection{Else Condition Precondition Extraction Prompt}
\label{else_pre_prompt}
\begin{tcolorbox}[ colback=gray!10, colframe=black, arc=5pt]
\begin{lstlisting}[language=, basicstyle=\ttfamily\footnotesize]
You are assigned the role of a program verifier, responsible for finding the **postcondition** of an `else` statement based on the condition of the corresponding `if` statement. You will be given:
- A **precondition**, describing the initial state of the program variables before evaluating the `if` condition.
- An **if condition**, which determines whether the program enters the `if` block. **In this case, we do not enter the `if` block but instead follow the `else` block**.

Your task is to determine the **postcondition**, which describes the state of the program variables after entering the `else` block. The postcondition must extend the precondition by ensuring that the `if` condition is **false**. This description should include both the values of the variables and their relationships.

- **Do not explain how the program runs**; only focus on variable values and relationships.
- Ensure that the postcondition retains all valid information from the precondition while incorporating the negation of the `if` condition.
- Follow the format: **Postcondition: `calculated postcondition`**.

I am giving you some examples to understand the task better. Then I am giving you your task.

---

### Example 1  
**Precondition:** `str` is a string  
**If condition:**  
```python
if len(str) < 3:
```
**Postcondition:** **`str` is a string, and the length of `str` is greater than or equal to 3**

---

### Example 2  
**Precondition:** `n` can have any value  
**If condition:**  
```python
if isinstance(n, int):
```
**Postcondition:** **`n` can have any value except an integer**

---

### Example 3  
**Precondition:** `x` is a positive integer  
**If condition:**  
```python
if x < 2:
```
**Postcondition:** **`x` is a positive integer greater than or equal to 2**

---

### Example 4  
**Precondition:** `m` is an integer, `n` is an integer, `a` is a list of integers.  
**If condition:**  
```python
if n <= 0:
```
**Postcondition:** **`m` and `n` are integers, `n` is greater than 0, `a` is a list of integers**

\end{lstlisting}
\end{tcolorbox}
\begin{tcolorbox}[ colback=gray!10, colframe=black, arc=5pt]
\begin{lstlisting}[language=, basicstyle=\ttfamily\footnotesize]
### Example 5  
**Precondition:** `x` is an integer, `a` is a list of integers.  
**If condition:**  
```python
if a[0] != 0:
```
**Postcondition:** **`x` is an integer, `a` is a list of integers. The first element of `a` is 0**

### **Your Task**

**Precondition:**  
{precondition}

**If condition:**  
```python
{program_fragment}
```

Your task is to complete the **postcondition**. Ensure that:
- All valid information from the **precondition** is retained.
- The **if condition** is now assumed to be **false**.
- If variables have values related to previous conditions, put that early in the postcondition and state the **current** value of the variable clearly.
- The final state of the program **after entering the else block** is fully described.

Strictly adhere to the format: **Postcondition: `the postcondition you calculate`**, and describe this postcondition in **natural language easily understandable by humans**.

\end{lstlisting}
\end{tcolorbox}

\subsection{If-Else Statement Postcondition Extraction Prompt}
\label{if_else_prompt}
\begin{tcolorbox}[ colback=gray!10, colframe=black, arc=5pt]
\begin{lstlisting}[language=, basicstyle=\ttfamily\footnotesize]
You are assigned the role of a program verifier, responsible for completing the **overall postcondition** of Hoare triples for `if` statements based on the postconditions of both the `if` and `else` parts. You will be given:
- A **precondition**, describing the initial state of the program variables before the execution of the program fragment.
- A **program fragment**, which contains an `if` condition and possibly an `else` block.
- The **postcondition of the `if` part**, describing the state after entering the `if` block.
- The **postcondition of the `else` part**, describing the state after entering the `else` block (if an `else` exists).

Your task is to **combine the postconditions of both the `if` and `else` parts** (if an `else` exists), taking into account the `if` condition, to derive the **overall postcondition of the entire `if-else` block**.

- **Do not explain how the program runs**; only focus on variable values and relationships.
- Ensure that the postcondition retains all valid information from the precondition while incorporating both branches of the `if-else` logic.
- Summarize the `if` statement in a **coherent paragraph**, rather than discussing it in separate sections.
- Follow the format: **Postcondition: `calculated postcondition`**.

I am giving you some examples to understand the task better. Then I am giving you your task.

---

### Example 1  
**Precondition:** `str` is a string  
**Program Fragment:**  
```python
if len(str) < 3:
```
**If part:** `str` is a string with length less than 3, the function returns `None`  
**Else part:** There is no `else` part in the code  

**Postcondition:** **`str` is a string. If the length of `str` is less than 3, the function returns `None`.**

---

### Example 2  
**Precondition:** `n` can have any value  
**Program Fragment:**  
```python
if isinstance(n, int):
    
else:
```
**If part:** `n` is an integer of any value, and the function returns `n`  
**Else part:** `n` can have any value except an integer, and the function returns `int(n)`  

**Postcondition:** **If `n` is an integer, the function returns `n` itself. Otherwise, the function returns `int(n)`.**

---

### Example 3  
**Precondition:** `x` is a positive integer  
**Program Fragment:**  
```python
if x < 2:
    
else:
```
**If part:** `x` is a positive integer less than 2, and the function returns `False`  
**Else part:** `x` is a positive integer greater than or equal to 2, and the function returns `True`  

**Postcondition:** **`x` is a positive integer. If `x` is less than 2, the function returns `False`. Otherwise, the function returns `True`.**


\end{lstlisting}
\end{tcolorbox}

\begin{tcolorbox}[ colback=gray!10, colframe=black, arc=5pt]
\begin{lstlisting}[language=, basicstyle=\ttfamily\footnotesize]
### Example 4  
**Precondition:** `m` is an integer, `n` is an integer  
**Program Fragment:**  
```python
if n < 0:
    
else:
```
**If part:** The integer `n` was originally negative, `n` is updated to its negation, and `m` is increased by 1  
**Else part:** If `n` is 0, the function returns `m`. Otherwise, `n` is decreased by 13, and `m` is increased by 1  

**Postcondition:** **`m` and `n` are integers. If `n < 0`, `m` is increased by 1 and `n` is negated. If `n == 0`, the function returns `m`. Otherwise, `n` is decreased by 13 and `m` is increased by 1.**

---

### Example 5  
**Precondition:** `x` is an integer, `y` is zero  
**Program Fragment:**  
```python
if x > 0:
```
**If part:** `x` is a positive integer. If `x > 10`, `y` is set to twice the value of `x`. Otherwise, `y` is set to `x + 5`.  
**Else part:** There is no `else` part in the code  

**Postcondition:** **`x` is an integer. If `x > 10`, `y` is set to twice the value of `x`. If `x` is greater than 0 but less than 10, `y` is set to `x + 5`.**

---

### **Your Task**

**Precondition:**  
{precondition}

**Program Fragment:**  
```python
{program_fragment}
```

**If part:**  
{postconditions_if}

**Else part:**  
{postconditions_else}

Your task is to **complete the overall postcondition of the entire if-else block**. Summarize all cases and describe the final state of the program **after executing the if-else statement**.

Strictly adhere to the format: **Postcondition: `the postcondition you calculate`**, and describe this postcondition in **natural language easily understandable by humans**.

\end{lstlisting}
\end{tcolorbox}

\subsection{While Loop First Execution Necessary Condition Prompt}
\label{while_first_prompt}
\begin{tcolorbox}[ colback=gray!10, colframe=black, arc=5pt]
\begin{lstlisting}[language=, basicstyle=\ttfamily\footnotesize]
You have been assigned the task of a program verifier, responsible for modifying the **program state** so that the first iteration of the `while` loop can proceed. You will be provided with the **program state right before the loop**, which you need to modify, and the `while` loop statement.

- If the loop is a `while True` loop or if the loop **can certainly execute at least once**, simply repeat the program state right before the loop.
- **Do not make any assumptions** beyond the provided information.
- Only modify the states of **objects in the loop condition** that affect whether the loop executes.
- **Follow the text format:** **State: `state`**.

I am giving you some examples to understand the task better. Then I am giving you your task.

### Example 1  
**State right before the while loop:** `total` is 10, `i` is 0, `n` is an integer.  
**While loop:**  
```python
while i < n:
    # the loop body is omitted
```
Now, please think step by step: Which states need to be adjusted for the loop to execute the first time? **Only the states of objects in the loop head can be adjusted.**

**Example Answer:**  
The variables in the loop condition are `i` and `n`, so we can only adjust them. The loop executes while `i < n`. Right before the loop, `i` is 0 and `n` is an integer. Since `n` being an integer does not ensure the loop executes, it needs to be modified to `n` is greater than 0.

**State:** **`total` is 10, `i` is 0, `n` must be greater than 0**

### Example 2  
**State right before the while loop:** `total` is 0, `students` is 2 less than its initial value.  
**While loop:**  
```python
while students >= 1:
    # the loop body is omitted
```
Now, please think step by step: Which states need to be adjusted for the loop to execute the first time? **Only the states of objects in the loop head can be adjusted.**

**Example Answer:**  
The only variable in the loop condition is `students`. The loop executes while `students >= 1`. Right before the loop, `students` is **2 less than its initial value**, so for the loop to execute at least once, the initial value of `students` must have been at least 3, and its current value must be greater than or equal to 1.

**State:** **`total` is 0, `students` is 2 less than its initial value, and `students` must currently be greater than or equal to 1**

### **Your Task**

**State right before the while loop:**  
{post}

```python
{loop_head}
    # the loop body is omitted
```

Now, please think step by step: Which states need to be adjusted for the loop to execute the first time? **Only modify the states of objects in the loop head.*
Strictly adhere to the format: **State: `the state you calculate`**, and describe this state in **natural language easily understandable by humans**.

\end{lstlisting}
\end{tcolorbox}



\subsection{While Loop Prerequisite for Next Execution Prompt}
\label{while_next_prompt}
\begin{tcolorbox}[ colback=gray!10, colframe=black, arc=5pt]
\begin{lstlisting}[language=, basicstyle=\ttfamily\footnotesize]
You have been assigned the task of a program verifier, responsible for modifying the **program state** so that the next iteration of the `while` loop can proceed. You will be provided with the **program state after the previous iteration**, which you need to modify, and the `while` loop statement.

- If the loop is a `while True` loop or if the loop **can certainly execute one more time**, simply repeat the program state at the end of the previous iteration.
- **Do not make any assumptions** beyond the provided information.
- Only modify the states of **objects in the loop condition** that affect whether the loop executes.
- **Follow the text format:** **State: `state`**.

I am giving you some examples to understand the task better. Then I am giving you your task.

### Example 1  
**State at the end of the previous iteration:** `total` is 10, `i` is 4, `n` is greater than 3.  
**While loop:**  
```python
while i < n:
    # the loop body is omitted
```
Now, please think step by step: Which states need to be adjusted for the loop to execute one more time? **Only the states of objects in the loop head can be adjusted.**

**Example Answer:**  
The variables in the loop condition are `i` and `n`, so we can only adjust them. The loop executes while `i < n`. At the end of the last iteration, `i` is 4, `n` is greater than 3. `n` being greater than 3 does not ensure another execution, so it needs to be modified to `n` is greater than 4.

**State:** **`total` is 10, `i` is 4, `n` must be greater than 4**


### Example 2  
**State at the end of the previous iteration:** `total` is 0, `students` is 3 less than its initial value.  
**While loop:**  
```python
while students >= 1:
    # the loop body is omitted
```
Now, please think step by step: Which states need to be adjusted for the loop to execute one more time? **Only the states of objects in the loop head can be adjusted.**

**Example Answer:**  
The only variable in the loop condition is `students`. The loop executes while `students >= 1`. At the end of the last iteration, `students` is **3 less than its initial value**, so for the loop to execute again, the initial value of `students` must have been at least 4, and its current value must be greater than or equal to 1.

**State:** **`total` is 0, `students` is 3 less than its initial value, and `students` must currently be greater than or equal to 1**

### **Your Task**

**State at the end of the previous iteration:**  
{post}

```python
{loop_head}
    # the loop body is omitted
```

Now, please think step by step: Which states need to be adjusted for the loop to execute one more time? **Only modify the states of objects in the loop head.**
Strictly adhere to the format: **State: `the state you calculate`**, and describe this state in **natural language easily understandable by humans**.

\end{lstlisting}
\end{tcolorbox}

\subsection{For Loop First Execution Necessary Condition Prompt}
\label{for_first_prompt}

\begin{tcolorbox}[ colback=gray!10, colframe=black, arc=5pt]
\begin{lstlisting}[language=, basicstyle=\ttfamily\footnotesize]
You have been assigned the task of a program verifier, responsible for understanding the **program state at the start of the `for` loop**. You will be provided with the **program state before the first execution of the `for` loop**, which you need to modify, and the `for` loop statement.

- **Do not make any assumptions** beyond the provided information.
- Only modify the states of **objects in the loop condition** that affect whether the loop executes.
- **Follow the text format:** **State: `state`**.

I am giving you some examples to understand the task better. Then I am giving you your task.

### Example 1  
**State before the `for` loop:** `total` is 10.  
**For loop:**  
```python
for i in range(n):
    # the loop body is omitted
```
Now, please think step by step: Which states need to be adjusted for the loop to execute? **Only the states of objects in the loop head can be adjusted.**

**Example Answer:**  
The only variables in the loop head are `i` and `n`, so we can only adjust those. The loop executes while `n` is at least 1. Since `total` being 10 does not ensure that the loop will execute, `n` needs to be adjusted to be greater than 0, and `i` is initialized to 0.

**State:** **`total` is 10, `i` is 0, `n` must be greater than 0**


### Example 2  
**State before the loop starts:** `total` is 0, `students_list` is a list of students.  
**For loop:**  
```python
for index, student in enumerate(students_list):
    # the loop body is omitted
```
Now, please think step by step: Which states need to be adjusted for the loop to execute? **Only the states of objects in the loop head can be adjusted.**

**Example Answer:**  
The only objects in the loop head are `index`, `student`, and `students_list`, so we can only adjust those. The loop executes if `students_list` has at least 1 student. Since the `total` value does not impact execution, we focus on ensuring that `students_list` contains at least 1 student. The `index` starts at 0, and `student` is set to the first student in the list.

**State:** **`total` is 0, `students_list` is a list of students that must have at least 1 student, `student` is the first student in the list, `index` is 0**

### **Your Task**

**State before the loop starts:**  
{post}

```python
{loop_head}
    # the loop body is omitted
```

Now, please think step by step: Which states need to be adjusted for the loop to execute? **Only modify the states of objects in the loop head.**

Strictly adhere to the format: **State: `the state you calculate`**, and describe this state in **natural language easily understandable by humans**.

\end{lstlisting}
\end{tcolorbox}

\subsection{For Loop Prerequisite for Next Execution Prompt}
\label{for_next_prompt}
\begin{tcolorbox}[ colback=gray!10, colframe=black, arc=5pt]
\begin{lstlisting}[language=, basicstyle=\ttfamily\footnotesize]
You have been assigned the task of a program verifier, responsible for understanding the **program state at the start of the next iteration of a `for` loop**. You will be provided with the **program state after the previous iteration**, which you need to modify, and the `for` loop statement.
- **Do not make any assumptions** beyond the provided information.
- Only modify the states of **objects in the loop condition** that affect whether the loop executes again.
- **Follow the text format:** **State: `state`**.

I am giving you some examples to understand the task better. Then I am giving you your task.

### Example 1  
**State at the end of the previous iteration:** `total` is 10, `i` is 3, `n` must be greater than 3.  
**For loop:**  
```python
for i in range(n):
    # the loop body is omitted
```
Now, please think step by step: Which states need to be adjusted at the start of the next iteration of the loop? **Only the states of objects in the loop head can be adjusted.**

**Example Answer:**  
The only variables in the loop head are `i` and `n`, so we can only adjust those. The loop executes while `i` is less than `n`. At the end of the last iteration, `i` is 3, `n` is greater than 3. Since `i` is incremented by 1 in each iteration, for the loop to execute again, `i` must be 4 and `n` must be greater than 4.

**State:** **`total` is 10, `i` is 4, `n` must be greater than 4**

### Example 2  
**State at the end of the previous iteration:** `total` is 0, `students_list` is a list of students that must have at least 2 students, `student` is the second student in the list, `index` is 1.  
**For loop:**  
```python
for index, student in enumerate(students_list):
    # the loop body is omitted
```
Now, please think step by step: Which states need to be adjusted for the loop to execute one more time? **Only the states of objects in the loop head can be adjusted.**

**Example Answer:**  
The only objects in the loop head are `index`, `student`, and `students_list`, so we can only adjust those. The loop executes if `students_list` has at least 3 students. At the end of the last iteration, `students_list` has at least 2 students, `student` is the second student in the list, and `index` is 1. So, for the loop to execute one more time, `students_list` must have at least 3 students, `index` must be 2, and `student` must be the third student in the list.

**State:** **`total` is 0, `students_list` is a list of students that must have at least 3 students, `student` is the third student in the list, `index` is 2**

### **Your Task**

**State at the end of the previous iteration:**  
{post}

```python
{loop_head}
    # the loop body is omitted
```

Now, please think step by step: Which states need to be adjusted for the loop to execute one more time? **Only modify the states of objects in the loop head.**
Strictly adhere to the format: **State: `the state you calculate`**, and describe this state in **natural language easily understandable by humans**.
\end{lstlisting}
\end{tcolorbox}


\subsection{Total While Loop Execution with Unrolling Prompt}
\label{total_while_prompt}
\begin{tcolorbox}[ colback=gray!10, colframe=black, arc=5pt]
\begin{lstlisting}[language=, basicstyle=\ttfamily\footnotesize]
Given a Python loop, an initial execution state, and the output states after the first {times} iterations of the loop, determine the output state after all the executions of the loop have finished.

You must adhere to the text format: Output State: **output state.**

Initial State: {pre}
Code of the loop:
{loop_code}

The output state after the loop executes the first {times} times includes what needed to be true for the loop to execute at least that number of times:
{loop_unrolled}

What is the ouput state after the loop executes all the iterations? Change the values of only the variables in the loop head and body. The state of the other variables in the precondition that are not affected by the loop head and body must remain unchanged.
In your response strictly use the format: Output State: **the output state you calculate.**, and describe this output state in Natural language easily understandable by humans.

\end{lstlisting}
\end{tcolorbox}


\subsection{Total While Loop Execution Without Unrolling Prompt}
\label{while_NO_UNROLL_PROMPT}
\begin{tcolorbox}[ colback=gray!10, colframe=black, arc=5pt]
\begin{lstlisting}[language=, basicstyle=\ttfamily\footnotesize]
Given a Python loop and an initial execution state, determine the output state after all the executions of the loop have finished.

You must adhere to the text format: Output State: **output state.**

Initial State: {pre}
Code of the loop:
{loop_code}

What is the output state after the loop executes all the iterations? Change the values of only the variables in the loop head and body. The state of the other variables in the precondition that are not affected by the loop head and body must remain unchanged.
In your response strictly use the format: Output State: **the output state you calculate.**, and describe this output state in Natural language easily understandable by humans.

\end{lstlisting}
\end{tcolorbox}


\subsection{Total For Loop Execution with Unrolling Prompt}
\label{total_for_prompt}
\begin{tcolorbox}[ colback=gray!10, colframe=black, arc=5pt]
\begin{lstlisting}[language=, basicstyle=\ttfamily\footnotesize]
Given a Python loop, an initial execution state, and the output states after the first {times} iterations of the loop, determine the output state after all the executions of the loop have finished.

You must adhere to the text format: Output State: **output state.**

Initial State: {pre}
Code of the loop:
{loop_code}

The output state after the loop executes the first {times} of times includes what needed to be true for the loop to execute at least that number of times:
{loop_unrolled}

What is the ouput state after the loop executes all the iterations? Change the values of only the variables in the loop head and body. The state of the other variables in the precondition that are not affected by the loop head and body must remain unchanged.
In your response strictly use the format: Output State: **the output state you calculate.**, and describe this output state in Natural language easily understandable by humans.

\end{lstlisting}
\end{tcolorbox}

\subsection{Total For Loop Execution Without Unrolling Prompt}
\label{for_NO_UNROLL_PROMPT}
\begin{tcolorbox}[ colback=gray!10, colframe=black, arc=5pt]
\begin{lstlisting}[language=, basicstyle=\ttfamily\footnotesize]
Given a Python loop, and an initial execution state, determine the output state after all the executions of the loop have finished. 

You must adhere to the text format: Output State: **output state.**

Initial State: {pre}
Code of the loop:
{loop_code}

What is the ouput state after the loop executes all the iterations? Change the values of only the variables in the loop head and body. The state of the other variables in the precondition that are not affected by the loop head and body must remain unchanged.
The output state must be in a similar format as the initial execution state. Describe this output state in Natural language easily understandable by humans. In your response strictly use the format: Output State: **the output state you calculate.**.

\end{lstlisting}
\end{tcolorbox}

\end{document}


